\section{Boundary open--closed type theorem}\label{sec:typed-master}

The formal-Darboux stalk theorem of \S\ref{ssec:local-theorem} gives
one fibre of the mixed holomorphic--topological open--closed
correspondence.  The Hamiltonian BF bulk, evaluated at a Dirac brane
vacuum, acts on the boundary Koszul algebra through the finite-window
trace coordinates \(J_{\leq W}\) and their stable scalar-reduced inverse
limit \(J\); the resulting
Chevalley--Eilenberg/polyvector chart is the trace-sector coordinate of
the local derived-centre comparison.  The global comparison of
Theorem~\ref{thm:main-global-deligne-criterion} is the descent problem for
these fibres over the relative Costello--Gwilliam brane-link boundary
\(\partial X:=\partial_{\mathrm{link}}X\).

The descent problem separates six typed objects.  The \(E_1\)-bar coalgebra of
the brane algebra, its Koszul dual algebra, the derived mixed-HT
centre, the physical bulk factorization algebra, the protected scalar
trace with raw Rees/Weyl Capelli anomaly, and the filtered boundary operator algebra
live in different categories and are related by definite morphisms.
The functor \(\Phi^{\mathrm{HT}}\), the supplied open--closed comparison
morphism \(\mathrm{OCA}\), the mixed-HT Deligne--Swiss-cheese action
datum, and scalar non-faithfulness are the structure maps between these
objects.  A topologisation datum, when supplied, adds operators
\(G_i\) trivialising holomorphic translations by
\(T_i=[Q,G_i]\).  The type table records the category in which each
object and map is defined.

\subsection{The six objects and their ambient categories}\label{ssec:six-typed-objects}

\begin{defn}[Six boundary open--closed objects]\label{defn:six-boundary-open-closed-objects}
Fix the formal-Darboux local model
\(X=\R^2_{\mathrm{top}}\times\widehat{\C^2}_{\mathrm{hol}}\),
brane line \(L=\R_t\times\{0_s\}\times\{b\}\), reduced Hamiltonian Lie algebra
\(\mathfrak h=\C[[z_1,z_2]]/\C\cdot 1\), residue dual
\[
  P=\mathfrak h^\vee_{\mathrm{cont}}
  =\operatorname{Ann}_{\operatorname{Res}}(\C\cdot1)
  \subset H^2_{(z_1,z_2)}(\C[[z_1,z_2]])\,dz_1dz_2
  =\bigoplus_{a+b>0}\C\,\rho_{a,b},
\]
shifted-cotangent BF Lie
algebra \(\mathfrak g=\mathfrak h\ltimes P[1]\), canonical
constant-Hamiltonian central extension
\(\widehat{\mathfrak g}_{\mathrm{cent}}\), non-effective boundary
representative
\[
  \mathcal R_N^{\mathrm{neff}}
  =
  \mathrm{Sym}(\mathfrak{gl}_N^{\oplus2}\oplus\mathfrak{gl}_N[1]),
  \qquad
  A^{\mathrm{cl}}_{\partial,N}=
  C^\bullet_{\mathrm{CE}}(\mathfrak{gl}_N,\mathcal R_N^{\mathrm{neff}})
\]
with \(Q\psi=[\phi_1,\phi_2]\), and admissible category
\(\mathcal B^{\mathrm{adm}}\) of
Definition~\ref{defn:admissible-host-category}.  The effective Hamiltonian
version replaces \(\mathfrak{gl}_N[1]\) by \(\mathfrak{sl}_N[1]\) and
\(\mathfrak{gl}_N\) by \(\mathfrak{pgl}_N\); the scalar-stable trace theorem
removes the central \(\operatorname{Tr}\psi\) excess.  The formal completion
theorem is Appendix~\ref{app:completion-formal-local-theorem}.  Put
\(\bar A=\C[z_1,z_2]/\C\cdot1\).
The normalization of the central line is
\[
  0\longrightarrow\C K\longrightarrow
  \widehat{\mathfrak g}_{\mathrm{cent}}
  \longrightarrow\mathfrak g\longrightarrow0,
  \qquad
  [(f,\eta,aK),(g,\xi,bK)]_{\mathrm{cent}}
  =
  \bigl(\overline{\{f,g\}},\operatorname{ad}^{\vee}_f\xi
  -\operatorname{ad}^{\vee}_g\eta,\omega(f,g)K\bigr),
\]
where
\(\omega(f,g)=[\{f,g\}]_0\), computed on the zero-constant
representatives, and the central component vanishes if either argument
lies in \(P[1]\).  Thus the source extension is unscaled and has class
\([\bar c]=[\omega|_{\bar A}]\).  The rank enters only through the trace
character
\[
  \chi_N(K)=\operatorname{Tr}I_N=N.
\]
In the Rees/Weyl algebra the raw supercommutator satisfies
\[
  xy-(-1)^{|x||y|}yx
  =\hbar_{\mathrm W}[x,y]_{\mathrm{cent}}.
\]
Consequently its central component after the pushout by \(\chi_N\) is
\(\hbar_{\mathrm W}N\omega(f,g)\mathbf1\), with class
\(\hbar_{\mathrm W}N[\bar c]\).  The normalized Moyal bracket obtained
by dividing the raw commutator by \(\hbar_{\mathrm W}\) has central
class \(N[\bar c]\).  Neither \(N\) nor \(\hbar_{\mathrm W}\) occurs in
the canonical source extension.

The six boundary open--closed objects are:
\begin{enumerate}[label=\textup{(C\arabic*)}]
  \item \(\bar B^{E_1}(A^{\mathrm{cl}}_{\partial,N})
        =T^{\mathrm{c}}(s^{-1}\bar A^{\mathrm{cl}}_{\partial,N})\),
        the conilpotent \(E_1^{\mathrm{top}}\)-bar coalgebra of the
        augmented brane algebra
        \(\bar A^{\mathrm{cl}}_{\partial,N}\):
        a coassociative conilpotent dg coalgebra in
        \(\mathcal B^{\mathrm{adm}}\) with deconcatenation coproduct and
        differential \(d_{\bar B}=d_{\mathrm{int}}+d_m\)
        of \eqref{eq:bar-differential-explicit};
  \item \((A^{\mathrm{cl}}_{\partial,N})^!\), the Verdier/Koszul
        dual algebra of \(A^{\mathrm{cl}}_{\partial,N}\) on the Koszul
        locus, in the sense of
        Theorem~\ref{thm:universal-ce-pv-koszul-criterion};
  \item \(Z^{\mathrm{der},\mathrm{Mor}}_{E_2^{\mathrm{HT}}}\!\bigl(
        \mathcal C^{\mathrm{op}}_\partial\bigr)\), the derived mixed-HT
        \(E_2\)-centre of the primitive open boundary factorization
        category \(\mathcal C^{\mathrm{op}}_\partial\) on the
        Costello--Gwilliam brane-link boundary; the bulk action or OCA
        map is extra structure;
  \item \(\mathcal A^{\mathrm{cl}}_{\mathrm{bulk}}\), the physical
        bulk mixed-HT Hamiltonian BF factorisation algebra of
        Definition~\ref{def:closed-mixed-HT-factorization};
  \item the protected scalar trace together with the projective anomaly
        class
        \[
          \bigl(\chi,\kappa_{\mathrm{Cap}}\bigr),\qquad
          \chi=\operatorname{sTr}\,(-)\colon
          \mathcal A^{\mathrm{cl}}_{\mathrm{bulk}}\to\Z((q,\hbar_{\mathrm{W}})),
          \quad
          \kappa_{\mathrm{Cap}}=\hbar_{\mathrm{W}}N[\bar c]\in
          H^2_{\mathrm{Lie}}(\bar A,\C)[[\hbar_{\mathrm{W}}]] .
        \]
        Here \(\kappa_{\mathrm{Cap}}\) is the raw Rees/Weyl
        trace-commutator class obtained from the unscaled class
        \([\bar c]\) by \(\chi_N(K)=N\); the scalar trace is its protected
        character.  A determinant-line or modular-line curvature
        interpretation requires a supplied line, connection, and
        Atiyah-class pullback identifying its curvature with this Lie
        class.
  \item the filtered boundary operator algebra
        \(\mathrm{Op}^{\mathrm{HT}}_{\partial X}\), a filtered
        \(\mathrm{SC}^{\mathrm{HT}}\)-object on \(\partial X\)
        (\S\ref{ssec:open-modular-lift}).
\end{enumerate}
\end{defn}

\begin{defn}[Typed Morita centre and Hochschild chart]
\label{defn:typed-morita-centre-hochschild-chart}
Let \(\widehat{\mathcal C}^{\mathrm{op}}_{\partial,\mathrm{adm}}\)
be the admissible completion of the primitive open boundary
factorization category.  Here
\(\operatorname{Ran}(\partial X)\) denotes the analytic brane-link Ran
site of finite configurations of opens in the Costello--Gwilliam
brane-link boundary, with Weiss covers.  A Kato--Nakayama boundary model
enters only after a chosen log compactification and a comparison functor
from the brane-link model.  The analytic brane-link Ran site is not the
Beilinson--Drinfeld algebraic \(\mathcal D\)-module Ran space except
after a specified comparison datum.  Let
\[
  \mathrm{Bimod}^{\mathrm{HT}}_{\operatorname{Ran}(\partial X)}
  \bigl(\mathcal C^{\mathrm{op}}_\partial,
        \mathcal C^{\mathrm{op}}_\partial\bigr)
\]
be the stable \(\infty\)-category of complete mixed-HT factorization
bimodules on the brane-link Ran space, with continuous left and right
\(\widehat{\mathcal C}^{\mathrm{op}}_{\partial,\mathrm{adm}}\)-actions
and Weiss descent.  Let \(\Delta_\partial\) be the diagonal
\((\mathcal C^{\mathrm{op}}_\partial,
\mathcal C^{\mathrm{op}}_\partial)\)-bimodule.  The centre object in
\((C3)\) is
\[
  Z^{\mathrm{der},\mathrm{Mor}}_{E_2^{\mathrm{HT}}}
  \bigl(\mathcal C^{\mathrm{op}}_\partial\bigr)
  :=
  \operatorname{RHom}_{
  \mathrm{Bimod}^{\mathrm{HT}}_{\operatorname{Ran}(\partial X)}
  (\mathcal C^{\mathrm{op}}_\partial,
   \mathcal C^{\mathrm{op}}_\partial)}
  (\Delta_\partial,\Delta_\partial).
\]
This is the category-level centre before a brane vacuum is chosen.
For a dualisable object \(b\in\mathcal C^{\mathrm{op}}_\partial\), put
\(A_b=\operatorname{End}_{\mathcal C^{\mathrm{op}}_\partial}(b)\).  The
full subcategory generated by \(b\) gives a restriction functor on
bimodules and hence an evaluation morphism
\[
  \operatorname{ev}_b\colon
  Z^{\mathrm{der},\mathrm{Mor}}_{E_2^{\mathrm{HT}}}
  \bigl(\mathcal C^{\mathrm{op}}_\partial\bigr)
  \longrightarrow
  C^\bullet_{\mathrm{ch},\mathrm{HT},\mathrm{adm}}(A_b,A_b).
\]
At the \(N\)-Dirac-brane formal-Darboux vacuum \(b_p\), the algebra
\(A_{b_p}\) is represented by the quotient-stack Koszul cdga
\(A^{\mathrm{cl}}_{\partial,N}\).  The local theorem computes the
scalar-reduced stable trace-sector image of
\(\operatorname{ev}_{b_p}\); it does not identify the full Morita
centre with this Hochschild chart.
\end{defn}

\begin{prop}[Hochschild chart criterion for the Morita centre]
\label{prop:typed-morita-centre-hochschild-chart}
The evaluation morphism
\(\operatorname{ev}_b\) of
Definition~\ref{defn:typed-morita-centre-hochschild-chart} is an
equivalence only under the Morita-density condition that \(b\)
generates
\(\widehat{\mathcal C}^{\mathrm{op}}_{\partial,\mathrm{adm}}\) as an
idempotent-complete, colimit-closed mixed-HT factorization module
category.  Without that condition,
\[
  C^\bullet_{\mathrm{ch},\mathrm{HT},\mathrm{adm}}(A_b,A_b)
\]
is a Hochschild coordinate chart on the evaluation of the centre at
\(b\), not the definition of
\(Z^{\mathrm{der},\mathrm{Mor}}_{E_2^{\mathrm{HT}}}
(\mathcal C^{\mathrm{op}}_\partial)\).
\end{prop}

\begin{proof}
Restriction along the full subcategory generated by \(b\) sends the
diagonal bimodule \(\Delta_\partial\) to the diagonal
\((A_b,A_b)\)-bimodule.  Applying derived endomorphisms in the
corresponding admissible bimodule categories gives
\(\operatorname{ev}_b\).  If \(b\) Morita-generates, restriction of
complete factorization bimodules is fully faithful and essentially
surjective after idempotent completion and admissible colimits; derived
endomorphisms of the diagonal are therefore preserved.  If
Morita-density is not assumed, restriction forgets bimodule summands not
seen by \(b\).  The target remains the Hochschild cochain object of
\(A_b\), hence a chart on the centre at \(b\), not the category-level
centre itself.
\end{proof}

\begin{defn}[Morita-density quotient and chart cone]
\label{defn:morita-density-quotient-chart-cone}
Let
\[
  \langle b\rangle_{\mathrm{adm}}
  \subset
  \widehat{\mathcal C}^{\mathrm{op}}_{\partial,\mathrm{adm}}
\]
be the smallest full stable mixed-HT factorization module subcategory
which contains \(b\), is idempotent-complete, is closed under admissible
colimits, and is closed under the boundary factorization action.  The
\emph{Morita-density quotient} is the Verdier quotient in complete
stable \(\infty\)-categories
\[
  \mathcal Q_b^{\mathrm{Mor}}
  :=
  \widehat{\mathcal C}^{\mathrm{op}}_{\partial,\mathrm{adm}}
  /\langle b\rangle_{\mathrm{adm}} .
\]
The density obstruction is
\[
  \operatorname{Ob}^{\mathrm{dens}}_b
  =
  \mathcal Q_b^{\mathrm{Mor}};
\]
it vanishes exactly when \(b\) Morita-generates the admissible boundary
category.  The Hochschild-chart obstruction is the cone
\[
  \operatorname{Ob}^{\mathrm{chart}}_b
  =
  \operatorname{Cone}(\operatorname{ev}_b)
  \in\mathcal B^{\mathrm{adm}}.
\]
Thus \(\operatorname{ev}_b\) is an equivalence only after
\[
  \operatorname{Ob}^{\mathrm{dens}}_b=0,
  \qquad
  H^\bullet(\operatorname{Ob}^{\mathrm{chart}}_b)=0.
\]
At the Dirac brane \(b_p^{\oplus N}\), the local formal-Darboux theorem
computes the scalar-reduced stable trace-sector image in the target of
\(\operatorname{ev}_{b_p}\); it does not kill
\(\operatorname{Ob}^{\mathrm{dens}}_{b_p}\) or
\(\operatorname{Ob}^{\mathrm{chart}}_{b_p}\).
\end{defn}

\begin{defn}[Typed Hochschild-centre lift datum]
\label{defn:typed-hochschild-centre-lift-datum}
Let \(A=A^{\mathrm{cl}}_{\partial,N}\), and write
\[
  \overline C^{\mathrm{CE},\mathrm{cont}}_\bullet(\mathfrak g)
  =
  \ker\!\left(
  C^{\mathrm{CE},\mathrm{cont}}_\bullet(\mathfrak g)
  \xrightarrow{\varepsilon}\C\right),
  \qquad
  \mathcal H_N=
  C^\bullet_{\mathrm{Hoch},\mathrm{adm}}(A,A)[1].
\]
A typed Hochschild action is encoded covariantly by a continuous
twisting cochain
\[
  \tau_N\colon
  \overline C^{\mathrm{CE},\mathrm{cont}}_\bullet(\mathfrak g)
  \longrightarrow \mathcal H_N .
\]
Its restriction to the cogenerators of the Chevalley--Eilenberg chain
coalgebra, followed by desuspension, is the first Taylor coefficient
\begin{equation}\label{eq:typed-centre-first-taylor}
  \tau_{N,1}\colon
  \mathfrak g\longrightarrow
  C^\bullet_{\mathrm{Hoch},\mathrm{adm}}(A,A)[1].
\end{equation}
The formal-Darboux theorem prescribes the Hamiltonian trace-sector part
\(\tau_{N,1}^{\mathrm{tr}}\) of this coefficient.  Only after pairing it
with the continuous linear coordinates in
\(C^\bullet_{\mathrm{CE},\mathrm{cont}}(\mathfrak g)\) does one obtain
the coordinate-dual shorthand
\[
  c_f\longmapsto\theta_f,\qquad
  u_f\longmapsto M_{J(f)},\qquad
  J(f)=\varprojlim_WJ_{\leq W}(f\bmod\mathfrak m^{W+1}).
\]
We denote this coordinate expression by \(\Phi_N^{\mathrm{gen}}\); it
is the coordinate shadow of \(\tau_{N,1}^{\mathrm{tr}}\), not a map
whose source is the set of generators of a CE cochain algebra.  The
component of \(\tau_{N,1}\) on the cotangent-current summand \(P[1]\)
is present only after the datum of
Definition~\ref{defn:typed-cotangent-current-lift-datum} has supplied
the corresponding continuous current cochain.

Let \(\mathcal K^{\mathrm{cent}}_{\partial,N}\) be the completed
Hochschild-centrality convolution \(L_\infty\)-algebra whose target is
the shifted Hochschild complex
\(C^\bullet_{\mathrm{Hoch},\mathrm{adm}}(A,A)[1]\), with
Gerstenhaber bracket of degree \(0\), completed in the arity,
open-input, finite-window, and admissible-topology filtrations.  More
precisely, it is the closed \(L_\infty\)-subalgebra
\[
  \mathcal K^{\mathrm{cent}}_{\partial,N}
  \subset
  \underline{\mathrm{Hom}}_{\mathcal B^{\mathrm{adm}}}
  \left(
    \overline C^{\mathrm{CE},\mathrm{cont}}_\bullet
      (\mathfrak g_{\mathrm{src}}),\,
    C^\bullet_{\mathrm{Hoch},\mathrm{adm}}(A,A)[1]
  \right)
\]
of continuous filtered maps preserving the Hochschild operation list
\[
  \mathcal O^{\mathrm{cent}}_{\partial,N}
  =
  \{d_{\mathrm{Hoch}},\,\smile,\,[\, ,\,]_G,\,
    M_{(-)},\,\rho,\,\rho^{P_0},\,\pi_{M,M'},\,r_{\alpha\beta}\},
\]
where \(\pi_{M,M'}\) are finite-window projections and
\(r_{\alpha\beta}\) denotes arity/open-input restriction.  For the
ordinary Hochschild rows
\(\mathfrak g_{\mathrm{src}}=\mathfrak g\) and no kernel is declared.
For the restricted \(P_0\) row
\(\mathfrak g_{\mathrm{src}}=\mathfrak g_w\), the kernel list is
\[
  \mathcal K^{\mathrm{ker,cent}}_{\partial,N}
  =
  \{C_{\mathfrak h,w},\,K^{(2)}_\Delta\},
\]
with \(C_{\mathfrak h,w}\) the weighted diagonal Casimir and
\(K^{(2)}_\Delta=d\zeta_1d\zeta_2/(\zeta_1\zeta_2)\) the diagonal
local-cohomology kernel.  Thus a centrality primitive is a cochain in
this displayed Hom complex; it is not an abstract homotopy in an
unspecified completion.  A
typed Hochschild-centre lift is the continuous twisting cochain, written
as the Maurer--Cartan element
\[
  \tau_N=\mathcal C_{\Phi_N^{\mathrm{Hoch}}}
  =
  \tau_{N,1}
  +H^{\mathrm{Lie}}+H^{\mathrm{act}}
  +H^{\mathrm{prod}}+H^{P_0}+H^3+\sum_{m\ge4}H^m
  \in \mathcal K^{\mathrm{cent}}_{\partial,N}
\]
satisfying
\[
  d\mathcal C_{\Phi_N^{\mathrm{Hoch}}}
  +\frac12[\mathcal C_{\Phi_N^{\mathrm{Hoch}}},
           \mathcal C_{\Phi_N^{\mathrm{Hoch}}}]
  +\frac{1}{3!}\ell_3
    (\mathcal C_{\Phi_N^{\mathrm{Hoch}}}^{\otimes3})
  +\cdots=0.
\]
In unsuspended Hochschild degrees, the primary homotopies have degree
\(-1\), the next Stasheff primitive \(H^3\) has degree \(-2\), and in
general \(H^m\) has degree \(1-m\); after operadic suspension they are
the components of the degree-one Maurer--Cartan element above.

For homogeneous \(x,y\in\mathfrak g_{\mathrm{src}}\) and an open
observable \(a\),
put
\[
  \begin{aligned}
  \delta^{\mathrm{Lie}}_{x,y}
    &=[\tau_{N,1}(x),\tau_{N,1}(y)]_G
      -\tau_{N,1}([x,y]),\\
  \delta^{\mathrm{act}}_{x|y}
    &=[\tau_{N,1}(x),M_{O_y}]_G-M_{\rho_x(O_y)},\\
  \delta^{\mathrm{prod}}_{x,a}
    &=[\tau_{N,1}(x),M_a]_G-M_{\rho_x(a)},\\
  \delta^{P_0}_{x,a}
    &=\{\tau_{N,1}(x),M_a\}_{P_0}
      -M_{\rho_x^{P_0}(a)} .
  \end{aligned}
\]
Here \(O_y\) is the boundary observable attached to a degree-zero
closed generator, \(M_a\) is multiplication by \(a\), and
\(\rho_x,\rho_x^{P_0}\) are the closed-on-open Hochschild and
\(P_0\)-actions.  A strict centre class is the special case
\(\rho_x=\rho_x^{P_0}=0\); otherwise the datum is equivariance.
The primary primitive equations are
\[
  \begin{aligned}
  d_{\mathrm{cent}}H^{\mathrm{Lie}}_{x,y}
    &=-\delta^{\mathrm{Lie}}_{x,y},&
  d_{\mathrm{cent}}H^{\mathrm{act}}_{x|y}
    &=-\delta^{\mathrm{act}}_{x|y},\\
  d_{\mathrm{cent}}H^{\mathrm{prod}}_{x,a}
    &=-\delta^{\mathrm{prod}}_{x,a},&
  d_{\mathrm{cent}}H^{P_0}_{x,a}
    &=-\delta^{P_0}_{x,a}.
  \end{aligned}
\]
Their signed Stasheff Jacobiator is a secondary curvature
\(\delta^3_{x,y,a}\), killed by
\[
  d_{\mathrm{cent}}H^3_{x,y,a}=-\delta^3_{x,y,a}.
\]
The low-degree obstruction class is
\[
  \operatorname{Ob}^{\mathrm{cent}}_{\mathrm{Hoch},\le3}
  =
  [\delta^{\mathrm{Lie}},
    \delta^{\mathrm{act}},
    \delta^{\mathrm{prod}},
    \delta^{P_0},
    \delta^3]
  \in
  H^2\!\left(F^{\le3}\mathcal K^{\mathrm{cent}}_{\partial,N}\right).
\]
The full obstruction
\(\operatorname{Ob}^{\mathrm{cent}}_{\mathrm{Hoch}}\) is the tower of
these classes in all centrality filtrations.  A compatible system of
primitives satisfying the displayed equations, and commuting with the
arity restrictions, open-input restrictions, finite-window projections,
and Roos transition maps in
\(\mathcal O^{\mathrm{cent}}_{\partial,N}\), is precisely the
Maurer--Cartan datum
\(\tau_N=\mathcal C_{\Phi_N^{\mathrm{Hoch}}}\); the associated
Hochschild cochain action is \(\Phi_N^{\mathrm{Hoch}}\).  Its
coordinate-dual restriction is \(\Phi_N^{\mathrm{gen}}\).  This is the typed
version of
Definition~\ref{defn:admissible-hochschild-centre-lift},
Definition~\ref{defn:low-degree-hochschild-centrality-defects}, and
Proposition~\ref{prop:low-degree-centre-lift-obstruction}.
\end{defn}

\begin{prop}[Boundary open--closed type table]\label{prop:boundary-open-closed-type-table}
The six objects occupy six distinct cells of the type table below,
with pairwise distinct roles; \((C_1)\) is a coalgebra,
\((C_2)\) is a Koszul-dual algebra, \((C_3)\) is a centre carrying an
admissible Hochschild \(E_2\)-Gerstenhaber bracket, unshifted of degree
\(-1\) and shifted of degree \(0\), distinct from any \((C_i)\),
\((C_4)\) is a factorisation algebra compared through supplied
\(\mathrm{OCA}\),
\((C_5)\) is a graded scalar invariant valued in
\(\Z((q,\hbar_{\mathrm{W}}))\), and \((C_6)\) is a separately constructed filtered
\(\mathrm{SC}^{\mathrm{HT}}_{(1,2)}\)-object on \(\partial X\):
\[
  \renewcommand{\arraystretch}{1.25}
  \begin{array}{l|l|l|l}
    \mathsf{Object} & \mathsf{Type} & \mathsf{Role} & \mathsf{Comparison}\\\hline
    \bar B^{E_1}(A^{\mathrm{cl}}_{\partial,N}) &
      \text{conilpotent }E_1^{\mathrm{top}}\text{-coalgebra} &
      \text{classifies twisting morphisms} &
      \Omega^c\bar B\simeq A^{\mathrm{cl}}_{\partial,N}\\
    (A^{\mathrm{cl}}_{\partial,N})^! &
      \text{Verdier/Koszul dual algebra} &
      \text{controls line/open dual modules} &
      \text{Koszul paired with }\bar B\\
    Z^{\mathrm{der},\mathrm{Mor}}_{E_2^{\mathrm{HT}}}(\mathcal C^{\mathrm{op}}_\partial) &
      \text{derived }E_2^{\mathrm{HT}}\text{-centre object} &
      \text{closed centre before physical realisation} &
      \text{target of supplied }\mathrm{OCA}\\
    \mathcal A^{\mathrm{cl}}_{\mathrm{bulk}} &
      \text{mixed-HT BF factorisation algebra} &
      \text{physical bulk if realised} &
      \mathcal A^{\mathrm{cl}}_{\mathrm{bulk}}\xrightarrow{\mathrm{OCA}}Z^{\mathrm{der}}\\
    (\chi,\kappa_{\mathrm{Cap}}) &
      \text{protected index }\times\text{ raw Rees/Weyl Lie }2\text{-class} &
      \text{scalar invariant and raw projective anomaly} &
      \chi\text{ and }\kappa_{\mathrm{Cap}}\text{ lie in distinct targets}\\
    \mathrm{Op}^{\mathrm{HT}}_{\partial X} &
      \text{filtered }\mathrm{SC}^{\mathrm{HT}}\text{-object} &
      \text{boundary operator datum} &
      \operatorname{sTr}(\mathrm{Op}^{\mathrm{HT}}_{\partial X})=\chi
  \end{array}
\]
Each displayed identification factors through the comparison
column.  Bar coalgebras, derived centres, physical factorisation
algebras, scalar traces, and filtered \(\mathrm{SC}^{\mathrm{HT}}\)
objects occupy different ambient categories.
\end{prop}

\begin{proof}
The six cells are pairwise distinct.  \((C_1)\) is a coalgebra by
\eqref{eq:bar-differential-explicit} so cannot coincide with the
algebra \((C_2)\), the centre \((C_3)\), or the bulk \((C_4)\);
\((C_2)\) is the Verdier/Koszul dual algebra (different bracket from
the Hochschild Gerstenhaber bracket of shifted degree \(0\) of \((C_3)\) by
Theorem~\ref{thm:universal-ce-pv-koszul-criterion});
	\((C_3)\) is the derived \(E_2^{\mathrm{HT}}\)-centre object whose
	underlying formal-stalk endomorphism object carries the
	admissible Hochschild \(E_2\)-Gerstenhaber bracket of shifted degree \(0\)
	(Theorem~\ref{thm:main-local}),
and it is distinct from \((C_4)\)'s factorisation product (Definition
\ref{def:closed-mixed-HT-factorization}: \((C_4)\) acts on the
boundary through \(\mathrm{OCA}\), \((C_3)\) is the abstract centre
of the boundary category); \((C_5)\) is a graded scalar
invariant in \(\Z((q,\hbar_{\mathrm{W}}))\) and is non-faithful by
Lemma~\ref{lem:scalar-non-faithfulness}; \((C_6)\) is
a chain-level filtered \(\mathrm{SC}^{\mathrm{HT}}_{(1,2)}\)-object
on \(\partial X\) determined by an independent filtered construction
(Corollary~\ref{cor:mandatory-chain-supplements}), and whose scalar
\(\chi\) is a non-faithful image
(Lemma~\ref{lem:scalar-non-faithfulness}).  Distinct roles are read
off the column entries, and the comparison column records the allowed
maps between the six cells.
\end{proof}

\subsection{The boundary open--closed functor
\texorpdfstring{$\Phi^{\mathrm{HT}}$}{Phi-HT}}
\label{ssec:phi-ht-functor}

\begin{defn}[Boundary open--closed datum]\label{defn:hol-bd-datum-HT}
A \emph{boundary open--closed datum} on the mixed-HT brane line
\(L\subset X\) is a tuple
\[
  \bigl(\bar B,\,\Theta,\,r(z),\,Z^{\mathrm{der}},\,\nabla\bigr),
\]
where \(\bar B\) is a conilpotent dg coalgebra in
\(\mathcal B^{\mathrm{adm}}\); \(\Theta\in
\operatorname{Coder}^1_{\mathrm{cont}}(\bar B)\) is a continuous
degree-one coderivation whose corestriction to cogenerators is the
primitive part of the operation, satisfying the Maurer--Cartan equation
\[
  d_{\operatorname{Coder}}\Theta+\tfrac12[\Theta,\Theta]=0;
\]
when \(\bar B=\bar B^{E_1}(A)\) for an augmented algebra \(A\), \(r(z)\)
is shorthand for a pair consisting of the binary local-cohomology kernel
and its coefficient action,
\[
  r(z)\in
  H^2_{\Delta}\!\bigl(\C[[z,w]]\bigr)\,d\zeta_1d\zeta_2\otimes
  \underline{\mathrm{Hom}}(A^{\widehat\otimes2},A).
\]
Only after this action is fixed does \(r(z)\) determine an arity-two
chiral-collision corestriction in the coderivation algebra;
\(Z^{\mathrm{der}}\) is a derived \(E_2^{\mathrm{HT}}\)-centre object
whose formal-stalk endomorphism object carries a chiral-brace dg
algebra structure.

Let
\[
  G_{\mathrm{Darb}}
  =
  \mathrm{Symp}^{\mathrm{form}}
  (\widehat{\mathbb A}^{2}_{0},dz_1\wedge dz_2)
\]
be the formal symplectomorphism group and let
\(\mathcal T_{\mathrm{Darb}}(b)\) be the torsor of formal Darboux
coordinates at the brane vacuum \(b\).  Its Lie algebra is
\(\mathfrak h=\C[[z_1,z_2]]/\C\cdot1\) with the Poisson bracket.  The
Darboux-transport complex of the bar-centre family is the completed
Chevalley--Eilenberg total complex
\[
  \mathcal E^\bullet_{\mathrm{Darb}}
  =
  C^\bullet_{\mathrm{CE},\mathrm{cont}}\!\left(
  \mathfrak h,\,
  \underline{\mathrm{End}}_{\mathcal B^{\mathrm{adm}}}(\bar B)
  \oplus
  \underline{\mathrm{End}}_{\mathcal B^{\mathrm{adm}}}(Z^{\mathrm{der}})
  \oplus
  \underline{\mathrm{Hom}}_{\mathcal B^{\mathrm{adm}}}(\bar B,Z^{\mathrm{der}})
  \right),
\]
with differential
\[
  d_{\mathrm{Darb}}
  =
  d_{\mathrm{CE}}+d_{\mathrm{int}}+[\Theta,-],
\]
where the last bracket denotes the linearised bar/centre
Maurer--Cartan action on the Hom summand.  A preconnection
\(\nabla_0\) on the completed family over
\(\mathcal T_{\mathrm{Darb}}(b)\) is a degree-zero derivation lifting the
infinitesimal \(G_{\mathrm{Darb}}\)-action and preserving the completed
filtration.  Before flatness is assumed, it determines the two degree-two
obstruction components
\[
  \Omega_{\nabla}:=\nabla_0^2,\qquad
  \mathcal R_{\nabla,\Theta}:=\nabla_0\Theta .
\]
A flat transport datum is the choice of \(\nabla_0\) together with
degree-one primitives
\[
  h_{\nabla},h_{\nabla,\Theta}\in\mathcal E^1_{\mathrm{Darb}},
  \qquad
  d_{\mathrm{Darb}}h_{\nabla}=-\Omega_{\nabla},
  \qquad
  d_{\mathrm{Darb}}h_{\nabla,\Theta}=-\mathcal R_{\nabla,\Theta},
\]
compatible with the Bianchi identities in
\(\mathcal E^\bullet_{\mathrm{Darb}}\).  The Darboux-transport
obstruction is
\[
  \operatorname{Ob}^{\mathrm{Darb}}
  =
  \bigl([\Omega_{\nabla}],[\mathcal R_{\nabla,\Theta}]\bigr)
  \in
  H^2(\mathcal E^\bullet_{\mathrm{Darb}})
  \oplus H^2(\mathcal E^\bullet_{\mathrm{Darb}}).
\]
The formal-Darboux trace theorem supplies the infinitesimal Hamiltonian
action on the trace sector; Proposition~\ref{prop:trace-sector-darboux-flat-transport}
packages this action as a strict flat trace-sector transport datum.  The
bar-centre family requires the separate primitives
\((h_{\nabla},h_{\nabla,\Theta})\) satisfying the two displayed
\(d_{\mathrm{Darb}}\)-equations.
\end{defn}

\begin{prop}[Strict Darboux transport on the trace-sector chart]
\label{prop:trace-sector-darboux-flat-transport}
Let \(\mathfrak h=\C[[z_1,z_2]]/\C\cdot1\) act on the completed
trace-sector source
\[
  G_{\mathrm{tr}}
  =
  \langle c_g,u_g:g\in\mathfrak h\rangle_{\mathrm{CE/PV}}
\]
by
\[
  L^{\mathrm{src}}_f(c_g)=c_{\{f,g\}},
  \qquad
  L^{\mathrm{src}}_f(u_g)=u_{\{f,g\}},
\]
and on the completed trace-sector Hochschild chart
\[
  H^{\mathrm{tr}}_{N,\mathrm{adm}}
  =
  \langle \theta_g,M_{J(g)}:g\in\mathfrak h\rangle
  \subset C^\bullet_{\mathrm{Hoch},\mathrm{adm}}
  (A^{\mathrm{cl}}_{\partial,N},A^{\mathrm{cl}}_{\partial,N})
\]
by
\[
  L^{\mathrm{tar}}_f(\theta_g)=\theta_{\{f,g\}},
  \qquad
  L^{\mathrm{tar}}_f(M_{J(g)})=M_{J(\{f,g\})}.
\]
Both actions extend as continuous derivations for the completed
product and bracket filtrations.  After continuous dualization of
\(\tau_{N,1}^{\mathrm{tr}}\), they define a strict flat
Darboux-transport datum on the coordinate trace-sector chart
\[
  \zeta_{\mathrm{tr}}\colon
  G_{\mathrm{tr}}\longrightarrow H^{\mathrm{tr}}_{N,\mathrm{adm}},
  \qquad
  c_g\mapsto\theta_g,\quad u_g\mapsto M_{J(g)}.
\]
\[
  [L_f,L_g]=L_{\{f,g\}},\qquad
  [d_{\mathrm{tr}},L_f]=0,\qquad
  L^{\mathrm{tar}}_f\zeta_{\mathrm{tr}}
  -
  \zeta_{\mathrm{tr}}L^{\mathrm{src}}_f=0.
\]
Consequently the trace-sector curvature and trace-sector
Maurer--Cartan compatibility classes vanish strictly:
\[
  \Omega^{\mathrm{tr}}_{\nabla}=0,\qquad
  \mathcal R^{\mathrm{tr}}_{\nabla,\Theta}=0,
\]
with zero primitives.  The remaining row in
Definition~\ref{defn:hol-bd-datum-HT} is the extension of this strict
transport from \(G_{\mathrm{tr}}\to H^{\mathrm{tr}}_{N,\mathrm{adm}}\)
to the whole bar-centre family
\(\bar B\to Z^{\mathrm{der}}\).
\end{prop}

\begin{proof}
Continuity is checked on Taylor windows.  The coefficient of
\(\{f,g\}\) modulo \(\mathfrak m^{W+1}\) depends only on the
\((W+1)\)-jets of \(f\) and \(g\), so \(L_f^{\mathrm{src}}\) and
\(L_f^{\mathrm{tar}}\) are continuous for the completed Taylor, stable
trace, and admissible Hochschild filtrations.

The commutator identity is the Jacobi identity for the Poisson bracket:
on \(c_h,u_h,\theta_h\), and \(M_{J(h)}\).  Writing \(x_h\) for any
one of these four generators,
\[
  [L_f,L_g](x_h)
  =
  x_{\{f,\{g,h\}\}-\{g,\{f,h\}\}}
  =
  x_{\{\{f,g\},h\}}
  =
  L_{\{f,g\}}(x_h).
\]
The differentials are equivariant because the CE/PV trace-sector
complex uses the adjoint \(\mathfrak h\)-module structure and the
Koszul differential \(Q\psi=[\phi_1,\phi_2]\) is preserved by
Hamiltonian vector fields:
\[
  \theta_f([\phi_1,\phi_2])
  =
  [\theta_f(\phi_1),\phi_2]+[\phi_1,\theta_f(\phi_2)].
\]
Finally
\[
  L_f^{\mathrm{tar}}\zeta_{\mathrm{tr}}(c_g)
  =
  \theta_{\{f,g\}}
  =
  \zeta_{\mathrm{tr}}L_f^{\mathrm{src}}(c_g),
  \qquad
  L_f^{\mathrm{tar}}\zeta_{\mathrm{tr}}(u_g)
  =
  M_{J(\{f,g\})}
  =
  \zeta_{\mathrm{tr}}L_f^{\mathrm{src}}(u_g).
\]
Thus the trace-sector preconnection has zero curvature and preserves
the trace-sector chart strictly.  No statement about the non-trace
Hochschild summands or the full Morita centre follows from this
calculation.
\end{proof}

\begin{defn}[Relative Darboux extension obstruction]
\label{defn:relative-darboux-extension-obstruction}
Fix the strict trace-sector transport of
Proposition~\ref{prop:trace-sector-darboux-flat-transport}.  The
trace-sector Darboux deformation complex is
\[
  \mathcal E^\bullet_{\mathrm{tr}}
  =
  C^\bullet_{\mathrm{CE},\mathrm{cont}}\!\left(
  \mathfrak h,\,
  \underline{\mathrm{End}}_{\mathcal B^{\mathrm{adm}}}(G_{\mathrm{tr}})
  \oplus
  \underline{\mathrm{End}}_{\mathcal B^{\mathrm{adm}}}
    (H^{\mathrm{tr}}_{N,\mathrm{adm}})
  \oplus
  \underline{\mathrm{Hom}}_{\mathcal B^{\mathrm{adm}}}
    (G_{\mathrm{tr}},H^{\mathrm{tr}}_{N,\mathrm{adm}})
  \right),
\]
with differential \(d_{\mathrm{tr}}\) induced by the source differential,
the Hochschild differential restricted to the strict trace image, and
the linearised action on \(\zeta_{\mathrm{tr}}\).

A trace receiver for the full bar-centre Darboux complex is a continuous
filtered cochain map
\[
  r_{\mathrm{tr}}\colon
  \mathcal E^\bullet_{\mathrm{Darb}}
  \longrightarrow
  \mathcal E^\bullet_{\mathrm{tr}}
\]
whose value on the coordinate-dual trace boundary is the image/coimage trace sector
of Definition~\ref{defn:trace-sector-reduction}, and whose value on a
preconnection extending the trace-sector action is
\((L^{\mathrm{src}},L^{\mathrm{tar}})\) on the two endomorphism summands,
with Hom component
\(L^{\mathrm{tar}}\zeta_{\mathrm{tr}}
-\zeta_{\mathrm{tr}}L^{\mathrm{src}}\).  The existence of
\(r_{\mathrm{tr}}\) is part of the Darboux extension datum; it is not
supplied by the coordinate-dual assignment
\(c_f\mapsto\theta_f,\ u_f\mapsto M_{J(f)}\) alone.

The relative Darboux extension complex is the homotopy fibre
\[
  \mathcal E^\bullet_{\mathrm{Darb}/\mathrm{tr}}
  :=
  \operatorname{Cone}(r_{\mathrm{tr}})[-1].
\]
In the model
\[
  \mathcal E^n_{\mathrm{Darb}/\mathrm{tr}}
  =
  \mathcal E^n_{\mathrm{Darb}}
  \oplus
  \mathcal E^{n-1}_{\mathrm{tr}},
  \qquad
  d_{\mathrm{rel}}(x,y)
  =
  \bigl(d_{\mathrm{Darb}}x,\,-d_{\mathrm{tr}}y-r_{\mathrm{tr}}x\bigr),
\]
a preconnection \(\nabla_0\) extending the strict trace transport has
\[
  r_{\mathrm{tr}}(\Omega_\nabla)=0,\qquad
  r_{\mathrm{tr}}(\mathcal R_{\nabla,\Theta})=0,
\]
and hence defines the relative obstruction pair
\[
  \operatorname{Ob}^{\mathrm{Darb}/\mathrm{tr}}
  =
  \bigl([(\Omega_\nabla,0)],
        [(\mathcal R_{\nabla,\Theta},0)]\bigr)
  \in
  H^2(\mathcal E^\bullet_{\mathrm{Darb}/\mathrm{tr}})
  \oplus
  H^2(\mathcal E^\bullet_{\mathrm{Darb}/\mathrm{tr}}).
\]
\end{defn}

\begin{prop}[Relative criterion for extending trace-sector transport]
\label{prop:relative-darboux-extension-criterion}
Fix a trace receiver \(r_{\mathrm{tr}}\) and a preconnection
\(\nabla_0\) whose restriction is the strict trace-sector transport of
Proposition~\ref{prop:trace-sector-darboux-flat-transport}.  A flat
Darboux transport on the full bar-centre family extending that trace
transport up to trace-sector gauge homotopy exists if and only if
\(\operatorname{Ob}^{\mathrm{Darb}/\mathrm{tr}}=0\).  In this case
relative primitives are pairs
\[
  (h_\nabla,\eta_\nabla),\qquad
  (h_{\nabla,\Theta},\eta_{\nabla,\Theta})
  \in\mathcal E^1_{\mathrm{Darb}/\mathrm{tr}}
\]
satisfying
\[
  d_{\mathrm{rel}}(h_\nabla,\eta_\nabla)
  =
  -(\Omega_\nabla,0),
  \qquad
  d_{\mathrm{rel}}(h_{\nabla,\Theta},\eta_{\nabla,\Theta})
  =
  -(\mathcal R_{\nabla,\Theta},0).
\]
Equivalently,
\[
  d_{\mathrm{Darb}}h_\nabla=-\Omega_\nabla,\qquad
  r_{\mathrm{tr}}(h_\nabla)=-d_{\mathrm{tr}}\eta_\nabla,
\]
and the same equations hold with
\((h_{\nabla,\Theta},\mathcal R_{\nabla,\Theta},
\eta_{\nabla,\Theta})\).  The extension is strict over the zero
trace-sector primitives precisely when one can choose
\(\eta_\nabla=\eta_{\nabla,\Theta}=0\), so
\[
  r_{\mathrm{tr}}(h_\nabla)=0,\qquad
  r_{\mathrm{tr}}(h_{\nabla,\Theta})=0.
\]
\end{prop}

\begin{proof}
The trace-sector curvature and trace-sector
Maurer--Cartan compatibility curvature vanish by
Proposition~\ref{prop:trace-sector-darboux-flat-transport}.  Therefore a
preconnection extending that trace action gives closed degree-two
elements \((\Omega_\nabla,0)\) and
\((\mathcal R_{\nabla,\Theta},0)\) in the homotopy fibre
\(\mathcal E^\bullet_{\mathrm{Darb}/\mathrm{tr}}\).  Unwinding the
differential in the displayed cone model shows that a null-homotopy of
\((\Omega_\nabla,0)\) is exactly a full Darboux primitive
\(h_\nabla\) together with a trace-sector gauge homotopy
\(\eta_\nabla\) from its restriction to zero.  The same calculation
applies to \((\mathcal R_{\nabla,\Theta},0)\).  Vanishing of the two
relative cohomology classes is therefore equivalent to the existence of
the two displayed primitive equations.  Setting the trace homotopies to
zero is exactly the condition that the chosen full primitives restrict
to the zero primitives of the strict trace-sector transport.
\end{proof}

\begin{defn}[Decorated brane and boundary-data categories]
\label{defn:decorated-brane-boundary-data-categories}
The category
\[
  \mathbf{E_1}^{\mathrm{top}}\text{-}
  \mathbf{Brane}^{\mathrm{HT}}_{\R_t,\mathrm{Kosz}}
\]
has as objects tuples
\[
  (A,d_A,\epsilon_A,\alpha_\Delta,\Theta_A,\tau_A,
    \mathcal H_A,\nabla_A).
\]
Here \(A\) is an augmented \(E_1^{\mathrm{top}}\)-algebra in
\(\mathcal B^{\mathrm{adm}}\) of Dirac-brane Koszul form; for the
Dirac \(N\)-brane object,
\[
  A=A^{\mathrm{cl}}_{\partial,N},\qquad Q\psi=[\phi_1,\phi_2].
\]
The map
\[
  \alpha_\Delta\colon
  K^{(2)}_\Delta\widehat\otimes A^{\widehat\otimes2}\longrightarrow A
\]
is the chosen continuous two-variable singular-product coefficient
action, of total degree \(0\) in the local-cohomology convention of
Proposition~\ref{prop:diagonal-cohomology-operad}.  The element
\(\Theta_A\in\operatorname{Coder}^1_{\mathrm{cont}}(\bar B^{E_1}A)\)
satisfies the Maurer--Cartan equation
\[
  d_{\operatorname{Coder}}\Theta_A+\tfrac12[\Theta_A,\Theta_A]=0.
\]
The centre-action datum is a continuous twisting cochain
\[
  \tau_A\colon
  \overline C^{\mathrm{CE},\mathrm{cont}}_\bullet(\mathfrak g)
  \longrightarrow C^\bullet_{\mathrm{Hoch},\mathrm{adm}}(A,A)[1],
  \qquad
  \tau_{A,1}\colon\mathfrak g\longrightarrow
  C^\bullet_{\mathrm{Hoch},\mathrm{adm}}(A,A)[1].
\]
Its trace-sector coordinate dual is the chart
\[
  \zeta_A\colon
  \langle c_f,u_f:f\in\bar A\rangle_{\mathrm{CE/PV}}
  \longrightarrow C^\bullet_{\mathrm{Hoch},\mathrm{adm}}(A,A),
  \qquad
  c_f\mapsto\theta_f,\quad u_f\mapsto M_{J(f)}.
\]
The datum \(\mathcal H_A\) is the family of centrality primitives:
for each finite-window centrality defect
\(\mathfrak D_{\gamma,M}^A\) occurring in the mixed-HT
\(E_2^{\mathrm{HT}}\)-centre
lift, it contains a degree \(-1\) continuous homotopy
\[
  H_{\gamma,M}^A,\qquad
  d_{\mathrm{Hoch}}H_{\gamma,M}^A
  +H_{\gamma,M}^A d_{\mathrm{CE/PV}}
  =-\mathfrak D_{\gamma,M}^A,
\]
compatible with the window transition maps.  Finally
\(\nabla_A\) is a flat transport datum along the formal Darboux
coordinate torsor: either a strict flat connection satisfying
\(\nabla_A^2=0\) and \(\nabla_A\Theta_A=0\), or a preconnection
\(\nabla_{A,0}\) together with primitives killing
\(\Omega_{\nabla}=\nabla_{A,0}^2\) and
\(\mathcal R_{\nabla,\Theta}=\nabla_{A,0}\Theta_A\) in the completed
endomorphism complex.

A morphism
\[
  F\colon
  (A,\alpha_\Delta,\Theta_A,\tau_A,\mathcal H_A,\nabla_A)
  \longrightarrow
  (A',\alpha'_\Delta,\Theta_{A'},\tau_{A'},\mathcal H_{A'},\nabla_{A'})
\]
is an \(E_1^{\mathrm{top}}\)-algebra morphism \(F_A:A\to A'\)
together with continuous comparison homotopies
\[
  h_\Delta,\quad h_\Theta,\quad h_Z,\quad h_{\mathrm{cent}},\quad
  h_\nabla
\]
which solve
\[
\begin{aligned}
  F_A\alpha_\Delta-\alpha'_\Delta(1\otimes F_A^{\otimes2})
    &=d h_\Delta+h_\Delta d,\\
  \bar B(F_A)\Theta_A-\Theta_{A'}\bar B(F_A)
    &=d h_\Theta+h_\Theta d,\\
  Z_F\tau_A-\tau_{A'}
    &=d h_Z+h_Zd,\\
  F_A\mathcal H_A-\mathcal H_{A'}F_A
    &=d h_{\mathrm{cent}}+h_{\mathrm{cent}}d,\\
  \bar B(F_A)\nabla_A-\nabla_{A'}\bar B(F_A)
    &=d h_\nabla+h_\nabla d .
\end{aligned}
\]
Here \(Z_F\) is part of the morphism datum: it is the chosen
Hochschild-centre comparison map induced by \(F_A\) in the admissible
brace model; the displayed comparison induces the corresponding one on
the coordinate-dual charts \(\zeta_A\) and \(\zeta_{A'}\).  Arrows are
taken modulo second homotopy, and composition
is the usual pasting of the five displayed homotopies.

The category \(\mathbf{BdOC}^{\mathrm{HT}}_X\) has objects the boundary
open--closed data of Definition~\ref{defn:hol-bd-datum-HT} equipped
with the same \(\alpha_\Delta,\zeta,\mathcal H\) fields, and has
morphisms defined by the displayed equations with the bar-coalgebra,
centre, and connection components in place of the underlying algebra
component.  Thus an object of either category includes the indicated
primitives killing the singular-product, centre-lift, centrality, and
flat-transport obstruction classes.
\end{defn}

\begin{thm}[Typed boundary open--closed realisation
\texorpdfstring{$\Phi^{\mathrm{HT}}$}{Phi-HT} at \(N\) Dirac branes]
\label{thm:typed-boundary-open-closed-realisation}

In $\mathcal B^{\mathrm{adm}}$, suppose the decorated brane object of
Definition~\ref{defn:decorated-brane-boundary-data-categories} is fixed:
bar--cobar admissibility, the two-variable diagonal-kernel action,
the chiral brace operations, the centrality hierarchy, the trace-sector
centre chart, and the flat connection primitives are part of the
input.  Then the
assignment \eqref{eq:phi-ht-explicit} extends from the
$E_1^{\mathrm{top}}$-brane algebra
$A^{\mathrm{cl}}_{\partial,N}$ on the worldline $\R_t$ (with Koszul
differential $Q\psi=[\phi_1,\phi_2]$) to the corresponding boundary
open--closed datum.

There is a functor
\begin{equation}\label{eq:phi-ht-functor}
  \Phi^{\mathrm{HT}}\colon
  \mathbf{E_1}^{\mathrm{top}}\text{-}\mathbf{Brane}^{\mathrm{HT}}_{\R_t,\mathrm{Kosz}}
    \longrightarrow
  \mathbf{BdOC}^{\mathrm{HT}}_{X}
\end{equation}
from the category of \(E_1^{\mathrm{top}}\)-algebras on \(\R_t\) with
Koszul differential and the displayed HT decorations to the category of
boundary open--closed data of Definition~\ref{defn:hol-bd-datum-HT},
sending \(A^{\mathrm{cl}}_{\partial,N}\) to the boundary open--closed
datum
\begin{equation}\label{eq:phi-ht-explicit}
  \Phi^{\mathrm{HT}}\!\bigl(A^{\mathrm{cl}}_{\partial,N}\bigr)
  =\Bigl(\bar B^{E_1}\!\bigl(A^{\mathrm{cl}}_{\partial,N}\bigr),\,
        \Theta_{A^{\mathrm{cl}}_{\partial,N}},\,
        K^{(2)}_\Delta,\,
        Z^{\mathrm{der},\mathrm{Mor}}_{E_2^{\mathrm{HT}}}\!\bigl(
          \mathcal C^{\mathrm{op}}_\partial\bigr),\,
        \nabla^{\mathrm{HT}}\Bigr),
\end{equation}
where
\begin{enumerate}[label=\textup{(\roman*)}]
  \item \(\bar B^{E_1}(A^{\mathrm{cl}}_{\partial,N})
        =T^{\mathrm{c}}(s^{-1}\bar A^{\mathrm{cl}}_{\partial,N})\) is
        the conilpotent \(E_1\)-bar coalgebra with deconcatenation
        coproduct and differential
\begin{equation}\label{eq:bar-differential-explicit}
          d_{\bar B}=d_{\mathrm{int}}+d_m
\end{equation}
	        where \(d_{\mathrm{int}}\) is the internal differential of
	        \(A^{\mathrm{cl}}_{\partial,N}\) (the Koszul differential
	        \(Q\psi=[\phi_1,\phi_2]\) plus the Chevalley--Eilenberg
	        differential for \(\mathfrak{gl}_N\)).  If
	        \(\pi:T^c(s^{-1}\bar A^{\mathrm{cl}}_{\partial,N})
	        \to s^{-1}\bar A^{\mathrm{cl}}_{\partial,N}\) is projection to
	        cogenerators, \(\bar A=\bar A^{\mathrm{cl}}_{\partial,N}\),
	        and \(s:s^{-1}\bar A\to\bar A\) is suspension,
	        then \(d_m\) is the unique coderivation with
	        \[
	          \pi d_m|_{(s^{-1}\bar A)^{\otimes2}}
	            =
	          s^{-1}\circ m_2\circ(s\otimes s),
	          \qquad
	          \pi d_m|_{(s^{-1}\bar A)^{\otimes r}}=0
	          \quad(r\ne2),
	        \]
	        where the Koszul sign is the one carried by the suspension
	        maps.  The diagonal
	        local-cohomology kernel
	        \(K^{(2)}_\Delta=d\zeta_1d\zeta_2/(\zeta_1\zeta_2)\) of
        Proposition~\ref{prop:diagonal-cohomology-operad} is a
        chiral-kernel decoration of the bar object; it becomes a
        coderivation only after a coefficient action map
        \(K^{(2)}_\Delta\otimes
          (A^{\mathrm{cl}}_{\partial,N})^{\otimes2}\to
          A^{\mathrm{cl}}_{\partial,N}\) is fixed;
  \item \(\Theta_{A^{\mathrm{cl}}_{\partial,N}}\in
        \operatorname{Coder}^1_{\mathrm{cont}}\!\bigl(
          \bar B^{E_1}(A^{\mathrm{cl}}_{\partial,N})\bigr)\) is the
        continuous coderivation Maurer--Cartan element representing the
        \(E_1\)-product on \(A^{\mathrm{cl}}_{\partial,N}\) inside its
        bar.  Its primitive corestriction to cogenerators is
        \[
          \pi\Theta_{A^{\mathrm{cl}}_{\partial,N}}
            =
            \sum_{n\ge 1}
            s^{-1}m_n^{A^{\mathrm{cl}}_{\partial,N}}s^{\otimes n}
              \colon
              (s^{-1}\bar A^{\mathrm{cl}}_{\partial,N})^{\otimes n}
              \to s^{-1}\bar A^{\mathrm{cl}}_{\partial,N},
        \]
        with \(m_2\) the multiplication, \(m_3\) the associator
        homotopy (zero in the strict case), and higher \(m_n\) the
        \(A_\infty\) corrections coming from boundary collision
        contractions; \(\Theta\) satisfies
\begin{equation}\label{eq:mc-bar-equation}
          d_{\operatorname{Coder}}\Theta_{A^{\mathrm{cl}}_{\partial,N}}
            +\tfrac12[\Theta_{A^{\mathrm{cl}}_{\partial,N}},
                       \Theta_{A^{\mathrm{cl}}_{\partial,N}}]=0;
\end{equation}
  \item \(K^{(2)}_\Delta\in
        H^2_\Delta(\C[[z_1,z_2,w_1,w_2]])\,d\zeta_1d\zeta_2\) is the
        scalar binary collision kernel.  It is not a residue of the
        ordinary \(E_1\)-bar Maurer--Cartan element
        \(\Theta_{A^{\mathrm{cl}}_{\partial,N}}\).  After a coefficient
        action map
        \(K^{(2)}_\Delta\otimes
        (A^{\mathrm{cl}}_{\partial,N})^{\otimes2}\to
        A^{\mathrm{cl}}_{\partial,N}\) has been fixed, its
        line/Koszul-dual restriction gives the scalar Arnold--Jacobi
        degree-three relation
\begin{equation}\label{eq:arnold-jacobi-from-mc-degree-three}
          \eta_{12}\wedge\eta_{23}
          +\eta_{23}\wedge\eta_{31}
          +\eta_{31}\wedge\eta_{12}=0,\qquad
          \eta_{ij}=K^{(2)}_\Delta(z_i,z_j),
\end{equation}
        as the scalar kernel-associativity constraint.  No \(r\)-matrix
        equation is asserted: an \(r\)-matrix structure is independent
        Lie-coefficient data, and the scalar diagonal-kernel datum
        contains no such coefficient;
  \item \(Z^{\mathrm{der},\mathrm{Mor}}_{E_2^{\mathrm{HT}}}(
        \mathcal C^{\mathrm{op}}_\partial)\) denotes the derived
        \(E_2^{\mathrm{HT}}\)-centre object.  On the formal-Darboux
        stalk the covariant action is a lifting problem for a continuous
        twisting cochain from the Chevalley--Eilenberg chain coalgebra
        to the shifted admissible Hochschild target
\begin{equation}\label{eq:bar-resolution-derived-centre}
        \begin{gathered}
          H_{N,\mathrm{adm}}^{\mathrm{Hoch}}=
          \RHom_{A^{\mathrm{cl}}_{\partial,N}\otimes
                 (A^{\mathrm{cl}}_{\partial,N})^{\mathrm{op}}}\!\bigl(
	            A^{\mathrm{cl}}_{\partial,N},
	            A^{\mathrm{cl}}_{\partial,N}\bigr)
	          \;\simeq\;
	          C^\bullet_{\mathrm{Hoch},\mathrm{adm}}\!\bigl(
	            A^{\mathrm{cl}}_{\partial,N},
	            A^{\mathrm{cl}}_{\partial,N}\bigr),
          \\
	          \tau_N^{\mathrm{tr}}\colon
	          \overline C^{\mathrm{CE},\mathrm{cont}}_\bullet(\mathfrak g)
	          \longrightarrow H_{N,\mathrm{adm}}^{\mathrm{Hoch}}[1],
	          \qquad
	          \tau_{N,1}^{\mathrm{tr}}\colon
	          \mathfrak g\longrightarrow
	          H_{N,\mathrm{adm}}^{\mathrm{Hoch}}[1],\\
	          G_{\mathrm{tr}}=
	          \langle c_f,u_f:f\in\bar A\rangle_{\mathrm{CE/PV}},
	          \qquad
	          \Phi_N^{\mathrm{gen}}(c_f)=\theta_f,\qquad
	          \Phi_N^{\mathrm{gen}}(u_f)=M_{J(f)}.
        \end{gathered}
\end{equation}
        The last line is obtained from the first Taylor coefficient by
        continuous coordinate duality; it is not the type of the
        twisting cochain itself.
        A full centre identification is the typed datum consisting of
        the admissible centre-lift, equivariance homotopies, mixed-HT
        \(E_2^{\mathrm{HT}}\)-centre operations, and descent data of
        Theorem~\ref{thm:main-global-deligne-criterion}; and
  \item \(\nabla^{\mathrm{HT}}\) is a supplied flat connection on the
        bar fibres over the formal Darboux coordinate torsor at \(b\).
        It is part of the boundary open--closed datum, not a consequence
        of the trace-sector computation.  Its defining equations are
        \[
          (\nabla^{\mathrm{HT}})^2=0,\qquad
          \nabla^{\mathrm{HT}}\Theta_{A^{\mathrm{cl}}_{\partial,N}}=0,
        \]
        together with the induced flat connection on the chosen
        Morita-centre chart
        \(Z^{\mathrm{der},\mathrm{Mor}}_{E_2^{\mathrm{HT}}}
        (\mathcal C^{\mathrm{op}}_\partial)\) through
        \eqref{eq:bar-resolution-derived-centre}.
\end{enumerate}
The projection to \(\bar B^{E_1}(A^{\mathrm{cl}}_{\partial,N})\) is
the bar construction; the projection to
  \((\bar B^{E_1}(A^{\mathrm{cl}}_{\partial,N}),\Theta)\) is the
  Maurer--Cartan twisting datum; the \(K^{(2)}_\Delta\)-component
  is the external binary singular-product kernel; the projection to
\(Z^{\mathrm{der},\mathrm{Mor}}_{E_2^{\mathrm{HT}}}(\mathcal C^{\mathrm{op}}_\partial)\)
is the centre model.  None of these projections produces
the physical bulk \(\mathcal A^{\mathrm{cl}}_{\mathrm{bulk}}\); that
identification is Theorem~\ref{thm:oca-formal-darboux-fibre}
below.
\end{thm}

\begin{proof}
The construction of \eqref{eq:phi-ht-explicit} uses the HT decorations
and the CE/PV trace-sector map of Theorem~\ref{thm:main-local}.

	\emph{(i) Cofree conilpotent bar coalgebra.}
	Set \(\bar A=\bar A^{\mathrm{cl}}_{\partial,N}\).
	\(\bar B^{E_1}(A^{\mathrm{cl}}_{\partial,N})\) is the cofree
	conilpotent \(E_1\)-coalgebra on
	\(s^{-1}\bar A^{\mathrm{cl}}_{\partial,N}\) by the bar/cobar
	adjunction in \(\mathcal B^{\mathrm{adm}}\)
		(Lemma~\ref{lem:continuous-bar-cobar},
			Theorem~\ref{thm:tate-bar-cobar-quillen}).  The internal differential
			\(d_{\mathrm{int}}\) is the Koszul-CE differential of
			\(A^{\mathrm{cl}}_{\partial,N}\), denoted \(d_A\), transported to cogenerators by
		\(\pi d_{\mathrm{int}}s^{-1}a=-s^{-1}(d_Aa)\), with \(\pi\)
		projection to cogenerators.  The term \(d_m\) is
	the coderivation with binary corestriction
	\(s^{-1}m_2(s\otimes s)\) and no other corestriction.  The diagonal
	kernel \(K^{(2)}_\Delta\) enters as a chiral collision coefficient on
	the closed-colour action; an additional action map is required before
	it defines a coderivation of the bar coalgebra.

	\emph{(ii) \(d_{\bar B}^2=0\) and kernel associativity.}
	Since coderivations of \(T^c(s^{-1}\bar A)\) are determined by their
	corestrictions to \(s^{-1}\bar A\), \(d_{\bar B}^2=0\) is checked on
	tensor lengths \(1,2,3\).  The length-one corestriction is
	\(s^{-1}d_A^2s=0\).  The length-two corestriction is the transported
	Leibniz identity
	\[
	  d_A m_2
	    =
	  m_2(d_A\otimes 1+1\otimes d_A).
	\]
	The length-three corestriction is the transported associativity identity
	\[
	  m_2(m_2\otimes1)=m_2(1\otimes m_2).
	\]
	All signs are the suspension signs in the displayed corestrictions.  If
	the boundary product is \(A_\infty\) rather than strict, these three
	equations are replaced by the arity components of the Maurer--Cartan
	equation for \(\Theta\).  The scalar
	chiral-kernel decoration separately satisfies the iterated three-point
relation on \(\Delta_3\subset(\C^2)^3\):
\[
  \eta_{12}\wedge\eta_{23}+\eta_{23}\wedge\eta_{31}
   +\eta_{31}\wedge\eta_{12}=0,\qquad
  \eta_{ij}=K^{(2)}_\Delta(z_i,z_j),
\]
the two-complex-dimensional scalar Arnold coefficient identity on three points, proved as
the residue-form identity at \(\Delta_3\) by
Proposition~\ref{prop:diagonal-cohomology-operad}.  The cross-term
with the Koszul-CE differential is not a formal consequence of this
scalar identity; it belongs to the chosen chiral-kernel action map.
Under that action-map hypothesis, the Koszul homotopy for adjacent
transpositions in traces (Lemma~\ref{lem:koszul-adjacent-transpositions})
gives the required compatibility.  Hence \(d_{\bar B}^2=0\) for the
ordinary bar coalgebra.  The singular-product residue data do not live
as elements of this coalgebra; after the coefficient action is fixed,
they give an arity-two primitive corestriction in the completed
coderivation dg Lie algebra
\(\operatorname{Coder}_{\mathrm{cont}}(\bar B^{E_1}
(A^{\mathrm{cl}}_{\partial,N}))\), with scalar coefficient constrained
by \eqref{eq:arnold-jacobi-from-mc-degree-three}.

\emph{(iii) Binary collision residue equals \(K^{(2)}_\Delta\).}
The binary collision residue of \(\Theta\) at \(\Delta\subset
\C^2\times\C^2\) is the arity-two local-cohomology singular coefficient
\(K^{(2)}_\Delta=d\zeta_1d\zeta_2/(\zeta_1\zeta_2)\) of
\eqref{eq:singular-product-linear}.  The degree-three component of
\eqref{eq:mc-bar-equation}, evaluated on three free generators
\(x_1,x_2,x_3\) of \(\bar A^{\mathrm{cl}}_{\partial,N}\), produces the
scalar Arnold--Jacobi coefficient identity
\eqref{eq:arnold-jacobi-from-mc-degree-three}.  This is the scalar
operadic associativity constraint for the diagonal kernel.  It is not an
\(r\)-matrix equation; an \(r\)-matrix structure is independent
Lie-coefficient data.

\emph{(iv) Derived centre via bar resolution.}
The underlying bimodule endomorphism object of the derived
\(E_2^{\mathrm{HT}}\)-centre on the formal-Darboux stalk
identifies with the chiral Hochschild cochain brace dg algebra by the
bar resolution
\(\Omega^{\mathrm{c}}\bar B^{E_1}(A^{\mathrm{cl}}_{\partial,N})
\to A^{\mathrm{cl}}_{\partial,N}\), the admissible Tate bar--cobar
theorem (Theorem~\ref{thm:tate-bar-cobar-quillen}), the universal
filtered cobar recognition criterion
(Theorem~\ref{thm:universal-filtered-cobar-recognition}), and the
\(E_2\)-Gerstenhaber bracket match of shifted degree \(0\) in
\eqref{eq:ce-pv-dictionary-strong}.  Equation
\eqref{eq:bar-resolution-derived-centre} is the local form of the
trace-sector first Taylor coefficient and its coordinate-dual expression
in Hochschild cochains
		\(\operatorname{Obs}^{\mathrm{cl}}_b(\mathcal A^{\mathrm{cl}}_{\mathrm{bulk}})\to
		C^\bullet_{\mathrm{Hoch},\mathrm{adm}}\!\bigl(A^{\mathrm{cl}}_{\partial,N},
		A^{\mathrm{cl}}_{\partial,N}\bigr)\) of
Theorem~\ref{thm:main-local}; a global derived-centre quasi-isomorphism
is governed by the typed data of \eqref{eq:typed-total-obstruction-ledger},
including the Hochschild-centre lift, descent/QME/anomaly/locality rows,
unreduced cotangent-current row when the full BF source is used, and
comparison-cone contraction.

\emph{(v) Flat transport compatibility.}
The connection component is the separate flat transport datum included in
Definition~\ref{defn:hol-bd-datum-HT}.  The required equations are
\[
  (\nabla^{\mathrm{HT}})^2=0,\qquad
  \nabla^{\mathrm{HT}}\Theta_{A^{\mathrm{cl}}_{\partial,N}}=0
\]
in the completed endomorphism complex of the bar family, and the same
equations after applying the chosen centre chart
\eqref{eq:bar-resolution-derived-centre}.  If one starts from a
connection \(\nabla_0\), the obstructions are
\(\Omega_{\nabla}=\nabla_0^2\) and
\(\mathcal R_{\nabla,\Theta}=\nabla_0\Theta\); their primitives are
additional data of the boundary open--closed object.  The local
formal-Darboux theorem supplies the Hamiltonian infinitesimal action on
the trace-sector chart.  The extension from this trace-sector action to
flat transport on the full bar family is exactly the vanishing of
\(\Omega_{\nabla}\) and \(\mathcal R_{\nabla,\Theta}\), together with
chosen primitives in the completed endomorphism complex.

\emph{(Functoriality.)}
A morphism in
\(\mathbf{E_1}^{\mathrm{top}}\text{-}\mathbf{Brane}^{\mathrm{HT}}_{\R_t,\mathrm{Kosz}}\)
is an \(E_1^{\mathrm{top}}\)-morphism
\(f\colon A^{\mathrm{cl}}_{\partial,N}\to A^{\mathrm{cl}}_{\partial,N'}\)
together with homotopies identifying the pulled-back diagonal residue,
centre lift, centrality hierarchy, and flat transport connection.  The bar
functor sends \(f\) to a coalgebra morphism
\(\bar B(f)\colon\bar B^{E_1}(A^{\mathrm{cl}}_{\partial,N})\to
\bar B^{E_1}(A^{\mathrm{cl}}_{\partial,N'})\), and the chosen
homotopies make it a morphism of boundary open--closed data.  Thus
\(\Phi^{\mathrm{HT}}\) is a functor on the decorated category.  The
decoration consists of the chiral \(E_2\), centrality, descent, and
connection data attached to the underlying \(E_1\)-morphism.
\end{proof}

\subsection{The open--closed comparison morphism}
\label{ssec:oca}

\begin{defn}[Open--closed comparison datum]\label{defn:oca-comparison}
At \(N\) Dirac branes, and relative to a supplied
\(\mathrm{SC}^{\mathrm{HT}}_{(1,2)}\)-action datum, a
\emph{pre-open--closed comparison datum} is the bulk-to-boundary
morphism
\begin{equation}\label{eq:pre-oca-definition}
  \mathrm{preOCA}_N=
  \widetilde{\rho}_{\mathrm{bulk}\to\partial}\colon
    \mathcal A^{\mathrm{cl}}_{\mathrm{bulk}}
    \;\longrightarrow\;
    \underline{\mathrm{End}}_{\mathrm{SC}^{\mathrm{HT}}_{(1,2)}}\!
      \bigl(\mathcal C^{\mathrm{op}}_\partial\bigr)
\end{equation}
in \(\mathcal B^{\mathrm{adm}}\), together with its formal-Darboux
trace-sector fibre.  If a chiral \(E_2\) centrality morphism
\[
  \iota_{\mathrm{cent}}\colon
    \underline{\mathrm{End}}_{\mathrm{SC}^{\mathrm{HT}}_{(1,2)}}\!
      \bigl(\mathcal C^{\mathrm{op}}_\partial\bigr)
    \longrightarrow
    Z^{\mathrm{der},\mathrm{Mor}}_{E_2^{\mathrm{HT}}}\!\bigl(
      \mathcal C^{\mathrm{op}}_\partial\bigr)
\]
and the required \(\mathrm{SC}^{\mathrm{HT}}_{(1,2)}\)-centrality
homotopies are chosen, the datum lifts to the
\emph{open--closed comparison morphism}
\begin{equation}\label{eq:oca-definition}
  \mathrm{OCA}_N\colon
    \mathcal A^{\mathrm{cl}}_{\mathrm{bulk}}
    \;\longrightarrow\;
    Z^{\mathrm{der},\mathrm{Mor}}_{E_2^{\mathrm{HT}}}\!\bigl(
      \mathcal C^{\mathrm{op}}_\partial\bigr),
  \qquad
  \mathrm{OCA}_N
   =
   \iota_{\mathrm{cent}}\circ\mathrm{preOCA}_N .
\end{equation}
Here
\begin{enumerate}[label=\textup{(\roman*)}]
  \item \(\widetilde{\rho}_{\mathrm{bulk}\to\partial}\) is the supplied
        mixed-HT bulk--boundary factorisation pairing, written in the
        Costello--Gwilliam factorization vocabulary of
        \cite{costello-gwilliam}*{Vol.~I Secs.~6.1, 7.2}, adjointed to the
        endomorphism algebra of \(\mathcal C^{\mathrm{op}}_\partial\)
        as an \(\mathrm{SC}^{\mathrm{HT}}_{(1,2)}\)-module: for nested
        \(U\supset V\subset\partial X\), the pairing
        \(\mathcal A^{\mathrm{cl}}_{\mathrm{bulk}}(U)\otimes
        \mathcal C^{\mathrm{op}}_\partial(V)\to
        \mathcal C^{\mathrm{op}}_\partial(V)\) is the closed-acting-on-open
        operation of Theorem~\ref{thm:sc-ht-deligne-swiss-cheese} item~(iii),
        and \(\widetilde\rho\) is its currying to the endomorphism algebra;
  \item \(\iota_{\mathrm{cent}}\) is the centrality morphism into the
        derived centre object.  Its universal property is part of the
        mixed-HT \(E_2^{\mathrm{HT}}\)-centrality datum; the formal-Darboux trace
        theorem gives its trace-sector coordinate.
\end{enumerate}
With these choices the composition is a morphism of
\(\mathrm{SC}^{\mathrm{HT}}_{(1,2)}\)-algebras on \(\partial X\).  The
additional assertion that this morphism is the open--closed
correspondence quasi-isomorphism is the comparison-cone datum of
Definition~\ref{defn:oca-comparison-cone}.
\end{defn}

\begin{defn}[Open--closed comparison cone]
\label{defn:oca-comparison-cone}
For an open--closed comparison morphism
\[
  \mathrm{OCA}_N\colon
  \mathcal A^{\mathrm{cl}}_{\mathrm{bulk}}
  \longrightarrow
  Z_\partial,\qquad
  Z_\partial=
  Z^{\mathrm{der},\mathrm{Mor}}_{E_2^{\mathrm{HT}}}
  (\mathcal C^{\mathrm{op}}_\partial),
\]
the \emph{open--closed comparison cone} is the object
\[
  \operatorname{Cone}(\mathrm{OCA}_N)
  =
  Z_\partial\oplus
  \mathcal A^{\mathrm{cl}}_{\mathrm{bulk}}[1]
\]
of \(\mathcal B^{\mathrm{adm}}\), with differential
\[
  d_{\operatorname{Cone}}(z,a)
  =
  \bigl(d_{Z_\partial}z+\mathrm{OCA}_N(a),\,
        -d_{\mathcal A}a\bigr).
\]
The comparison-cone obstruction is the cohomology object
\[
  \operatorname{Ob}_{\operatorname{Cone}}
  =
  H^\bullet\!\left(\operatorname{Cone}(\mathrm{OCA}_N)\right).
\]
A contracting homotopy is a degree \(-1\) endomorphism
\(\kappa\) satisfying
\[
  d_{\operatorname{Cone}}\kappa+\kappa d_{\operatorname{Cone}}
  =\operatorname{id}_{\operatorname{Cone}}.
\]
The equation \(\operatorname{Ob}_{\operatorname{Cone}}=0\), together
with such a \(\kappa\), is the datum that \(\mathrm{OCA}_N\) is a
quasi-isomorphism in \(\mathcal B^{\mathrm{adm}}\).  It is not part of
the definition of \(\mathrm{OCA}_N\).
\end{defn}

\begin{lem}[Complete filtered contraction lift]
\label{lem:complete-filtered-contraction-lift}
Let \(C\) be a complete separated descending filtered cochain complex
with differential \(d\), and assume that the filtration is bounded
below in each cohomological degree.  Let
\[
  h\colon C\longrightarrow C[-1]
\]
be a continuous filtration-preserving degree \(-1\) endomorphism.  If
the contraction defect
\[
  E=\operatorname{id}_C-(dh+hd)
\]
satisfies \(E(F^rC)\subset F^{r+1}C\) for every \(r\), then
\[
  H=h\sum_{n\ge0}E^n
\]
is a continuous contracting homotopy of \(C\):
\[
  dH+Hd=\operatorname{id}_C .
\]
\end{lem}

\begin{proof}
The operator \(E\) has degree zero and raises filtration by one.  Hence
for each cohomological degree the sequence of partial sums
\(\sum_{0\le n\le M}E^n\) is Cauchy in the complete filtered topology,
and the limit \(S=\sum_{n\ge0}E^n\) is continuous.  Since
\[
  E=\operatorname{id}_C-[d,h],
  \qquad [d,h]=dh+hd,
\]
and \(d^2=0\), the graded Jacobi identity gives
\([d,E]=-[d,[d,h]]=0\).  Thus \(E\), and therefore \(S\), commute with
\(d\).  The filtration-raising condition gives
\(\lim_{M\to\infty}E^{M+1}=0\), so
\[
  (\operatorname{id}_C-E)S
  =
  \lim_{M\to\infty}
  (\operatorname{id}_C-E)\sum_{0\le n\le M}E^n
  =
  \operatorname{id}_C .
\]
Consequently
\[
  d(hS)+(hS)d
  =
  (dh+hd)S
  =
  (\operatorname{id}_C-E)S
  =
  \operatorname{id}_C .
\]
\end{proof}

\begin{thm}[Filtered comparison-cone criterion]
\label{thm:filtered-oca-cone-criterion}
Let \(\mathrm{OCA}_N\colon\mathcal A\to Z_\partial\) be an
open--closed comparison morphism in \(\mathcal B^{\mathrm{adm}}\).
Assume \(\mathcal A\) and \(Z_\partial\) carry complete, separated,
descending filtrations \(F^\bullet\) such that
\[
  d_{\mathcal A}F^r\mathcal A\subset F^r\mathcal A,\qquad
  d_{Z_\partial}F^rZ_\partial\subset F^rZ_\partial,\qquad
  \mathrm{OCA}_N(F^r\mathcal A)\subset F^rZ_\partial ,
\]
and the filtration is bounded below in each cohomological degree.  Put
\[
  F^r\operatorname{Cone}(\mathrm{OCA}_N)
  =
  F^rZ_\partial\oplus F^r\mathcal A[1].
\]
If the associated-graded cone
\[
  \operatorname{gr}\operatorname{Cone}(\mathrm{OCA}_N)
  =
  \operatorname{Cone}\bigl(
    \operatorname{gr}\mathrm{OCA}_N\colon
    \operatorname{gr}\mathcal A\to\operatorname{gr}Z_\partial\bigr)
\]
is acyclic, then \(\operatorname{Cone}(\mathrm{OCA}_N)\) is acyclic.

If, in addition, \(\operatorname{gr}\operatorname{Cone}(\mathrm{OCA}_N)\)
has a degree \(-1\) contracting homotopy \(\kappa_0\) whose components
lift to a continuous filtration-preserving endomorphism
\(\widetilde\kappa_0\) of
\(\operatorname{Cone}(\mathrm{OCA}_N)\), and if its contraction defect
\[
  E_{\widetilde\kappa}
  =
  \operatorname{id}-
  (d_{\operatorname{Cone}}\widetilde\kappa_0
   +\widetilde\kappa_0d_{\operatorname{Cone}})
\]
satisfies
\[
  E_{\widetilde\kappa}
  \bigl(F^r\operatorname{Cone}(\mathrm{OCA}_N)\bigr)
  \subset
  F^{r+1}\operatorname{Cone}(\mathrm{OCA}_N)
  \qquad\text{for every }r,
\]
then there is a continuous contracting homotopy
\[
  \kappa\colon
  \operatorname{Cone}(\mathrm{OCA}_N)
  \longrightarrow
  \operatorname{Cone}(\mathrm{OCA}_N)[-1]
  \]
with
\[
  d_{\operatorname{Cone}}\kappa+\kappa d_{\operatorname{Cone}}
  =\operatorname{id}_{\operatorname{Cone}} .
\]
Thus the filtered cone data above supply the quasi-isomorphism datum for
the open--closed comparison.
\end{thm}

\begin{proof}
The cone filtration is complete and separated because the two summands
are complete and separated in \(\mathcal B^{\mathrm{adm}}\), and
\(\mathrm{OCA}_N\) is filtered.  The spectral sequence of the filtered
cone has
\[
  E_1\cong
  H^\bullet\!\left(
    \operatorname{gr}\operatorname{Cone}(\mathrm{OCA}_N)
  \right).
\]
The bounded-below hypothesis in each cohomological degree gives strong
convergence to
\(H^\bullet(\operatorname{Cone}(\mathrm{OCA}_N))\).  If the
associated-graded cone is acyclic, \(E_1=0\), hence the cone cohomology
vanishes.

For the homotopy statement, write the cone differential as
\(d=d_0+\delta\), where \(d_0\) is a filtered lift of the
associated-graded differential and \(\delta(F^r)\subset F^{r+1}\).
Let \(\widetilde\kappa_0\) be a continuous lift of \(\kappa_0\).  On the
associated graded,
\[
  d_0\widetilde\kappa_0+\widetilde\kappa_0d_0=\operatorname{id}
  \quad\text{mod }F^1.
\]
The error
\[
  E=\operatorname{id}-
  (d\widetilde\kappa_0+\widetilde\kappa_0d)
\]
raises filtration by at least one.  Lemma~\ref{lem:complete-filtered-contraction-lift}
applied to \(C=\operatorname{Cone}(\mathrm{OCA}_N)\) and
\(h=\widetilde\kappa_0\) gives the continuous contracting homotopy
\[
  \kappa=\widetilde\kappa_0\sum_{n\ge0}E^n .
\]
\end{proof}

\begin{defn}[Brane-link descent and comparison-cone datum]
\label{defn:typed-brane-link-descent-cone-datum}
Choose a collared formal-Darboux Weiss cover
\(\mathfrak U=\{U_a\}_{a\in A}\) of the brane-link boundary
\(\partial X\), in the sense of
Definition~\ref{app:defn:collared-darboux-weiss-cover}.  Put
\[
  Z_\partial=
  Z^{\mathrm{der},\mathrm{Mor}}_{E_2^{\mathrm{HT}}}
  (\mathcal C^{\mathrm{op}}_\partial),
  \qquad
  \mathcal A=\mathcal A^{\mathrm{cl}/\hbar}_{\mathrm{bulk}} .
\]
The Costello--Gwilliam descent datum used here is the pair of Weiss
cosheaf diagrams
\(U\mapsto \mathcal A|_U\) and \(U\mapsto Z_\partial|_U\) in
\(\mathcal B^{\mathrm{adm}}\).  The complex below is the internal Hom
complex between their restrictions on the Weiss nerve of
\(\mathfrak U\).
For \(U_{a_0\cdots a_p}=U_{a_0}\cap\cdots\cap U_{a_p}\), let
\[
  \mathcal E^\bullet_{\mathrm{desc}}(U_{a_0\cdots a_p})
  =
  \underline{\mathrm{Hom}}^\bullet_{\mathcal B^{\mathrm{adm}}}
  \bigl(\mathcal A|_{U_{a_0\cdots a_p}},
        Z_\partial|_{U_{a_0\cdots a_p}}\bigr),
  \qquad
  d_{\mathcal E}\varphi
  =
  d_{Z_\partial}\varphi-(-1)^{|\varphi|}\varphi d_{\mathcal A}.
\]
The descent complex is
\[
  \operatorname{Tot}^m\check C^\bullet_{\mathrm{Weiss}}
  (\mathfrak U;\mathcal E^\bullet_{\mathrm{desc}})
  =
  \prod_{p+q=m}
  \check C^p_{\mathrm{Weiss}}
  (\mathfrak U;\mathcal E^q_{\mathrm{desc}}),
  \qquad
  d_{\mathrm{desc}}=\delta_{\check C}+(-1)^p d_{\mathcal E}
  \text{ on }\check C^p .
\]

A brane-link descent datum for local comparison maps is a total
degree-zero cocycle
\[
  \mathfrak F_{\mathrm{desc}}
  =
  F^{(0)}+H^{(1)}+H^{(2)}+\cdots
  \in
  \operatorname{Tot}^0\check C^\bullet_{\mathrm{Weiss}}
  (\mathfrak U;\mathcal E^\bullet_{\mathrm{desc}}),
  \qquad
  d_{\mathrm{desc}}\mathfrak F_{\mathrm{desc}}=0,
\]
whose \(p=0\) component is the family of local mixed-HT
closed-to-open maps
\[
  F^{(0)}_a=
  \mathrm{OCA}_{N,a}\colon
  \mathcal A|_{U_a}\longrightarrow Z_\partial|_{U_a}.
\]
The \(p=1\) component consists of overlap homotopies
\[
  H^{(1)}_{ab}\in
  \mathcal E^{-1}_{\mathrm{desc}}(U_{ab}),\qquad
  d_{\mathcal E}H^{(1)}_{ab}
  =
  F^{(0)}_b|_{U_{ab}}-F^{(0)}_a|_{U_{ab}},
\]
and the \(p=2\) component kills the triple-overlap cocycle
\[
  (\delta_{\check C}H^{(1)})_{abc}
  =
  H^{(1)}_{bc}-H^{(1)}_{ac}+H^{(1)}_{ab}.
\]
Let \(F^r_{\check C}\operatorname{Tot}\) denote the descending
\v{C}ech-filtration subcomplex obtained by retaining only components of
\v{C}ech degree \(p\ge r\).  If the datum has been chosen through
\v{C}ech degree \(k\), its next descent curvature is
\[
  \mathcal R^{\mathrm{desc}}_{k+1}
  =
  \bigl(d_{\mathrm{desc}}
  (F^{(0)}+H^{(1)}+\cdots+H^{(k)})\bigr)_{k+1}
  \in
  \check C^{k+1}_{\mathrm{Weiss}}
  (\mathfrak U;\mathcal E^{-k}_{\mathrm{desc}}).
\]
It is a cycle in the associated graded complex
\[
  \operatorname{gr}^{k+1}_{F_{\check C}}
  \operatorname{Tot}\check C^\bullet_{\mathrm{Weiss}}
  (\mathfrak U;\mathcal E^\bullet_{\mathrm{desc}}),
\]
and defines the stage-\((k+1)\) descent obstruction
\[
  \operatorname{Ob}^{\mathrm{desc}}_{k+1}
  =
  [\mathcal R^{\mathrm{desc}}_{k+1}]
  \in
  H^1\!\left(
  \operatorname{gr}^{k+1}_{F_{\check C}}
  \operatorname{Tot}\check C^\bullet_{\mathrm{Weiss}}
  (\mathfrak U;\mathcal E^\bullet_{\mathrm{desc}})
  \right).
\]
A primitive \(H^{(k+1)}\) satisfies
\[
  d_{\mathrm{desc}}H^{(k+1)}
  =-\mathcal R^{\mathrm{desc}}_{k+1}
\]
with
\[
  H^{(k+1)}\in
  \check C^{k+1}_{\mathrm{Weiss}}
  (\mathfrak U;\mathcal E^{-k-1}_{\mathrm{desc}}).
\]
The total descent obstruction
\(\operatorname{Ob}^{\mathrm{desc}}_{\partial}\) is the compatible tower
of the classes \(\operatorname{Ob}^{\mathrm{desc}}_{k+1}\), including
the triple-overlap class at \(k+1=2\).  This is the body-level form of
the obstruction complex of
Definition~\ref{app:defn:descent-obstruction-complex}.

After the cocycle \(\mathfrak F_{\mathrm{desc}}\) descends to a global
map \(\mathrm{OCA}_N\), a comparison-cone datum is the filtered datum
\[
  \bigl(
  F^\bullet\operatorname{Cone}(\mathrm{OCA}_N),
  \kappa_0,\widetilde\kappa_0,E
  \bigr)
\]
of Theorem~\ref{thm:filtered-oca-cone-criterion}: \(\kappa_0\) is a
degree \(-1\) contraction of the associated-graded cone,
\(\widetilde\kappa_0\) is a continuous lift, and \(E\) raises
filtration so that \(\sum_{n\ge0}E^n\) converges.  The resulting
homotopy
\[
  \kappa=\widetilde\kappa_0\sum_{n\ge0}E^n
\]
contracts \(\operatorname{Cone}(\mathrm{OCA}_N)\).  Thus descent
produces the global morphism, while the cone datum proves it is a
quasi-isomorphism.
\end{defn}

\begin{thm}[\(\mathrm{OCA}\) at the formal-Darboux stalk]
\label{thm:oca-formal-darboux-fibre}

This is Theorem~\ref{thm:main-local} viewed as the formal-Darboux
trace-sector fibre of an open--closed comparison datum.  In
\(\mathcal B^{\mathrm{adm}}\) its covariant boundary value is the
continuous first Taylor coefficient
\[
  \tau_{N,1}^{\mathrm{tr}}\colon
  \mathfrak g\longrightarrow
  C^\bullet_{\mathrm{Hoch},\mathrm{adm}}\!
    (A^{\mathrm{cl}}_{\partial,N},A^{\mathrm{cl}}_{\partial,N})[1],
\]
the arity-one part of a candidate continuous twisting cochain on
\(\overline C^{\mathrm{CE},\mathrm{cont}}_\bullet(\mathfrak g)\).
Continuous coordinate duality then gives
\[
  c_f\longmapsto\theta_f,\qquad
  u_f\longmapsto J(f),
\]
with \(J(f)\) represented in Hochschild cochains by multiplication
\(M_{J(f)}\).
The degrees are \(|c_f|=|\theta_f|=1\) and \(|u_f|=|J(f)|=0\).
After stable scalar reduction and admissible completion this
	  restriction is the CE/PV coordinate isomorphism.  At fixed finite \(N\) it is not
		the full derived-centre quasi-isomorphism; that assertion is the typed
	ledger problem of \eqref{eq:typed-total-obstruction-ledger}.  Its
	formal-Darboux coordinate
	is \(J(f)=\overline{\operatorname{Tr}\,f(\phi_1,\phi_2)}\) for
polynomial \(f\), extended to formal \(f\) only after completion, with
raw Rees/Weyl projective trace-commutator class
\(\hbar_{\mathrm{W}}\chi_N(K)[\bar c]
=\hbar_{\mathrm{W}}N[\bar c]\).  The underlying
constant-Hamiltonian extension remains the unscaled class \([\bar c]\).
\end{thm}

\begin{proof}
The pre-open--closed datum of Definition~\ref{defn:oca-comparison}, evaluated
at the formal-Darboux stalk \(b\), contains the bulk-boundary factorisation
pairing on the contractible polydisk neighbourhood of \(b\).  On the
trace sector Theorem~\ref{thm:main-local} computes this pairing by the
first Taylor coefficient \(\tau_{N,1}^{\mathrm{tr}}\); its
coordinate-dual CE/PV expression is
\(c_f\mapsto\theta_f,\ u_f\mapsto J(f)\).  The binary diagonal kernel,
the unscaled constant-Hamiltonian cocycle, and its raw Rees/Weyl trace
pushout are local data.  The all-arity
mixed-HT \(E_2^{\mathrm{HT}}\)-centre extension, descent, and all-order
QME counterterm tower are the separate extension rows of
Definition~\ref{defn:typed-mixed-ht-deligne-equivalence-datum}.
\end{proof}

\subsection{The mixed-HT chiral Deligne--Swiss-cheese action}
\label{ssec:sc-ht-deligne}

The normal-defect geometry is encoded by a two-coloured mixed-HT
Swiss-cheese operadic model.  The open colour is
\(E_1^{\mathrm{top}}\) along the brane line \(\R_t\subset L\).  The
normal closed colour is the
mixed holomorphic-topological factorisation operad on
\(\R_s\times\C^2\).  The full bulk colour on
\(\R_t\times\R_s\times\C^2\) is the tensor of this normal operad with
the tangential \(E_1^{\mathrm{top}}\) factor, equipped with the mixed
coherence maps.

\begin{defn}[Normal mixed-HT Swiss-cheese operad]
\label{defn:sc-ht-operad}
The \emph{normal mixed-HT Swiss-cheese operad of bidegree \((1,2)\)},
denoted \(\mathrm{SC}^{\mathrm{HT}}_{(1,2)}\), is the two-coloured
coloured dg operad with:
\begin{itemize}
  \item closed colour \(\mathsf c\): the \(\mathrm{HT}^{(1,2)}\)
        framed little disks operad on \(\R^1_{\mathrm{top}}\times\C^2_{\mathrm{hol}}\),
        equivalently \(E_1^{\mathrm{top}}\widehat\otimes E_2^{\mathrm{ch,hol}}\),
        modelled on the framed Fulton--MacPherson compactification of
        configurations of disks in
        \(\R^1_{\mathrm{top}}\times\C^2_{\mathrm{hol}}\);
  \item open colour \(\mathsf o\): the \(E_1^{\mathrm{top}}\) little
        intervals operad on the brane line \(\R_t\);
  \item mixed operations recording configurations of normal closed disks
        in \(\R_{s\ge 0}\times\C^2_{\mathrm{hol}}\), parametrised along
        \(\R_t\), approaching the boundary
        \(\R_t\times\{s=0\}\times\C^2_{\mathrm{hol}}\), with marked open
        intervals on the brane line
        \(\R_t\times\{s=z_1=z_2=0\}\), with composition by nesting
        closed disks and substitution of open intervals.
\end{itemize}
This is the geometric two-coloured operad.  The formal
Hamiltonian-current representation used below is additional data: its
closed arity-two restriction on a formal Darboux chart is represented by
the diagonal local-cohomology kernel \(K^{(2)}_\Delta\) on \(\C^2\), its
open arity-two restriction is the \(E_1\)-product on the brane line, and
its mixed arity-two restriction is the bulk-to-boundary insertion.  The
higher closed and mixed operations are not determined by
\(K^{(2)}_\Delta\) alone; their compatibility is the Maurer--Cartan
lifting problem of Definition~\ref{defn:sc-ht-action-datum}.  A
comparison with a topological \(E_3\) Swiss-cheese object requires the
topologisation datum of
Theorem~\ref{thm:topologisation-by-stress-tensor}.
\end{defn}

\begin{rmk}[Dunn additivity and mixed holomorphicity]\label{rmk:no-dunn-collapse}
Dunn additivity \cite{lurie-ha}*{Theorem~5.1.2.2} identifies
\(E_p^{\mathrm{top}}\otimes E_q^{\mathrm{top}}\simeq E_{p+q}^{\mathrm{top}}\)
in topological algebras.  The mixed-HT operad
\(E_1^{\mathrm{top}}\otimes E_2^{\mathrm{ch,hol}}\) is the tensor of a
topological \(E_1\) factor with a holomorphic/chiral
\(E_2^{\mathrm{ch,hol}}\) factor.  The passage from
\(E_2^{\mathrm{ch,hol}}\) to \(E_2^{\mathrm{top}}\) requires a
supplied topologisation datum \(T_i=[Q,G_i]\) and the vanishing
obstruction tower of Theorem~\ref{thm:topologisation-by-stress-tensor}.
Thus \(\mathrm{SC}^{\mathrm{HT}}_{(1,2)}\) is a two-coloured mixed-HT
operadic model; its \(E_3^{\mathrm{top}}\) form exists only after that
transfer datum is fixed.
\end{rmk}

\begin{defn}[Mixed-HT Swiss-cheese action datum]
\label{defn:sc-ht-action-datum}
Put
\[
  Z_\partial=
  Z^{\mathrm{der},\mathrm{Mor}}_{E_2^{\mathrm{HT}}}\!
  \bigl(\mathcal C^{\mathrm{op}}_\partial\bigr).
\]
Let \(S\) be a proposed closed-colour object in
\(\mathcal B^{\mathrm{adm}}\).  The two cases used below are
\[
  S=Z_\partial,\qquad
  S=\mathcal A^{\mathrm{cl}/\hbar}_{\mathrm{bulk}} .
\]
The deformation complex controlling a mixed-HT Swiss-cheese action on
the pair \((S,\mathcal C^{\mathrm{op}}_\partial)\) is the complete
arity-filtered operadic convolution dg Lie algebra
\[
  \mathfrak{Def}^{\mathrm{HT}}_{\mathrm{SC}}
  (S,\mathcal C^{\mathrm{op}}_\partial)
  =
  \prod_{\mathsf x\in\{\mathsf c,\mathsf o\}}
  \prod_{\substack{r,s\ge0\\ r+s\ge1}}
  \underline{\mathrm{Hom}}_{\Sigma_r\times\Sigma_s}\!\left(
    \overline{\mathrm{SC}}^{\mathrm{HT}}_{(1,2)}
    (\mathsf c^{\,r},\mathsf o^{\,s};\mathsf x),\,
    \mathrm{End}_{(S,\mathcal C^{\mathrm{op}}_\partial)}
    (\mathsf c^{\,r},\mathsf o^{\,s};\mathsf x)\right),
\]
with differential induced by the differentials of
\(\overline{\mathrm{SC}}^{\mathrm{HT}}_{(1,2)}\),
\(S\), and \(\mathcal C^{\mathrm{op}}_\partial\), and bracket
induced by coloured operadic insertion.  The filtration \(F^m\) is by
total arity \(r+s\ge m\).  A mixed-HT Swiss-cheese action datum on
\((S,\mathcal C^{\mathrm{op}}_\partial)\) is a Maurer--Cartan element
\[
  \Theta_{\mathrm{SC}}\in
  \mathfrak{Def}^{\mathrm{HT},1}_{\mathrm{SC}}
  (S,\mathcal C^{\mathrm{op}}_\partial),
  \qquad
  d\Theta_{\mathrm{SC}}+
  \frac12[\Theta_{\mathrm{SC}},\Theta_{\mathrm{SC}}]=0 .
\]
For an \(n\)-truncated lift \(\Theta_{\le n}\), define the arity
\((n+1)\) curvature representative
\[
  \mathcal R^{\mathrm{SC}}_{n+1}(\Theta_{\le n})
  =
  \left(d\Theta_{\le n}+\frac12[\Theta_{\le n},\Theta_{\le n}]\right)_
  {F^{n+1}/F^{n+2}}
  \in
  \left(
    F^{n+1}\mathfrak{Def}^{\mathrm{HT}}_{\mathrm{SC}}/
    F^{n+2}\mathfrak{Def}^{\mathrm{HT}}_{\mathrm{SC}}
  \right)^2 .
\]
It is closed for the twisted differential \(d_{\Theta_{\le n}}\).  The
obstruction class is
\[
  \mathfrak o^{\mathrm{SC}}_{n+1}(\Theta_{\le n})
  =
  [\mathcal R^{\mathrm{SC}}_{n+1}(\Theta_{\le n})]
  \in
  H^2\!\left(
    F^{n+1}\mathfrak{Def}^{\mathrm{HT}}_{\mathrm{SC}}/
    F^{n+2}\mathfrak{Def}^{\mathrm{HT}}_{\mathrm{SC}},
    d_{\Theta_{\le n}}\right).
\]
Choosing a primitive
\[
  H_{n+1}\in
  \left(F^{n+1}\mathfrak{Def}^{\mathrm{HT}}_{\mathrm{SC}}/
  F^{n+2}\mathfrak{Def}^{\mathrm{HT}}_{\mathrm{SC}}\right)^1,
  \qquad
  d_{\Theta_{\le n}}H_{n+1}
  =-\mathcal R^{\mathrm{SC}}_{n+1}(\Theta_{\le n}),
\]
is the arity-\((n+1)\) null-homotopy.
The specialization \(S=Z_\partial\) is the centre-action complex.  The
specialization
\(S=\mathcal A^{\mathrm{cl}/\hbar}_{\mathrm{bulk}}\) is the bulk-action
complex.  A closed-colour morphism
\(\alpha\colon\mathcal A^{\mathrm{cl}/\hbar}_{\mathrm{bulk}}\to
Z_\partial\) pulls the centre action back to a bulk action exactly when
\(\alpha\) satisfies the mixed operations and higher centrality equations
in the open--closed comparison complex of
Definition~\ref{defn:oca-comparison}.
\end{defn}

\begin{thm}[Mixed-HT chiral Deligne--Swiss-cheese action criterion]
\label{thm:sc-ht-deligne-swiss-cheese}

Assume that the arity filtration in
Definition~\ref{defn:sc-ht-action-datum} is separated and complete, and
fix \(S\in\{Z_\partial,\mathcal A^{\mathrm{cl}/\hbar}_{\mathrm{bulk}}\}\)
together with an arity-\(\le2\) boundary datum in
\(\mathfrak{Def}^{\mathrm{HT}}_{\mathrm{SC}}
(S,\mathcal C^{\mathrm{op}}_\partial)\): a closed binary operation on
\(S\) whose formal Hamiltonian-current representative is the diagonal
singular product, the \(E_1^{\mathrm{top}}\)-product on
\(A^{\mathrm{cl}}_{\partial,N}\), and the mixed trace-sector operation
of Theorem~\ref{thm:main-local}.  Extending this boundary datum to an
algebra over \(\mathrm{SC}^{\mathrm{HT}}_{(1,2)}\) on
\[
  (S,\mathcal C^{\mathrm{op}}_\partial)
\]
is equivalent to supplying a Maurer--Cartan lift
\(\Theta_{\mathrm{SC}}\) in
\(\mathfrak{Def}^{\mathrm{HT}}_{\mathrm{SC}}
(S,\mathcal C^{\mathrm{op}}_\partial)\).
Such a lift exists if and only if every obstruction class
\[
  \mathfrak o^{\mathrm{SC}}_{n+1}(\Theta_{\le n})
  \in
  H^2\!\left(
    F^{n+1}\mathfrak{Def}^{\mathrm{HT}}_{\mathrm{SC}}/
    F^{n+2}\mathfrak{Def}^{\mathrm{HT}}_{\mathrm{SC}},
    d_{\Theta_{\le n}}\right),
  \qquad n\ge2,
\]
vanishes and the chosen primitives \(H_{n+1}\) are compatible under the
completed inverse limit.  For \(S=Z_\partial\), the closed
Morita-centre component of this obstruction tower is
\(\operatorname{Ob}^{\mathrm{ch}\,E_2}_{\mathrm{Deligne}}\).  The global
bulk form is the specialization \(S=\mathcal A^{\mathrm{cl}/\hbar}_{\mathrm{bulk}}\);
it is governed by tangential \(E_1^{\mathrm{top}}\) coherence, the
open--closed comparison criterion of
Theorem~\ref{thm:main-global-deligne-criterion}, and the
comparison-cone contraction of
Definition~\ref{defn:oca-comparison-cone}.

The pair
\begin{equation}\label{eq:sc-ht-action-on-pair}
  \bigl(Z^{\mathrm{der},\mathrm{Mor}}_{E_2^{\mathrm{HT}}}\!\bigl(
	    \mathcal C^{\mathrm{op}}_\partial\bigr),\,
	    \mathcal C^{\mathrm{op}}_\partial\bigr)
\end{equation}
has, at the formal-Darboux stalk, the following centre-action
arity-\(\le2\) boundary components in \(\mathcal B^{\mathrm{adm}}\):
\begin{enumerate}[label=\textup{(\roman*)}]
  \item closed colour acting on
        \(Z^{\mathrm{der},\mathrm{Mor}}_{E_2^{\mathrm{HT}}}(
        \mathcal C^{\mathrm{op}}_\partial)\) by the chiral brace
        operations on
        \(C^\bullet_{\mathrm{Hoch},\mathrm{adm}}(A^{\mathrm{cl}}_{\partial,N},
        A^{\mathrm{cl}}_{\partial,N})\) of
        \eqref{eq:bar-resolution-derived-centre}, with binary singular
        product \eqref{eq:singular-product-linear};
  \item open colour acting on \(\mathcal C^{\mathrm{op}}_\partial\)
        through the \(E_1^{\mathrm{top}}\)-product of
        \(A^{\mathrm{cl}}_{\partial,N}\);
  \item trace-sector mixed binary operation
        \(\mu^{\mathrm{mix}}_{\mathsf c,\mathsf o}\bigl(c\otimes a\bigr)
        =c\cdot a\), representing the action of the trace-sector closed
        generators on the open sector; the higher mixed operations,
        when the Maurer--Cartan lift exists, are the equivariance
        homotopy hierarchy
        \(\{H^{\mathrm{prod}}_{f,-},H^{P_0}_{f,-},H^3_{f,-,-},\ldots\}\)
        of \eqref{eq:ce-pv-dictionary-strong}.
\end{enumerate}
The trace-sector action is bulk--boundary compatible: the Hochschild
component sends \(c_f,u_f\) to \(\theta_f,J(f)\cdot(-)\).  For
polynomial \(f\), \(J(f)\) is represented by
\(\overline{\operatorname{Tr}\,\tilde f(\phi_1,\phi_2)}\) for any
ordinary-word lift \(\tilde f\); for \(f\in\mathfrak h\) it is the
completed stable scalar-reduced inverse limit.
After continuous coordinate duality, the scalar-contact component sends
\(c_K\) to \(\gamma_{\mathbf 1}\), while the rank-
\(N\) trace character is \(\chi_N(K)=\operatorname{Tr}I_N=N\).
The raw Rees/Weyl commutator contributes the separate factor
\(\hbar_{\mathrm W}\).  A full
closed action of
\(Z^{\mathrm{der},\mathrm{Mor}}_{E_2^{\mathrm{HT}}}\) requires the
centrality hierarchy, equivalently the all-arity Maurer--Cartan lift
above.

If, in addition, an open--closed correspondence quasi-isomorphism
\[
  \mathrm{OCA}_N\colon\mathcal A^{\mathrm{cl}}_{\mathrm{bulk}}
  \xrightarrow{\sim}
  Z^{\mathrm{der},\mathrm{Mor}}_{E_2^{\mathrm{HT}}}
  (\mathcal C^{\mathrm{op}}_\partial)
\]
is supplied as in Theorem~\ref{thm:typed-global-mixed-ht-deligne}, with
the mixed-HT Swiss-cheese action Maurer--Cartan lift, Hochschild-centre
chart lift, representatives \(\mathcal R^{\mathrm{global},5}\)
satisfying \eqref{eq:global-five-primitive-equations}, flat
Darboux-torsor transport when the datum varies in formal coordinates,
the unreduced cotangent-current row when the full BF source is used, and
a contracting homotopy for \(\operatorname{Cone}(\mathrm{OCA}_N)\), then
\((\mathcal A^{\mathrm{cl}}_{\mathrm{bulk}},
\mathcal C^{\mathrm{op}}_\partial)\) carries the supplied physical
\(\mathrm{SC}^{\mathrm{HT}}_{(1,2)}\)-bulk--boundary system of the
mixed-HT theory on \(X\) at \(N\) Dirac branes.
\end{thm}

\begin{proof}
For a cofibrant two-coloured dg operad, algebra structures on a fixed
pair are Maurer--Cartan elements in the completed operadic convolution
dg Lie algebra from the reduced cooperad to the coloured endomorphism
operad of the pair.  Applied to
\(\mathrm{SC}^{\mathrm{HT}}_{(1,2)}\) this is
Definition~\ref{defn:sc-ht-action-datum}.  The arity filtration is
pronilpotent by hypothesis, so the Maurer--Cartan equation is solved
recursively.

Assume \(\Theta_{\le n}\) satisfies the Maurer--Cartan equation modulo
\(F^{n+1}\).  Its curvature in
\(F^{n+1}/F^{n+2}\) is closed for the twisted differential
\(d_{\Theta_{\le n}}\) by the Bianchi identity
\(d_\Theta(d\Theta+\frac12[\Theta,\Theta])=0\).  Changing the
arity-\((n+1)\) component by a degree-one term \(H_{n+1}\) changes this
curvature by \(d_{\Theta_{\le n}}H_{n+1}\).  Hence the extension exists
at arity \(n+1\) exactly when
\(\mathfrak o^{\mathrm{SC}}_{n+1}(\Theta_{\le n})\) vanishes, and the
primitive \(H_{n+1}\) is the required higher homotopy.  Completeness
then identifies a compatible system of arity lifts with a genuine
Maurer--Cartan element \(\Theta_{\mathrm{SC}}\).

The prescribed arity-\(\le2\) boundary components consist of the local
formal-Darboux data together with representation choices for the chosen
closed colour \(S\).  For \(S=Z_\partial\), the closed binary piece is
the Hochschild-chart representative of the chiral brace singular product
with diagonal local-cohomology kernel \(K^{(2)}_\Delta\) of
\eqref{eq:singular-product-linear}; it is a category-level Morita centre
operation only after the centre chart and its centrality homotopies have
been fixed.  For \(S=\mathcal A^{\mathrm{cl}/\hbar}_{\mathrm{bulk}}\),
the closed binary piece is the Hamiltonian BF singular product on closed
observables; its trace-sector image under the open--closed comparison is
the same diagonal kernel \(K^{(2)}_\Delta\).  The open binary piece is the
\(E_1^{\mathrm{top}}\)-product on
\(A^{\mathrm{cl}}_{\partial,N}\).  The trace-sector mixed binary piece
is the centrality arrow of the closed colour acting on
\(A^{\mathrm{cl}}_{\partial,N}\); its first Taylor coefficient is
\(\tau_{N,1}^{\mathrm{tr}}\).  Under continuous coordinate duality and
\eqref{eq:ce-pv-dictionary-strong} its expression is
\[
  c_f\mapsto\theta_f,\qquad
  u_f\mapsto J(f),\qquad
  c_K\mapsto\gamma_{\mathbf 1},
\]
and the rank-
\(N\) trace character evaluates
\(\chi_N(K)=\operatorname{Tr}I_N=N\).  The source cocycle remains
unscaled; its raw Rees/Weyl trace commutator is
\(\hbar_{\mathrm W}N[\bar c]\).  These formulas determine the
trace-sector arity-two boundary condition for the obstruction recursion;
they do not by themselves construct the all-arity Swiss-cheese action.

The global comparison is the obstruction criterion of
Theorem~\ref{thm:main-global-deligne-criterion} and its typed
refinement Theorem~\ref{thm:typed-global-mixed-ht-deligne}.  After the
Swiss-cheese curvature representatives
\(\mathcal R^{\mathrm{SC}}_{n+1}(\Theta_{\le n})\), the
Hochschild-centre lift, the components of
\(\mathcal R^{\mathrm{global},5}\), the Darboux-torsor transport row,
and the cotangent-current row admit the primitive equations stated in
Theorem~\ref{thm:typed-global-mixed-ht-deligne}, and after
\(\operatorname{Cone}(\mathrm{OCA}_N)\) is contracted,
this contraction is the quasi-isomorphism datum for \(\mathrm{OCA}_N\)
on \(\partial X\) in the supplied
\(\mathrm{SC}^{\mathrm{HT}}_{(1,2)}\)-category.
\end{proof}

\subsection{Scalar non-faithfulness}
\label{ssec:scalar-non-faithfulness}

The graded Euler character is the protected scalar trace
\(\chi=\operatorname{sTr}\colon\mathrm{FiltCh}_{
\mathrm{SC}^{\mathrm{HT}}_{(1,2)}}\to\Z((q,\hbar_{\mathrm{W}}))\) defined on
chain-level filtered \(\mathrm{SC}^{\mathrm{HT}}_{(1,2)}\)-algebras with
a compatible homological grading and a filtration by charge / weight /
height.  It is a scalar invariant, not a chain-level invariant; scalar
non-faithfulness is the elementary obstruction to recovering a chain-level
\(\mathrm{SC}^{\mathrm{HT}}_{(1,2)}\)-object from its scalar trace.

\begin{lem}[Scalar non-faithfulness]\label{lem:scalar-non-faithfulness}
The graded Euler supertrace
\[
  \chi=\operatorname{sTr}(-)\colon
  \mathrm{FiltCh}_{\mathrm{SC}^{\mathrm{HT}}_{(1,2)}}
  \longrightarrow\Z((q,\hbar_{\mathrm{W}}))
\]
has nontrivial fibres: equality of scalar traces is compatible with
inequivalent quasi-isomorphism types, with inequivalent
\(\mathrm{SC}^{\mathrm{HT}}_{(1,2)}\)-operations, and with distinct
Hall-style products, Drinfeld-double pairings, current-envelope
locality structures, derived-centre trace maps, and open--closed
comparison morphisms.  Consequently every scalar-to-chain lift of
automorphic data includes a chain-level
\(\mathrm{SC}^{\mathrm{HT}}_{(1,2)}\)-object.
\end{lem}

\begin{proof}
Let \(C=\C\) be concentrated in degree \(0\) with the trivial filtered
\(\mathrm{SC}^{\mathrm{HT}}_{(1,2)}\)-structure.  Let
\[
  V=\C e_0\oplus \C e_1,\qquad |e_0|=0,\quad |e_1|=1,\quad d_V=0,
\]
and set \(C'=\C\oplus V\) with \(V^2=0\), unit \(1\in\C\), and all
mixed operations on \(V\) square-zero.  Then
\[
  \chi(C')=1+\chi(V)=1+1-1=\chi(C),
\]
while \(H^\bullet(C')=\C\oplus \C e_0\oplus \C e_1\) is not
isomorphic to \(H^\bullet(C)=\C\).  Thus the scalar trace takes the
same value on two distinct quasi-isomorphism types.  Keeping the same graded
complex and changing square-zero products, Hall pairings, current
envelopes, or open--closed operations on \(V\) gives further points of
the same scalar fibre with different chain-level structures.
\end{proof}

\begin{cor}[Mandatory chain-level supplements]
\label{cor:mandatory-chain-supplements}
Every lift from automorphic scalar to physical operator --- in
particular of the Igusa
		cusp form \(\Phi_{10}\) or its Borcherds square root
		\(\Delta_5=\Phi_{10}^{1/2}\) --- consists of a chain-level
filtered \(\mathrm{SC}^{\mathrm{HT}}_{(1,2)}\)-object, for example the
filtered boundary operator algebra
\(\mathrm{Op}^{\mathrm{HT}}_{\partial X}\) of
Definition~\ref{defn:six-boundary-open-closed-objects}~(C6), together with an
identification of its protected trace with the scalar invariant.  The
scalar identity is the trace value; the chain-level operator is the
filtered object whose trace it is.
\end{cor}

\begin{proof}
Direct from Lemma~\ref{lem:scalar-non-faithfulness} applied to the
chain-level objects whose protected traces are the
scalar invariants in question.
\end{proof}

\subsection{Topologisation criterion}
\label{ssec:topologisation}

The normal mixed-HT operad \(\mathrm{SC}^{\mathrm{HT}}_{(1,2)}\) of
Definition~\ref{defn:sc-ht-operad} is a two-coloured operad by
construction.  Topologisation is the problem of replacing the
holomorphic closed colour by a topological closed colour through a
homotopy-coherent trivialisation of holomorphic translations.  When
operators \(G_i\) are supplied, the equations \(T_i=[Q,G_i]\),
\(i=1,2\), are the unary part of that problem.

\begin{defn}[Topologisation datum]\label{defn:topologisation-datum}
A \emph{topologisation datum} for the mixed-HT closed sector
\((\mathcal A^{\mathrm{cl}}_{\mathrm{bulk}},
\mathrm{SC}^{\mathrm{HT}}_{(1,2)})\) is a tuple
\((\mathcal C_\rho,Q,G_1,G_2,T_1,T_2)\) where
\(\mathcal C_\rho\) is a complete filtered ambient
symmetric-monoidal \(\infty\)-category refining
\(\mathcal B^{\mathrm{adm}}\) by a weight filtration \(\rho\); \(Q\) is
the BRST/BV differential of
Statement~\ref{stmt:costello-bv-construction}; \(G_i\) are
\(\rho\)-degree-\(-1\) operators on
\(\mathcal A^{\mathrm{cl}}_{\mathrm{bulk}}\); \(T_i\) are the two
holomorphic translation stress-tensor components; and they satisfy
\begin{equation}\label{eq:stress-tensor-null-homotopy}
  T_i=[Q,G_i]\quad(i=1,2)\quad\text{in }\mathcal C_\rho.
\end{equation}
\end{defn}

\begin{defn}[Topological transfer complex]
\label{defn:topological-transfer-complex}
Let
\(\mathfrak t_{\mathrm{hol}}=\C\partial_{z_1}\oplus
\C\partial_{z_2}\).  For a mixed-HT action
\(\Theta_{\mathrm{HT}}\) on \(\mathcal A^{\mathrm{cl}}_{\mathrm{bulk}}\),
write
\[
  \mathrm{Der}^{\rho}_{\mathrm{SC}^{\mathrm{HT}}}
  (\mathcal A^{\mathrm{cl}}_{\mathrm{bulk}})
\]
for the complete filtered dg Lie algebra of
\(\mathrm{SC}^{\mathrm{HT}}_{(1,2)}\)-derivations in
\(\mathcal C_\rho\).  The homotopy-coherent trivialisation of
holomorphic translations is governed by
\[
  C^\bullet_{\mathrm{CE}}\!\left(
    \mathfrak t_{\mathrm{hol}},
    \mathrm{Der}^{\rho}_{\mathrm{SC}^{\mathrm{HT}}}
    (\mathcal A^{\mathrm{cl}}_{\mathrm{bulk}})\right),
\]
with total degree equal to Chevalley--Eilenberg degree plus internal
degree.  The obstruction to factoring the closed colour through the
topological operad is measured by the mapping-cone deformation complex
\[
  \mathfrak{Def}^{\rho}_{\mathrm{SC},\mathrm{top}}
  (\mathcal A^{\mathrm{cl}}_{\mathrm{bulk}})
  =
  \operatorname{Cone}\!\left(
    \mathfrak{Def}^{\rho}_{\mathrm{SC}^{\mathrm{top}}}
    (\mathcal A^{\mathrm{cl}}_{\mathrm{bulk}})
    \longrightarrow
    \mathfrak{Def}^{\rho}_{\mathrm{SC}^{\mathrm{HT}}}
    (\mathcal A^{\mathrm{cl}}_{\mathrm{bulk}})
  \right)[-1],
\]
where the arrow is induced by forgetting topological constancy in the
two holomorphic directions.  The complete topological transfer complex is
\[
  \mathfrak T_\rho(\mathcal A^{\mathrm{cl}}_{\mathrm{bulk}})
  =
  C^\bullet_{\mathrm{CE}}\!\left(
    \mathfrak t_{\mathrm{hol}},
    \mathrm{Der}^{\rho}_{\mathrm{SC}^{\mathrm{HT}}}
    (\mathcal A^{\mathrm{cl}}_{\mathrm{bulk}})\right)
  \widehat\oplus
  \mathfrak{Def}^{\rho}_{\mathrm{SC},\mathrm{top}}
  (\mathcal A^{\mathrm{cl}}_{\mathrm{bulk}}).
\]
The translation cocycle \(\tau\in\mathfrak T_\rho^1\) is defined by
\(\tau(\partial_{z_i})=T_i\) and has zero operadic-cone component.  A
homotopy-coherent translation trivialisation is a primitive
\(\Gamma\in
C^\bullet_{\mathrm{CE}}(\mathfrak t_{\mathrm{hol}},
\mathrm{Der}^{\rho}_{\mathrm{SC}^{\mathrm{HT}}}
(\mathcal A^{\mathrm{cl}}_{\mathrm{bulk}}))^0\)
with
\[
  D_\rho\Gamma=\tau,
\]
so \(\Gamma(\partial_{z_i})=G_i\) and
\(\Gamma(\partial_{z_1}\wedge\partial_{z_2})\) is the secondary
commutation homotopy.  Equivalently, writing
\[
  G_{12}:=\Gamma(\partial_{z_1}\wedge\partial_{z_2})
\]
and letting \(\mathfrak t_{\mathrm{hol}}\) act on derivations by
commutator with the translation derivations \(T_i\), the equation
\(D_\rho\Gamma=\tau\) contains
\[
  [Q,G_i]=T_i,\qquad
  [Q,G_{12}]+[T_1,G_2]-[T_2,G_1]=0 .
\]
The second identity is part of the datum; it is not implied by the two
unary equations.  The first obstruction is
\[
  [\tau]\in
  H^1\!\left(
    C^\bullet_{\mathrm{CE}}\!\left(
      \mathfrak t_{\mathrm{hol}},
      \mathrm{Der}^{\rho}_{\mathrm{SC}^{\mathrm{HT}}}
      (\mathcal A^{\mathrm{cl}}_{\mathrm{bulk}})\right)\right).
\]
After such a translation primitive is fixed, the arity-filtered
Maurer--Cartan lift in
\(\mathfrak{Def}^{\rho}_{\mathrm{SC},\mathrm{top}}\) has obstruction
classes
\[
  \mathfrak o^{\mathrm{top}}_{n+1}(\Theta^{\mathrm{top}}_{\le n})
  \in
  H^2\!\left(
    F^{n+1}\mathfrak{Def}^{\rho}_{\mathrm{SC},\mathrm{top}}/
    F^{n+2}\mathfrak{Def}^{\rho}_{\mathrm{SC},\mathrm{top}},\,
    d_{\Theta^{\mathrm{top}}_{\le n}}\right).
\]
For a lift chosen through arity \(n\), a representative is
\[
  \mathcal R^{\mathrm{top}}_{n+1}
  =
  \left(
  d\Theta^{\mathrm{top}}_{\le n}
  +\frac12[\Theta^{\mathrm{top}}_{\le n},
          \Theta^{\mathrm{top}}_{\le n}]
  +\sum_{k\ge3}\frac1{k!}
    \ell_k((\Theta^{\mathrm{top}}_{\le n})^{\otimes k})
  \right)_{F^{n+1}/F^{n+2}},
\]
and the next operation is supplied by a primitive
\[
  d_{\Theta^{\mathrm{top}}_{\le n}}H^{\mathrm{top}}_{n+1}
  =-\mathcal R^{\mathrm{top}}_{n+1}.
\]
\end{defn}

\begin{thm}[Topologisation obstruction criterion]
\label{thm:topologisation-by-stress-tensor}
Inside the complete filtered refinement $\mathcal C_\rho$ of
$\mathcal B^{\mathrm{adm}}$, the equations
\eqref{eq:stress-tensor-null-homotopy} are the first row of the
topologisation problem.  They do not by themselves construct an
\(E_3^{\mathrm{top}}\)-algebra structure.  A topological closed colour
\(\mathrm{SC}^{\mathrm{top}}_{(1,2)}\), hence an
\(E_3^{\mathrm{top}}\)-structure on the closed sector by Dunn
additivity, is obtained exactly when the topological transfer complex of
Definition~\ref{defn:topological-transfer-complex} has vanishing
obstruction tower:
\[
  [\tau]=0,\qquad
  \mathfrak o^{\mathrm{top}}_{n+1}(\Theta^{\mathrm{top}}_{\le n})=0
  \quad(n\ge 1),
\]
with primitives compatible under the transition maps of the completed
inverse system.

Suppose a topologisation datum \((\mathcal C_\rho,Q,G_1,G_2,T_1,T_2)\)
of Definition~\ref{defn:topologisation-datum} is given.  It supplies
only the unary translation primitives \(G_i\).  The completed closed
sector is topological in the two holomorphic directions, and the
transferred closed colour has an \(E_3^{\mathrm{top}}\)-algebra
structure, if and only if the following chain-level data are supplied:
\begin{enumerate}[label=\textup{(T\arabic*)}]
  \item the operators \(G_1,G_2\) are constructed and satisfy
        \([Q,G_i]=T_i\) as identities of chain-level operators in
        \(\mathcal C_\rho\);
  \item the identities \(T_i=[Q,G_i]\) are compatible with the weight
        filtration \(\rho\) and with the completion in
        \(\mathcal C_\rho\);
  \item there is a continuous secondary operator
        \(G_{12}\in
        \mathrm{Der}^{\rho}_{\mathrm{SC}^{\mathrm{HT}}}
        (\mathcal A^{\mathrm{cl}}_{\mathrm{bulk}})^{-2}\)
        satisfying
        \[
          [Q,G_{12}]+[T_1,G_2]-[T_2,G_1]=0,
        \]
        and its transition defects under arity truncation and admissible
        weight transport are killed by primitives in the same
        Chevalley--Eilenberg/Roos total complex;
  \item the topologisation has finite holomorphic propagation in
        \(\mathcal C_\rho\): for some constants
        \(c_1,c_2,c_{12}\) and every weight level \(p\),
        \[
          G_i(F^p_\rho)\subset F^{p-c_i}_\rho,\qquad
          G_{12}(F^p_\rho)\subset F^{p-c_{12}}_\rho,
        \]
        and, in the analytic local-functional model, these inclusions are
        continuous estimates.  Namely, for \(a\in\{1,2,12\}\), \(G_a\)
        is defined on a closed microlocal subcomplex
        \(\mathcal C^{\mathrm{top,WF}}_{\rho,w,M,L}(I)\subset
        \mathcal Q^\bullet_{w,r,M,L}(I)\) generated by regular density
        labels and by the graphwise wavefront-admissible current labels
        \(\mathcal D'_{c,\mathrm{WF}}(I;\Gamma,M)\) of
        Definition~\ref{def:app-wavefront-admissible-defect-labels}, not
        on all of \(\mathcal D'_c(I)\).  For every defining seminorm of
        this locally convex finite-scale model there are constants
        \(C_{a,p,w,M,L}\) with
        \[
          \|G_a x\|_{\rho,p-c_a;w,M,L}
          \leq C_{a,p,w,M,L}\|x\|_{\rho,p;w,M,L}.
        \]
        Writing \(K^{M,G_a}_{\Gamma,\alpha}\) for the finite-window
        Hadamard coefficient of a Costello graph with one \(G_a\)-insertion
        and \(B_\Gamma=e_\Gamma^*(B_1\boxtimes\cdots\boxtimes B_r)\), the
        estimates required at every graph, coefficient label, and
        collision diagonal are
        \[
          N_{e_\Gamma}\cap
          \mathrm{WF}(B_1\boxtimes\cdots\boxtimes B_r)=\emptyset,
        \]
        \[
          \bigl(\mathrm{WF}(K^{M,G_a}_{\Gamma,\alpha})
          +\mathrm{WF}(B_\Gamma)\bigr)
          \cap T^*_{X_\Gamma(I)}X_\Gamma(I)=\emptyset,
        \]
        \[
          \operatorname{sd}_{\Delta}
          (K^{M,G_a}_{\Gamma,\alpha}B_\Gamma)<\infty,\qquad
          \mathrm{WF}\bigl(p_{\Gamma,*}
          (K^{M,G_a}_{\Gamma,\alpha}B_\Gamma)\bigr)
          \subset
          p_{\Gamma,\#}\mathrm{WF}
          (K^{M,G_a}_{\Gamma,\alpha}B_\Gamma).
        \]
        The local extensions across \(\Delta\) and their counterterm
        ambiguities must be compatible with interval extension by zero,
        finite-window truncation, and admissible weight transport as in
        Proposition~\ref{prop:app-wavefront-admissible-costello-graph-kernel};
        if any displayed condition fails, the topologisation operation is
        not an element of the declared local-functional complex and an
        additional renormalization rule for the corresponding
        distributional product is required;
  \item the transferred \(E_3^{\mathrm{top}}\)-operations are
        independent of the small-disks SDR transfer choices used to
        carry \eqref{eq:stress-tensor-null-homotopy} across the bar
        resolution: for two SDR choices \((p,j,h)\) and
        \((p',j',h')\), the defect
        \[
          \mathcal D^{\mathrm{SDR}}_{n+1}
          =
          \Theta^{\mathrm{top}}_{n+1}(p,j,h)
          -\Theta^{\mathrm{top}}_{n+1}(p',j',h')
        \]
        is killed by a continuous primitive
        \(K^{\mathrm{SDR}}_{n+1}\) with
        \[
          d_{\Theta^{\mathrm{top}}_{\le n}}K^{\mathrm{SDR}}_{n+1}
          =-\mathcal D^{\mathrm{SDR}}_{n+1};
        \]
  \item the anomaly class \(\operatorname{Ob}^{\mathrm{anom}}\) has a
        representative \(\mathcal R^{\mathrm{anom}}\) whose reduced part
        satisfies
        \(d_{\mathrm{Conv}}h_{\mathrm{anom}}
        =-(\mathcal R^{\mathrm{anom}})_{\mathrm{red}}\) in the
        operadic convolution complex of
        Definition~\ref{defn:convolution-complex},
        \(\operatorname{Ob}^{\mathrm{anom}}\in
        H^2\bigl(\mathrm{Conv}_{\mathrm{SC}}\!\bigl(
          \mathcal A^{\mathrm{cl}}_{\mathrm{bulk}},
          Z^{\mathrm{der},\mathrm{Mor}}_{E_2^{\mathrm{HT}}}\!\bigl(
            \mathcal C^{\mathrm{op}}_\partial\bigr)\bigr)\bigr)=0\),
        with secondary coherence class in
        \(H^3\bigl(\mathrm{Conv}_{\mathrm{SC}}(\cdots)\bigr)\).
\end{enumerate}
For every \(n\ge1\) the topological arity curvature
\(\mathcal R^{\mathrm{top}}_{n+1}\) must also satisfy
\[
  d_{\Theta^{\mathrm{top}}_{\le n}}H^{\mathrm{top}}_{n+1}
  =-\mathcal R^{\mathrm{top}}_{n+1}
\]
with the same arity, weight, SDR, and boundary-restriction coherence.
Under these hypotheses the unscaled extension class is \([\bar c]\),
and its raw Rees/Weyl trace pushout
\(\hbar_{\mathrm{W}}\chi_N(K)[\bar c]
=\hbar_{\mathrm{W}}N[\bar c]\) is carried only as the trace-sector
projective Lie anomaly.  Its interpretation as determinant-line or
modular-line curvature is the separate supplied line, connection,
Atiyah-class pullback, and lift killing \(\operatorname{Ob}^{\det}_J\)
of Definition~\ref{defn:capelli-determinant-line-obstruction}.  Its
interpretation as a BCOV stress-tensor anomaly is a separate matched
BCOV stress-tensor comparison datum.
\end{thm}

\begin{proof}
The two operators \(T_i\) define a Chevalley--Eilenberg cocycle
\(\tau\) for the abelian holomorphic translation Lie algebra
\(\mathfrak t_{\mathrm{hol}}\) with coefficients in the filtered
derivations of the mixed-HT action.  Equations
\([Q,G_i]=T_i\) say that the unary part of this cocycle is exact.  They
therefore give local constancy of the bulk currents only after passing
to cohomology and only in the derivation complex where the \(G_i\) are
continuous.  The two unary equations do not imply homotopy-coherent
triviality of the \(\mathfrak t_{\mathrm{hol}}\)-action: the two-simplex
face is represented by \([T_1,G_2]-[T_2,G_1]\), and \(G_{12}\) is the
chain homotopy killing this face.

Chain-level topological constancy is the stronger statement that the
translation action is homotopy-coherently trivial and that the
\(\mathrm{SC}^{\mathrm{HT}}_{(1,2)}\)-action factors through the
topological closed colour.  This is precisely the lifting problem in
\(\mathfrak T_\rho(\mathcal A^{\mathrm{cl}}_{\mathrm{bulk}})\).  The
first obstruction is \([\tau]\).  After it is killed by \(\Gamma\), the
remaining problem is the arity-filtered Maurer--Cartan lifting problem in
\(\mathfrak{Def}^{\rho}_{\mathrm{SC},\mathrm{top}}\).  If
\(\Theta^{\mathrm{top}}_{\le n}\) solves the equation modulo
\(F^{n+1}\), its curvature in
\(F^{n+1}/F^{n+2}\) is closed for
\(d_{\Theta^{\mathrm{top}}_{\le n}}\).  Changing the next component by a
degree-one primitive changes the curvature by an exact term.  Hence the
lift exists through arity \(n+1\) exactly when
\(\mathfrak o^{\mathrm{top}}_{n+1}(\Theta^{\mathrm{top}}_{\le n})\)
vanishes.  Completeness of \(\mathcal C_\rho\) identifies a compatible
system of finite-arity lifts with a full transfer datum.  The finite
  propagation, wave-front, scaling-degree, and SDR-choice clauses state
  that the constructed primitives remain in the same completed
  local-functional/Roos complex after passing between weights, heat-kernel
  scales, and small-disks contractions.  H\"ormander's product criterion is
  load-bearing here: the graphwise conditions of
  Definition~\ref{def:app-wavefront-admissible-defect-labels} are exactly
  the conditions under which the \(G_a\)-inserted Hadamard coefficient can
  be multiplied by the brane current, extended across collision diagonals,
  and pushed forward to the brane-defect local-functional target.  Without
  these clauses the construction is only a formal unary stress-tensor
  identity.

With this transfer datum fixed, the \(E_2^{\mathrm{ch,hol}}\)-closed colour
has been replaced by the \(E_2^{\mathrm{top}}\)-closed colour.  Dunn
additivity then identifies
\(E_1^{\mathrm{top}}\otimes E_2^{\mathrm{top}}\) with
\(E_3^{\mathrm{top}}\)
\cite{lurie-ha}*{Theorem~5.1.2.2}.  The operadic transfer changes neither
the unscaled constant-Hamiltonian extension class \([\bar c]\) nor its
raw Rees/Weyl trace pushout
\(\hbar_{\mathrm{W}}N[\bar c]\).  A determinant-line or modular-line
curvature interpretation requires the separate supplied line,
connection, and Atiyah-class pullback.
\end{proof}

\begin{rmk}[Class \(\mathrm L\) versus class \(\mathrm M\) topologisation]
\label{rmk:class-L-vs-class-M}
For class \(\mathrm L\) (affine Kac--Moody at non-critical level)
boundary algebras the two translation null-homotopies
\(T_i=[Q,G_i]\) may be raw chain-level data, provided finite
holomorphic propagation is proved.  For
class \(\mathrm M\) (Virasoro / principal \(W\)-families) the admissible
topologisation is the weight-completed / pro / Banach ambient unless
finite holomorphic propagation is separately verified.  The
finite-jet PVA models exhibit this distinction: the classical BV action
is constructed, while the all-order
QME, the curve-restricted boundary vertex algebra produced by
\(\mathsf R_{L,\mathcal B_\perp}\), and the \(E_3^{\mathrm{top}}\)-lift
remain governed by the filtered formal SDR, sectorial analytic
regularity when analytic convergence is asserted, finite-window weight
symmetry, and the local-functional lifts \(T_i=d_{\mathrm{QME}}^{L}G_i\).
\end{rmk}

\subsection{The typed global mixed-HT Deligne criterion}
\label{ssec:typed-global-mixed-ht-deligne}

Theorem~\ref{thm:main-global-deligne-criterion} states the OCA
comparison on the Costello--Gwilliam brane-link boundary with the
mixed-HT Swiss-cheese action tower, the five-component obstruction
\(\operatorname{Ob}^{\mathrm{global},5}\), the comparison cone, and the
source complex in which each null-homotopy lives.  The
\(\mathrm{SC}^{\mathrm{HT}}_{(1,2)}\)-morphism is defined only after the
all-arity Swiss-cheese action datum of
Definition~\ref{defn:sc-ht-action-datum}; topologisation adds the
transfer complex of Definition~\ref{defn:topological-transfer-complex}.
Throughout this subsection
\[
  \mathcal A^{\mathrm{cl}/\hbar}_{\mathrm{bulk}}
\]
is a tagged notation for the chosen bulk source in
\(\mathcal B^{\mathrm{adm}}\).  The classical branch is the Hamiltonian
BF factorization algebra \(\mathcal A^{\mathrm{cl}}_{\mathrm{bulk}}\) of
Definition~\ref{def:closed-mixed-HT-factorization}.  The quantum branch
means the same source together with the finite-scale Costello QME tower,
propagators, BV Laplacians, counterterms, and scalar-contact projection
specified in Definition~\ref{defn:typed-brane-defect-qme-tower}.  The
slash in \(\mathrm{cl}/\hbar\) is not an equality between
\(\hbar_{\mathrm{QME}}\), \(\hbar_{\mathrm{W}}\), \(\hbar_\omega\), or
\(\hbar_{\mathrm{Rees}}\); these parameters remain typed separately.

\begin{defn}[Open--closed centrality convolution complexes]\label{defn:convolution-complex}
Let
\[
  Z_\partial =
  Z^{\mathrm{der},\mathrm{Mor}}_{E_2^{\mathrm{HT}}}\!\bigl(
    \mathcal C^{\mathrm{op}}_\partial\bigr).
\]
For a source object
\[
  S\in\Bigl\{
  \bar B^{E_1}(A^{\mathrm{cl}}_{\partial,N}),\,
  \mathcal A^{\mathrm{cl}}_{\mathrm{bulk}}\Bigr\}
\]
in \(\mathcal B^{\mathrm{adm}}\), its centrality convolution complex
with the derived centre has underlying Hom complex
\[
  \mathrm{Conv}(S,Z_\partial)
  =
  \underline{\mathrm{Hom}}_{\mathcal B^{\mathrm{adm}}}^\bullet(S,Z_\partial),
  \qquad
  d_{\mathrm{Conv}}\varphi
  =
  d_{Z_\partial}\varphi-(-1)^{|\varphi|}\varphi d_S .
\]
For \(S=\bar B^{E_1}(A^{\mathrm{cl}}_{\partial,N})\) we write
\begin{equation}\label{eq:convolution-complex-def}
	  \mathrm{Conv}_{\bar B}\!\bigl(
	    \bar B^{E_1}(A^{\mathrm{cl}}_{\partial,N}),
    Z^{\mathrm{der},\mathrm{Mor}}_{E_2^{\mathrm{HT}}}\!\bigl(
      \mathcal C^{\mathrm{op}}_\partial\bigr)\bigr)
   \;=\;
   \underline{\mathrm{Hom}}_{\mathcal B^{\mathrm{adm}}}^\bullet\!\bigl(
     \bar B^{E_1}\!\bigl(A^{\mathrm{cl}}_{\partial,N}\bigr),\,
     Z^{\mathrm{der},\mathrm{Mor}}_{E_2^{\mathrm{HT}}}\!\bigl(
       \mathcal C^{\mathrm{op}}_\partial\bigr)\bigr),
\end{equation}
The obstruction theory for \(\mathrm{SC}^{\mathrm{HT}}_{(1,2)}\)-algebra
morphisms is the corresponding operadic convolution dg Lie, or
\(L_\infty\), complex
\[
  \mathrm{Conv}_{\mathrm{SC}}(S,Z_\partial),
\]
whose differential is the linearisation of the
\(\mathrm{SC}^{\mathrm{HT}}_{(1,2)}\)-Maurer--Cartan equation at the
chosen source and target algebra structures, and whose bracket is induced
by the cooperadic decomposition.  The displayed Hom complex is its
underlying graded mapping complex; obstruction classes for centrality,
anomaly, and locality are taken in
\(\mathrm{Conv}_{\mathrm{SC}}\) after the operadic deformation problem has
been fixed.
\end{defn}

\begin{defn}[Anomaly and locality datum]
\label{defn:typed-anomaly-locality-datum}
Let
\[
  \mathrm{OCA}_N\colon
  \mathcal A^{\mathrm{cl}/\hbar}_{\mathrm{bulk}}
  \longrightarrow
  Z_\partial
\]
be a supplied mixed-HT open--closed comparison morphism in the operadic
convolution complex of Definition~\ref{defn:convolution-complex}.  Its
anomaly curvature is the degree-two operadic curvature
\[
  \mathcal R^{\mathrm{anom}}
  =
  d_{\mathrm{Conv}}\mathrm{OCA}_N
  +\frac12[\mathrm{OCA}_N,\mathrm{OCA}_N]_{\mathrm{SC}}
  +\frac{1}{3!}\ell_3(\mathrm{OCA}_N^{\otimes3})
  +\cdots
  \in
  \mathrm{Conv}_{\mathrm{SC}}^2(
  \mathcal A^{\mathrm{cl}/\hbar}_{\mathrm{bulk}},Z_\partial).
\]
The strict anomaly datum is a degree-one primitive
\[
  h_{\mathrm{anom}}\in
  \mathrm{Conv}_{\mathrm{SC}}^1(
  \mathcal A^{\mathrm{cl}/\hbar}_{\mathrm{bulk}},Z_\partial),
  \qquad
  d_{\mathrm{Conv}}h_{\mathrm{anom}}
  =-\mathcal R^{\mathrm{anom}} .
\]

The projective anomaly datum replaces strict vanishing by the target
central extension obtained from the canonical unscaled
constant-Hamiltonian extension by first taking the raw Rees/Weyl
commutator and then pushing out along \(\chi_N(K)=N\):
\[
  0\to \C[[\hbar_{\mathrm{W}}]]\,\gamma_{\mathbf 1}
  \to Z_\partial^{\mathrm{proj}}\to Z_\partial\to0
\]
whose scalar projection satisfies
\[
  \sigma_{\mathrm{Cap}}(\mathcal R^{\mathrm{anom}})
  =
  \hbar_{\mathrm{W}}\chi_N(K)\,\bar c\,\gamma_{\mathbf 1}
  =\hbar_{\mathrm{W}}N\,\bar c\,\gamma_{\mathbf 1}.
\]
The reduced anomaly curvature is
\[
  (\mathcal R^{\mathrm{anom}})_{\mathrm{red}}
  =
  \mathcal R^{\mathrm{anom}}
  -\hbar_{\mathrm{W}}N\,\bar c\,\gamma_{\mathbf 1},
\]
and the projective primitive satisfies
\[
  d_{\mathrm{Conv}}h_{\mathrm{anom}}
  =-(\mathcal R^{\mathrm{anom}})_{\mathrm{red}} .
\]
The source central extension in \(\widehat{\mathfrak g}_{\mathrm{cent}}\)
still has cocycle \(\bar c K\), with no factor of \(N\) or
\(\hbar_{\mathrm W}\).  No determinant-line or modular-line curvature
is asserted by this datum; a supplied line, connection, Atiyah-class
pullback to Lie cochains, and lift killing
\(\operatorname{Ob}^{\det}_J\) are separate data.

For locality, let \(\operatorname{Ran}_{\mathrm{disj}}(\partial X)\)
denote finite pairwise-disjoint configurations of opens in the
brane-link boundary.  For
\(\underline U=(U_1,\ldots,U_m)\) put
\[
  \mathcal E^\bullet_{\mathrm{loc}}(\underline U)
  =
  \operatorname{Cone}\!\left(
  \mathcal E^\bullet_{\mathrm{desc}}(\textstyle\coprod_iU_i)
  \longrightarrow
  \widehat\bigotimes_i
  \mathcal E^\bullet_{\mathrm{desc}}(U_i)
  \right)[-1],
\]
where the arrow is induced by factorization structure maps for the bulk
and centre.  The locality curvature is
\[
  \mathcal R^{\mathrm{loc}}
  =
  \partial_{\mathrm{loc}}\mathrm{OCA}_N
  \in
  \operatorname{Tot}^1\check C^\bullet_{\mathrm{Weiss,Ran}}
  (\mathfrak U;\mathcal E^\bullet_{\mathrm{loc}}).
\]
A locality datum is a degree-zero primitive
\[
  h_{\mathrm{loc}}\in
  \operatorname{Tot}^0\check C^\bullet_{\mathrm{Weiss,Ran}}
  (\mathfrak U;\mathcal E^\bullet_{\mathrm{loc}}),
  \qquad
  \partial_{\mathrm{loc}}h_{\mathrm{loc}}
  =-\mathcal R^{\mathrm{loc}}.
\]
The obstruction classes are
\[
  \operatorname{Ob}^{\mathrm{anom}}
  =
  [\mathcal R^{\mathrm{anom}}],
  \qquad
  \operatorname{Ob}^{\mathrm{loc}}
  =
  [\mathcal R^{\mathrm{loc}}],
\]
with the anomaly class interpreted either strictly or projectively as
above.  The trace-sector formal-Darboux theorem supplies the raw
Rees/Weyl projective scalar value
\(\hbar_{\mathrm{W}}\chi_N(K)[\bar c]
=\hbar_{\mathrm{W}}N[\bar c]\).  The anomaly and locality
rows require the separate primitives \(h_{\mathrm{anom}}\) and
\(h_{\mathrm{loc}}\) killing the two displayed curvature classes.
\end{defn}

\begin{defn}[Cotangent-current lift datum]
\label{defn:typed-cotangent-current-lift-datum}
For an interval \(I\subset\R_t\) put
\[
  \mathfrak g_I^{\mathrm{cur}}
  =
  \bigl(\Omega^\bullet_c(I)\widehat\otimes\mathfrak h\bigr)
  \ltimes
  \bigl(\mathcal C^{\mathrm{dR}}_{c,\mathrm{cur}}(I)
  \widehat\otimes P\bigr)[1],
\]
where \(\mathcal C^{\mathrm{dR}}_{c,\mathrm{cur}}(I)\) is the compactly
supported de Rham current cosheaf,
\(P=\mathfrak h^\vee_{\mathrm{cont}}\), and
\(\bar A=\C[z_1,z_2]/\C\cdot1\) is the polynomial Hamiltonian quotient
used in the interval current model.
Its reduced support-local boundary target is
\[
  A^{\mathrm{pp}}_\partial(I)
  =
  \widehat{\mathbf S}\bigl(\Omega^0_c(I)\widehat\otimes\bar A\bigr)
  \widehat\otimes
  \widehat{\mathbf S}\bigl(\mathcal D'_c(I)\widehat\otimes P[1]\bigr),
\]
the principal-part current algebra of
Theorem~\ref{thm:reduced-principal-part-boundary-current}.  The
unreduced target is
\[
  \widetilde{\mathrm{Obs}}^{\mathrm{cl,pp}}_{\partial,N}(I)
  =
  \mathrm{Obs}^{\mathrm{cl}}_{\mathrm{open},N}(I)
  \widehat\otimes
  \widehat{\mathbf S}\!\left(
    \mathcal C^{\mathrm{dR}}_{c,\mathrm{cur}}(I)
    \widehat\otimes P[1]\right).
\]
A cotangent-current lift datum is a continuous filtered
Maurer--Cartan morphism
\[
  \kappa^{\mathrm{cot}}_I\colon
  C^\bullet_{\mathrm{CE},\mathrm{cont}}(\mathfrak g_I^{\mathrm{cur}})
  \longrightarrow
  Z^{P_0}_{\mathrm{fact}}\!
  \bigl(\widetilde{\mathrm{Obs}}^{\mathrm{cl,pp}}_{\partial,N}(I)\bigr)
\]
whose reduced image is
\[
  u_{a(t)dt\otimes f}\longmapsto B_f(a),\qquad
  c_{B,\rho}\longmapsto\Theta_\rho(B),
\]
and whose reduced \(P_0\)-brackets are
\[
  \{B_f(a),\Theta_\rho(B)\}=\Theta_{f\cdot\rho}(aB),\qquad
  \{\Theta_\rho(B),\Theta_\sigma(C)\}=0 .
\]
Here \(a,b\in C^\infty_c(I)\), \(B,C\in\mathcal D'_c(I)\), and
\[
  (aB)(\varphi)=B(a\varphi),\qquad \varphi\in C^\infty_c(I).
\]
Thus the only distributional multiplication in the reduced
cotangent-current datum is smooth-times-current.  The expression
\(\{\Theta_\rho(B),\Theta_\sigma(C)\}=0\) declares the cotangent current
summand to be an abelian odd ideal; it does not form a product \(BC\).
The equal-time contact \(\delta(t-t')\) is the distribution on
\(I\times I\) supported on the diagonal \(\Delta_I=\{t=t'\}\), and it is
used in this reduced row only after pairing with smooth smearings
\(a(t)b(t')\), which gives the smooth product \(ab\).  Point-mass labels
\(\delta_{t_0}\) may be used as compactly supported currents, but a
Costello graph containing the singular boundary propagator against
\(\delta_{t_0}\boxtimes\delta_{t_0}\) belongs to the
wavefront-admissible/renormalized graph problem of
Proposition~\ref{prop:app-current-compatible-Dprimec-graph-criterion},
not to the reduced support-local current theorem.  Before reduction this
datum includes a primitive
projection
\[
  \pi_{\mathrm{prim}}\colon
  \mathrm{Obs}^{\mathrm{cl}}_{\mathrm{open},N}(I)
  \longrightarrow
  \Omega^\bullet_c(I)\widehat\otimes\bar A
\]
	compatible with \(Q_{\mathrm{open}}\), constants, decomposable traces,
	and factorization products, together with centrality homotopies against
	arbitrary unreduced open observables.
	Thus the datum includes a continuous \(P_0\)-action
	\[
	  \operatorname{act}^{P_0}_I\colon
	  C^\bullet_{\mathrm{CE},\mathrm{cont}}(\mathfrak g_I^{\mathrm{cur}})
	  \widehat\otimes
	  \widetilde{\mathrm{Obs}}^{\mathrm{cl,pp}}_{\partial,N}(I)
	  \longrightarrow
	  \widetilde{\mathrm{Obs}}^{\mathrm{cl,pp}}_{\partial,N}(I),
	\]
	restricting on the reduced current generators to the displayed
	Hamiltonian action on \(B_f(a)\) and \(\Theta_\rho(B)\).  The completed
	open-observable centrality complex is
	\[
	  \mathcal H^\bullet_{\mathrm{cot,open}}(I)
	  =
	  \underline{\mathrm{Hom}}_{\mathcal B^{\mathrm{adm}}}^{\bullet}\!\left(
	  C^\bullet_{\mathrm{CE},\mathrm{cont}}(\mathfrak g_I^{\mathrm{cur}})
	  \widehat\otimes
	  \widetilde{\mathrm{Obs}}^{\mathrm{cl,pp}}_{\partial,N}(I),
	  \widetilde{\mathrm{Obs}}^{\mathrm{cl,pp}}_{\partial,N}(I)\right),
	\]
	with differential induced by \(Q_{\mathrm{open}}\), the CE differential,
	and the \(P_0\)-bracket sign convention of
	Appendix~\ref{app:sign-conventions}.  For a partial unreduced lift
	\(\widetilde\kappa_I\), a source cochain \(x\), and an unreduced open
	observable \(A\), the open-centrality defect is
	\[
	  \Delta^{\mathrm{open}}_{\widetilde\kappa_I}(x;A)
	  =
	  \{\widetilde\kappa_I(x),A\}_{P_0}
	  -
	  \operatorname{act}^{P_0}_I(x;A)
	  \in
	  \mathcal H^\bullet_{\mathrm{cot,open}}(I).
	\]
	The required open-observable centrality primitive is a cochain
	\(H^{\mathrm{open}}_{x,A}\) in the preceding complex, of one lower
	degree, satisfying
	\[
	  d_{\mathrm{cot}}H^{\mathrm{open}}_{x,A}
	  =
	  -\Delta^{\mathrm{open}}_{\widetilde\kappa_I}(x;A),
	\]
	compatibly with interval inclusions, finite-window projections,
	weight changes, and disjoint factorization products.  The phrase
	``against arbitrary unreduced open observables'' means precisely this
	Hom-complex condition; the reduced current brackets alone do not supply
	it.

	Let \(\Theta_I^{\mathrm{red}}\) denote the reduced current \(P_0\)
	prefactorization datum of
Theorem~\ref{thm:reduced-principal-part-boundary-current}.  The
unreduced centrality convolution complex is
\[
  \mathcal C^\bullet_{\mathrm{unred},P_0}(I)
  =
  \underline{\mathrm{Hom}}_{\mathcal B^{\mathrm{adm}}}^{\bullet}\!\left(
  C^\bullet_{\mathrm{CE},\mathrm{cont}}(\mathfrak g_I^{\mathrm{cur}}),
  Z^{P_0}_{\mathrm{fact}}\!
  \bigl(\widetilde{\mathrm{Obs}}^{\mathrm{cl,pp}}_{\partial,N}(I)\bigr)
  \right),
\]
with the completed \(L_\infty\)-structure obtained by linearising the
\(P_0\)-centrality equations at the chosen unreduced source and target
structures.  The reduced current-bracket complex
\(\mathcal C^\bullet_{\mathrm{red},P_0}(I)\) is the corresponding
linearised \(P_0\)-morphism complex at \(\Theta_I^{\mathrm{red}}\), with
target \(A^{\mathrm{pp}}_\partial(I)\).  Primitive projection
\(\pi_{\mathrm{prim}}\) and scalar reduction induce a restriction
morphism
\[
  \rho_{\mathrm{cot},I}\colon
  \mathcal C^\bullet_{\mathrm{unred},P_0}(I)
  \longrightarrow
  \mathcal C^\bullet_{\mathrm{red},P_0}(I).
\]
The cotangent-current deformation complex is the relative homotopy fibre
\[
  \mathcal K^{\mathrm{cot}}_\partial(I)
  :=
  \operatorname{hofib}_{\Theta_I^{\mathrm{red}}}
  \bigl(\rho_{\mathrm{cot},I}\bigr),
\]
	equivalently the shifted cone
	\(\operatorname{Cone}(\rho_{\mathrm{cot},I})[-1]\) after subtracting the
	fixed reduced current datum, enlarged by
	\(\mathcal H^\bullet_{\mathrm{cot,open}}(I)\) for the
	open-observable centrality equations.  Its differential
	\(d_{\mathrm{cot}}\) is the linearised \(P_0\)-centrality differential.
	For a partial unreduced
	lift \(\widetilde\kappa_I\), the cotangent curvature is
	\[
	  \mathcal R^{\mathrm{cot}}_{\partial,I}
  =
  d_{\mathrm{cot}}\widetilde\kappa_I
  +\frac12[\widetilde\kappa_I,\widetilde\kappa_I]_{P_0}
  +\sum_{m\ge3}\frac1{m!}
    \ell_m^{P_0}(\widetilde\kappa_I^{\otimes m})
  \in
  \bigl(\mathcal K^{\mathrm{cot}}_\partial(I)\bigr)^2 .
\]
The cotangent-current obstruction is
\[
  \operatorname{Ob}^{\mathrm{cot}}_\partial(I)
  =
  [\mathcal R^{\mathrm{cot}}_{\partial,I}]
  \in
  H^2\bigl(\mathcal K^{\mathrm{cot}}_\partial(I)\bigr).
\]
Vanishing of this class with a degree-one primitive
\(h_{\mathrm{cot},I}\) is exactly the lift from the reduced
principal-part current realization to the unreduced cotangent current
sector.  The Costello graph version is the weighted refinement of
Definition~\ref{def:app-unreduced-nonscalar-current-lift-datum}.
\end{defn}

\begin{defn}[Brane-defect QME tower datum]
\label{defn:typed-brane-defect-qme-tower}
Fix an interval \(I\subset\R_t\), an admissible weight \(w\), a
scalar-contact filtration row \(r\), a finite Hamiltonian window \(M\),
and a Costello scale \(L\).  A brane-defect quantum master equation
(QME) tower datum consists of
\[
  \bigl(
  Q,Q^{\mathrm{GF}},K_t,P_{\varepsilon,L},\Delta_L,\{-,-\}_L,
  W(P_{L_1,L_2},-),\pi_{M,M'},\mathcal Q^\bullet_{w,r,M,L}(I)
  \bigr),
\]
where \(Q\) is the classical BV differential, \(Q^{\mathrm{GF}}\) is a
gauge-fixing operator, \(K_t\) its heat kernel, \(P_{\varepsilon,L}\) the
Costello propagator, \(\Delta_L\) and \(\{-,-\}_L\) are the finite-scale
BV Laplacian and bracket, \(W(P_{L_1,L_2},-)\) is
renormalisation-group flow, and \(\mathcal Q^\bullet_{w,r,M,L}(I)\) is a
complete filtered complex of local functionals with brane-defect support,
closed under these operations.  The order parameter is the formal
Costello parameter \(\hbar_{\mathrm{QME}}\).

The operation list is
\[
  \mathcal O^{\mathrm{QME}}_{w,r,M,L}
  =
  \{\,Q,\{-,-\}_L,\Delta_L,W(P_{L_1,L_2},-),\pi_{M,M'},
      \widehat\sigma_{\mathrm{sc},w},\partial_{\mathrm{loc}}\,\}.
\]
Here \(\widehat\sigma_{\mathrm{sc},w}\) is the lifted scalar-contact
projection and \(\partial_{\mathrm{loc}}\) denotes the Costello--Gwilliam
locality and factorisation-boundary maps on brane-defect local
functionals.  The kernel list is
\[
  \mathcal K^{\mathrm{ker,QME}}_{w,r,M,L}
  =
  \{\,K_t,P_{\varepsilon,L},P_{L_1,L_2}\,\},
\]
enlarged, when the unreduced cotangent source is invoked, by the supplied
bulk-to-defect kernel \(K_{\hbar_{\mathrm{QME}},w,I}\) of
Theorem~\ref{thm:completion-graph-qme}.  The total QME complex is
\[
  \mathcal Q^\bullet_{\mathrm{QME}}(I)
  =
  \operatorname{Roos}_{(w,r,M,L)}
  \bigl\{\mathcal Q^\bullet_{w,r,M,L}(I),
  \pi_{M,M'},\,W(P_{L_1,L_2},-)\bigr\},
\]
with the finite-scale differential on each vertex and the Roos transition
differential.  The residual non-scalar part is
\(\ker\widehat\sigma_{\mathrm{sc},w}\).  A compatible family
\(\{C_{n,M}\}\) satisfying only \(\pi_{M,N}C_{n,M}=C_{n,N}\) is therefore
only the finite-window face of the datum.

For a scale-\(L\) local interaction
\[
  I^w_M[L]
  =
  I^w_{0,M}[L]
  +\sum_{n\ge1}\hbar_{\mathrm{QME}}^n I^w_{n,M}[L],
  \qquad
  I^w_{n,M}[L]\in\mathcal Q^0_{w,r,M,L}(I),
\]
the finite-scale QME is
\[
  QI^w_M[L]
  +\frac12\{I^w_M[L],I^w_M[L]\}_L
  +\hbar_{\mathrm{QME}}\Delta_LI^w_M[L]=0 .
\]
Assume the equation has been solved through order
\(\hbar_{\mathrm{QME}}^{\,n-1}\).  With
\[
  d^{L}_{\mathrm{cl},M}
  =
  Q+\{I^w_{0,M}[L],-\}_L ,
\]
write
\[
  d^{L}_{\mathrm{QME},n,M}:=d^{L}_{\mathrm{cl},M}
\]
for the order-\(n\) linearised QME coefficient differential.  This is
the differential acting on the unknown \(n\)-th counterterm.  The
completed quantum differential
\(d_{\mathrm{QME}}^{L}=Q+\{I^w_M[L],-\}_L+\hbar_{\mathrm{QME}}\Delta_L\)
belongs to the assembled formal interaction, not to the primitive
equation for a single obstruction coefficient.
The order-\(n\) QME curvature is
\[
  \mathcal R^{\mathrm{QME}}_{n,w,r,M,L}
  =
  \frac12
  \sum_{\substack{i+j=n\\0<i,j<n}}
  \{I^w_{i,M}[L],I^w_{j,M}[L]\}_L
  +\Delta_LI^w_{n-1,M}[L]
  \in\mathcal Q^1_{w,r,M,L}(I).
\]
The Costello Bianchi identity gives
\[
  d^{L}_{\mathrm{QME},n,M}
  \mathcal R^{\mathrm{QME}}_{n,w,r,M,L}=0 .
\]
An order-\(n\) counterterm is a primitive
\[
  C_{n,w,r,M,L}\in\mathcal Q^0_{w,r,M,L}(I),
  \qquad
  d^{L}_{\mathrm{QME},n,M}C_{n,w,r,M,L}
  =-\mathcal R^{\mathrm{QME}}_{n,w,r,M,L}.
\]
The tower also contains scale and window coherence primitives
\[
  K^{\mathrm{RG}}_{n,w,r,M;L_1,L_2},\qquad
  K^{\mathrm{win}}_{n,w,r;M,M',L},
\]
whose differentials kill the defects of \(C_{n,w,r,M,L}\) under
\(W(P_{L_1,L_2},-)\) and \(\pi_{M,M'}\).
Explicitly,
\[
  \begin{aligned}
  \mathcal D^{\mathrm{RG}}_{n,w,r,M;L_1,L_2}
  &=
  W(P_{L_1,L_2},C_{n,w,r,M,L_1})-C_{n,w,r,M,L_2},\\
  \mathcal D^{\mathrm{win}}_{n,w,r;M,M',L}
  &=
  \pi_{M,M'}C_{n,w,r,M,L}
  -C_{n,w,r,M',L}\pi^{\mathrm{CE}}_{M,M'} .
  \end{aligned}
\]

The all-order QME counterterm obstruction is the class of the family
\[
  \left(
  \{\mathcal R^{\mathrm{QME}}_{n,w,r,M,L}\}_{n,w,r,M,L},
  \{\mathcal D^{\mathrm{RG}}_{n,w,r,M;L_1,L_2}\},
  \{\mathcal D^{\mathrm{win}}_{n,w,r;M,M',L}\}
  \right)
\]
in the corresponding total Roos complex
\(\mathcal Q^\bullet_{\mathrm{QME}}(I)\); it is denoted
\[
  \operatorname{Ob}^{\mathrm{QME}}_{\mathrm{all\mbox{-}order}}
  \in H^1(\mathcal Q^\bullet_{\mathrm{QME}}(I)).
\]
Vanishing means that all order, scale, and window curvatures have the
displayed primitives.  The Costello loop parameter
\(\hbar_{\mathrm{QME}}\) is distinct from the raw Rees/Weyl commutator
parameter \(\hbar_{\mathrm{W}}\); neither parameter occurs in the
canonical constant-Hamiltonian source extension.
\end{defn}

\begin{defn}[Comparison-map deformation complex]
\label{defn:comparison-map-deformation-complex}
Let \((\mathfrak K_i,Q_i)\) and
\((\mathfrak K_{\mathrm{HT}},Q_{\mathrm{HT}})\) be complete filtered
\(L_\infty\)-algebras, written in the suspended coalgebra convention.
Put
\[
  \mathcal C^\bullet_{\chi_i}
  =
  \prod_{k\ge1}
  \underline{\operatorname{Hom}}^{\bullet}_{\mathrm{filt}}
  \left(
    \widehat{\operatorname{Sym}}^k(\mathfrak K_i[1]),
    \mathfrak K_{\mathrm{HT}}[1]\right).
\]
A sequence
\(\chi_i=(\chi_{i,k})_{k\ge1}\in\mathcal C^0_{\chi_i}\) determines a
filtered coalgebra map
\[
  \mathcal X_i\colon
  \widehat{\operatorname{Sym}}^c(\mathfrak K_i[1])
  \longrightarrow
  \widehat{\operatorname{Sym}}^c(\mathfrak K_{\mathrm{HT}}[1]).
\]
Its \(L_\infty\)-morphism defect is the cogenerator projection of
\[
  \mathcal R_{\mathcal X_i}
  =
  Q_{\mathrm{HT}}\mathcal X_i-\mathcal X_iQ_i .
\]
Thus \(\chi_i\) is a filtered \(L_\infty\)-morphism exactly when
\(\mathcal R_{\mathcal X_i}=0\).  If the defect vanishes in all filtration
orders \(<m+1\), the arity/filtration-\((m+1)\) linearized differential
\[
  D^{(m)}_{\chi_i}\alpha
  =
  \left.
  \frac{d}{dt}\right|_{t=0}
  \mathcal R_{\mathcal X_i+t\alpha}\big|_{\operatorname{gr}^{m+1}}
  \qquad
  \left(
    \alpha\in
    \operatorname{gr}^{m+1}\mathcal C^\bullet_{\chi_i}
  \right)
\]
satisfies \((D^{(m)}_{\chi_i})^2=0\).  The corresponding comparison-map
obstruction is
\[
  \operatorname{Ob}_{\chi_i,m+1}
  =
  [\mathcal R_{\chi_i,m+1}]
  \in
  H^1\!\left(
    \operatorname{gr}^{m+1}\mathcal C^\bullet_{\chi_i},
    D^{(m)}_{\chi_i}\right).
\]
It vanishes exactly when an order-\((m+1)\) component
\(\chi_{i,m+1}\) exists with
\[
  D^{(m)}_{\chi_i}\chi_{i,m+1}
  =
  -\mathcal R_{\chi_i,m+1}.
\]
For the six native comparison maps at a fixed obstruction order \(m+1\),
set
\[
  \operatorname{gr}^{m+1}\mathcal C^\bullet_{\chi}
  =
  \prod_i\operatorname{gr}^{m+1}\mathcal C^\bullet_{\chi_i},
  \qquad
  D^{(m)}_\chi=\prod_iD^{(m)}_{\chi_i}.
\]
\end{defn}

\begin{lem}[Comparison-map defect cocycle]
\label{lem:comparison-map-defect-cocycle}
In the setting of
Definition~\ref{defn:comparison-map-deformation-complex}, assume
\(\mathcal R_{\mathcal X_i}\equiv0\) modulo \(F^{m+1}\).  Then the
\((m+1)\)-st defect \(\mathcal R_{\chi_i,m+1}\) is
\(D^{(m)}_{\chi_i}\)-closed.  Hence
\(\operatorname{Ob}_{\chi_i,m+1}\) is a well-defined cohomology class.
\end{lem}

\begin{proof}
The \(L_\infty\)-algebra structures are square-zero coderivations:
\[
  Q_i^2=0,\qquad Q_{\mathrm{HT}}^2=0 .
\]
For any filtered coalgebra map \(\mathcal X_i\),
\[
  Q_{\mathrm{HT}}\mathcal R_{\mathcal X_i}
  +\mathcal R_{\mathcal X_i}Q_i
  =
  Q_{\mathrm{HT}}(Q_{\mathrm{HT}}\mathcal X_i-\mathcal X_iQ_i)
  +(Q_{\mathrm{HT}}\mathcal X_i-\mathcal X_iQ_i)Q_i
  =
  0 .
\]
This is the \(L_\infty\)-morphism Bianchi identity.  Reducing modulo
\(F^{m+2}\), all terms nonlinear in the \((m+1)\)-st unknown contain a
lower defect and vanish by the hypothesis
\(\mathcal R_{\mathcal X_i}\equiv0\) modulo \(F^{m+1}\).  The remaining
degree-\((m+1)\) component is
\[
  D^{(m)}_{\chi_i}\mathcal R_{\chi_i,m+1}=0 .
\]
The identity \((D^{(m)}_{\chi_i})^2=0\) follows by applying the same
Bianchi identity to the first variation of the defect.  Therefore
\(\mathcal R_{\chi_i,m+1}\) defines the displayed cohomology class.
\end{proof}

\begin{defn}[Mixed-HT Deligne equivalence datum]
\label{defn:typed-mixed-ht-deligne-equivalence-datum}
Fix \(N\), a brane-link boundary \(\partial X\), the admissible host
category \(\mathcal B^{\mathrm{adm}}\), the bulk object
\(\mathcal A^{\mathrm{cl}/\hbar}_{\mathrm{bulk}}\), and the Morita
centre
\[
  Z_\partial=
  Z^{\mathrm{der},\mathrm{Mor}}_{E_2^{\mathrm{HT}}}
  (\mathcal C^{\mathrm{op}}_\partial).
\]
Let
\[
  \Theta^{\mathrm{tr}}_{\leq 2}
  =
  \bigl(\tau_{N,1}^{\mathrm{tr}},K^{(2)}_\Delta\bigr).
\]
The first entry has the coordinate-dual expression
\[
  \Phi_N^{\mathrm{gen}}(c_f)=\theta_f,\quad
  \Phi_N^{\mathrm{gen}}(u_f)=M_{J(f)},
\]
where \(J(f)=\varprojlim_W J_{\leq W}(f_{\leq W})\).  The formal
Hamiltonian-current tree-kernel suboperad
\[
  \mathsf Q=\mathsf{PP}^{(2)}_{\mathrm{tree}}
\]
acts on the closed Hamiltonian-current sector by
Theorem~\ref{app:thm:formal-polydisk-all-arity-current-operad}; denote
this constructed restricted action by \(\Theta_{\mathsf Q}^{\mathrm{cur}}\).
It is the closed part of \(\Theta^{\mathrm{tr}}_{\leq2}\).  The full
two-coloured \(\mathrm{SC}^{\mathrm{HT}}_{(1,2)}\) action on the
bulk--boundary pair and the Morita-centre action on \(Z_\partial\) are
the residual Deligne rows below.
Let
\[
  \mathcal A^{\mathrm{Ham\mbox{-}cur}}_{\mathrm{bulk}}
  \subset
  \mathcal A^{\mathrm{cl}/\hbar}_{\mathrm{bulk}}
\]
be the formal Hamiltonian-current factor on which \(\mathsf Q\) acts.
Let
\[
  \iota_{\mathsf Q}\colon
  \mathsf Q
  \longrightarrow
  \mathrm{SC}^{\mathrm{HT}}_{(1,2)}(\mathsf c^\bullet;\mathsf c)
\]
be the inclusion of the tree principal-part sector into the closed
colour, and let
\[
  \pi_{\mathrm{cur}}\colon
  \mathcal A^{\mathrm{cl}/\hbar}_{\mathrm{bulk}}
  \longrightarrow
  \mathcal A^{\mathrm{Ham\mbox{-}cur}}_{\mathrm{bulk}}
\]
be the formal Hamiltonian-current projection.  Let
\(\mathfrak{Def}_{\mathrm{tr},\leq2}\) denote the finite arity
deformation complex whose degree-one Maurer--Cartan points are the
closed binary diagonal kernel together with the first Taylor coefficient
\(\tau_{N,1}^{\mathrm{tr}}\).  Its coordinate-dual values are
\(c_f\mapsto\theta_f\) and \(u_f\mapsto M_{J(f)}\).  Precomposition
with \(\iota_{\mathsf Q}\), postcomposition with \(\pi_{\mathrm{cur}}\),
and projection to the trace arity-\(\leq2\) face give restriction
morphisms
\[
\begin{aligned}
  \rho_{\mathrm{tr},\leq2}&\colon
  F^1\mathfrak{Def}^{\mathrm{HT}}_{\mathrm{SC}}
  (\mathcal A^{\mathrm{cl}/\hbar}_{\mathrm{bulk}},
   \mathcal C^{\mathrm{op}}_\partial)
  \longrightarrow
  \mathfrak{Def}_{\mathrm{tr},\leq2},\\
  \rho_{\mathsf Q}&\colon
  F^1\mathfrak{Def}^{\mathrm{HT}}_{\mathrm{SC}}
  (\mathcal A^{\mathrm{cl}/\hbar}_{\mathrm{bulk}},
   \mathcal C^{\mathrm{op}}_\partial)
  \longrightarrow
  \mathrm{Conv}\bigl(\mathrm B\mathsf Q,
    \operatorname{End}_{\mathcal A^{\mathrm{Ham\mbox{-}cur}}_{\mathrm{bulk}}}
  \bigr).
\end{aligned}
\]
Here \(\rho_{\mathrm{tr},\leq2}\) records the closed binary diagonal
kernel and first Taylor boundary value, while \(\rho_{\mathsf Q}\) records
all closed tree-current arities.  Their common binary face is
\(K^{(2)}_\Delta\).  The relative Swiss-cheese deformation complex over
the formal-Darboux fibre is
\[
  \mathfrak{Def}^{\mathrm{HT}}_{\mathrm{SC,rel}}
  :=
  \operatorname{hofib}_{(\Theta^{\mathrm{tr}}_{\leq2},
  \Theta_{\mathsf Q}^{\mathrm{cur}})}
  \bigl(\rho_{\mathrm{tr},\leq2},\rho_{\mathsf Q}\bigr).
\]

A mixed-HT Deligne equivalence datum extending this formal-Darboux fibre
is a tuple
\[
\begin{aligned}
  \mathfrak D^{\mathrm{HT}}_N
  =
  (&\Theta_{\mathrm{SC}},\Theta_{\mathrm{Mor}},
    \mathcal C_{\Phi_N^{\mathrm{Hoch}}},
    \nabla_{\mathrm{Darb}},\{\kappa^{\mathrm{cot}}_I,h_{\mathrm{cot},I}\}_I,
    \mathfrak F_{\mathrm{desc}},\\
   &\{I^w_{n,M}[L],C_{n,w,r,M,L},
      K^{\mathrm{RG}}_{n,w,r,M;L_1,L_2},
      K^{\mathrm{win}}_{n,w,r;M,M',L}\},\\
   &h_{\mathrm{anom}},h_{\mathrm{loc}},
    (\mathfrak K_{\mathrm{HT}},\chi_{\mathrm{nat}}),
    \mathfrak C_{\mathrm{prim}},
    \kappa_{\operatorname{Cone}}).
\end{aligned}
\]
The entries satisfy the following equations.
\begin{enumerate}[label=\textup{(D\arabic*)}]
	  \item \emph{All-arity Swiss-cheese action.}
	    \[
	      \Theta_{\mathrm{SC}}
	      \in
	      \mathrm{MC}\bigl(
	      \mathfrak{Def}^{\mathrm{HT}}_{\mathrm{SC,rel}}\bigr)
	    \]
	    is a Maurer--Cartan element in the complete arity-filtered relative
	    mixed-HT Swiss-cheese deformation complex.  Equivalently, it is a
	    Maurer--Cartan element in the absolute complex whose restrictions
	    satisfy
	    \[
	      \rho_{\mathrm{tr},\leq2}(\Theta_{\mathrm{SC}})
	      =
	      \Theta^{\mathrm{tr}}_{\le2},
	      \qquad
	      \rho_{\mathsf Q}(\Theta_{\mathrm{SC}})
	      =
	      \Theta_{\mathsf Q}^{\mathrm{cur}} .
		    \]
		    Equivalently, after the trace arity-\(\le2\) part and the
		    constructed \(\mathsf Q\)-restricted current part are fixed, at the
		    inductive stage \(n\) the residual obstruction representative
	    \(\mathcal R^{\mathrm{SC}}_{n+1}(\Theta_{\leq n})\) defines the
	    relative class
	    \(\operatorname{Ob}^{\mathrm{SC}/(\mathrm{tr},\mathsf Q)}_{n+1}\)
	    of
	    Definition~\ref{defn:relative-swiss-cheese-extension-obstruction},
	    and this class has a degree-one primitive
    \[
      d_{\Theta_{\leq n}}H^{\mathrm{SC}}_{n+1}
      =-\mathcal R^{\mathrm{SC}}_{n+1}(\Theta_{\leq n}).
    \]
    Concretely, this row contains, for every finite colour-labelled set
    \(I\), continuous operation maps
      \[
        \mu^{\Theta}_{I}\colon
        \mathrm{SC}^{\mathrm{HT}}_{(1,2)}(I)\otimes^!
      \left(\boxtimes_{i\in I}\mathcal X_i\right)
      \longrightarrow
      \Delta_I^*\mathcal X_{\mathrm{out}},
    \]
    where each \(\mathcal X_i\) is either the bulk object
    \(\mathcal A^{\mathrm{cl}/\hbar}_{\mathrm{bulk}}\) or the boundary
    object \(\mathcal C^{\mathrm{op}}_\partial\), and
    \(\Delta_I\) is the small diagonal on the appropriate
    Costello--Gwilliam/Ran configuration space.  The maps are
    \(\mathfrak S_I\)-equivariant, unital, compatible with disjoint
    union, and compatible with every partition
    \(I=\bigsqcup_{a\in A}I_a\) through the operadic composition along
    the partial diagonal
    \(\Delta_{\{I_a\}}\subset(\C^2)^I\).  On the closed Hamiltonian-current
    face the coefficient object is the local-cohomology tree complex
    \[
      \mathcal P^{\mathrm{tree}}_I
      =
      \left(
      \bigoplus_{T\in\operatorname{Tree}(I)}
      \C\,K_T
      \right)\Big/
      (\mathrm{antisymmetry},\mathrm{Arnold\mbox{-}Jacobi},
       \check C\mathrm{ech}\mbox{-}\mathrm{coboundary}),
    \]
    where
    \[
      K_T=
      \bigwedge_{e\in E_{\mathrm{int}}(T)}
      \frac{d\zeta_{e,1}\,d\zeta_{e,2}}
           {\zeta_{e,1}\zeta_{e,2}}
      \in
      H^{2(|I|-1)}_{\Delta_T}
      \bigl(\mathcal O_{(\widehat{\C^2}_0)^I}\bigr)
      \otimes\omega_{(\C^2)^I/\C^2}.
    \]
    The restriction equation
    \(\rho_{\mathsf Q}(\Theta_{\mathrm{SC}})
    =\Theta_{\mathsf Q}^{\mathrm{cur}}\) means that for every rooted
    \(I\)-tree \(T\)
    \[
      \mu_T^\Theta(K_T;x_1,\ldots,x_{|I|})
      =
      K_T\,\mathsf J_{[-,\ldots,-]_T(x_1,\ldots,x_{|I|})},
      \qquad x_i\in\mathfrak g,
    \]
    as in Theorem~\ref{app:thm:formal-polydisk-all-arity-current-operad}.
    For a partition \(I=\bigsqcup_{a\in A}I_a\), rooted trees \(T_a\)
    on \(I_a\), and a rooted tree \(T_A\) on \(A\), partition
    compatibility contains the residue identity
    \[
      \operatorname{Res}_{T_A}
      \operatorname{Res}_{\{T_a\}_{a\in A}}
      \left(K_{T_A}\wedge\bigwedge_{a\in A}K_{T_a}\right)
      =
      \operatorname{Res}_{T_A\circ\{T_a\}}
      K_{T_A\circ\{T_a\}},
    \]
    together with the corresponding iterated bracket identity in
    \(\mathfrak g\).  The arity-three tree-current curvature vanishes by
    Grothendieck residue Fubini on the small diagonal of \((\C^2)^3\)
    and by the graded Jacobi identity; the remaining mixed/open
    arity-three component is killed by \(H^{\mathrm{SC}}_3\).  The
    arity-four tree-current pentagon is the five-parenthesization
    full-diagonal residue identity of
    Lemma~\ref{app:lem:arity-four-residue-pentagon}; the residual
    four-ary mixed/open face is killed by \(H^{\mathrm{SC}}_4\) subject
    to the pentagon relation.  The general
    obstruction \(\mathfrak o^{\mathrm{SC}}_{n+1}\)
    is the class of
    \[
      \mathcal R^{\mathrm{SC}}_{n+1}
      =
      \left(d\mu^\Theta_{\le n}
      +\frac12[\mu^\Theta_{\le n},\mu^\Theta_{\le n}]\right)_{F^{n+1}/F^{n+2}}
    \]
    after projection to the quotient by the closed tree-current summand
    already killed by residue Fubini and Jacobi.  Thus this row is not
    the arity-two kernel \(K^{(2)}_\Delta\); it is the full system of
    symmetric, partition-compatible, all-arity open--closed operations
    extending that kernel and the constructed tree-current face.

	  \item \emph{Morita mixed-HT \(E_2^{\mathrm{HT}}\)-centre action.}
    \[
      \Theta_{\mathrm{Mor}}
      \in
      \mathrm{MC}\!\left(
      F^1\mathfrak{Def}^{\mathrm{HT}}_{\mathrm{Mor}\mbox{-}E_2}
      (\mathcal C^{\mathrm{op}}_\partial,b_p)\right)
    \]
    is a Maurer--Cartan point of the shifted homotopy fibre
    \[
      \operatorname{fib}\!\left[
      \mathrm{Conv}(\mathrm B\mathsf P,\operatorname{End}_{Z_\partial})
      \xrightarrow{\operatorname{res}_{b_p}}
      \mathrm{Conv}\!\left(\mathrm B\mathsf Q,
      \operatorname{End}_{\operatorname{ev}_{b_p}(Z_\partial)^{\mathrm{tr}}}
      \right)\right][-1]
    \]
    from Definition~\ref{app:defn:morita-chiral-E2-extension-complex}.
    Thus it is a relative lift of the \(\mathsf Q\)-restricted formal
    Hamiltonian-current action \(\Theta_{\mathsf Q}^{\mathrm{cur}}\) to
    the Morita centre \(Z_\partial\).  At filtration stage \(r\), the
    residual four-face curvature defines
    \[
      \operatorname{Ob}^{\mathrm{Mor}\mbox{-}E_2}_{r+1}
      \in
      H^2\!\left(
      \operatorname{gr}^{r+1}_{F}
      \mathfrak{Def}^{\mathrm{HT}}_{\mathrm{Mor}\mbox{-}E_2}
      (\mathcal C^{\mathrm{op}}_\partial,b_p)\right)
    \]
    of
    Definition~\ref{defn:relative-morita-e2-extension-obstruction}.
    Vanishing means that the primitive \(h_{E_2}\) exists with
    \[
      d_{E_2}h_{E_2}
      =
      -\mathcal R^{\mathrm{ch}\,E_2}_{\mathrm{Deligne}},
      \qquad
      d_{E_2}=d_F\oplus d_C\oplus d_R\oplus d_{\mathrm{Mor}} .
    \]
    This row includes the construction of the mixed-HT factorization
    bimodule category
    \[
      \mathcal M^{\mathrm{HT}}_\partial
      =
      \mathrm{Bimod}^{\mathrm{HT}}_{\operatorname{Ran}(\partial X)}
      (\mathcal C^{\mathrm{op}}_\partial,
       \mathcal C^{\mathrm{op}}_\partial)
    \]
    of complete brane-link Ran factorization bimodules with continuous
    left and right
    \(\widehat{\mathcal C}^{\mathrm{op}}_{\partial,\mathrm{adm}}\)-actions
    and Weiss descent.  Its diagonal object is
    \(\Delta_\partial\), the identity
    \((\mathcal C^{\mathrm{op}}_\partial,
    \mathcal C^{\mathrm{op}}_\partial)\)-bimodule.  The Morita centre is
    the derived endomorphism object
    \[
      Z_\partial
      =
      \operatorname{REnd}_{\mathcal M^{\mathrm{HT}}_\partial}
      (\Delta_\partial)
      =
      \operatorname{RHom}_{\mathcal M^{\mathrm{HT}}_\partial}
      (\Delta_\partial,\Delta_\partial),
    \]
    before a brane vacuum is chosen.  The Dirac brane \(b_p^{\oplus N}\)
    gives \(A_{b_p}=A^{\mathrm{cl}}_{\partial,N}\) and a restriction
    functor
    \[
      r_{b_p}\colon
      \mathcal M^{\mathrm{HT}}_\partial
      \longrightarrow
      \mathrm{Bimod}^{\mathrm{HT}}_{\mathrm{adm}}
      (A_{b_p},A_{b_p}),
      \qquad
      r_{b_p}(\Delta_\partial)=\Delta_{A_{b_p}} .
    \]
    Applying derived endomorphisms gives the Hochschild chart
    \[
      \operatorname{ev}_{b_p}\colon
      Z_\partial
      \longrightarrow
      \operatorname{REnd}_{\mathrm{Bimod}^{\mathrm{HT}}_{\mathrm{adm}}
      (A_{b_p},A_{b_p})}(\Delta_{A_{b_p}})
      =
      C^\bullet_{\mathrm{ch},\mathrm{HT},\mathrm{adm}}
      (A^{\mathrm{cl}}_{\partial,N},A^{\mathrm{cl}}_{\partial,N}).
    \]
    The Hochschild chart is only the evaluation of the centre at
    \(b_p\); it becomes an equivalence only under the Morita-density
    condition of
    Proposition~\ref{prop:typed-morita-centre-hochschild-chart}.  The
    precise obstruction objects are the density quotient
    \(\mathcal Q_{b_p}^{\mathrm{Mor}}\) and the chart cone
    \(\operatorname{Cone}(\operatorname{ev}_{b_p})\) of
    Definition~\ref{defn:morita-density-quotient-chart-cone}.  These
    are not consequences of the trace-sector formal-Darboux theorem.

    The curvature killed by \(h_{E_2}\) has four faces:
    \[
      \mathcal R^{\mathrm{ch}\,E_2}_{\mathrm{Deligne}}
      =
      (\mathcal R_F,\mathcal R_C,\mathcal R_R,\mathcal R_{\mathrm{Mor}}).
    \]
    Here \(\mathcal R_F\) represents \(\mathrm{Ob}_F\), the obstruction
    to promoting \(\mathsf Q=\mathsf{PP}^{(2)}_{\mathrm{tree}}\) to the
    chosen cofibrant
    \(\mathsf P=\mathsf hE_2^{\mathrm{ch},\mathrm{HT}}\)-model;
    \(\mathcal R_C\) represents \(\mathrm{Ob}_C\), the obstruction to a
    homotopy-central action of \(Z_\partial\) on
    \(\mathcal C^{\mathrm{op}}_\partial\);
    \(\mathcal R_R\) is the class of the Ran/Weiss comparison cone
    \(\operatorname{Cone}(\chi_{\mathrm{Ran}})\); and
    \(\mathcal R_{\mathrm{Mor}}\) is the Maurer--Cartan curvature in
    \(\operatorname{fib}(\operatorname{res}_{b_p})[-1]\).  Thus
    \(h_{E_2}=(h_F,h_C,h_R,h_{\mathrm{Mor}})\) means
    \[
      d_F h_F=-\mathcal R_F,\qquad
      d_C h_C=-\mathcal R_C,\qquad
      d_R h_R=-\mathcal R_R,\qquad
      d_{\mathrm{Mor}} h_{\mathrm{Mor}}=-\mathcal R_{\mathrm{Mor}},
    \]
    together with compatibility under \(r_{b_p}\), finite arity
    truncation, admissible completion, and Weiss restriction.  The
    primitive \(h_{E_2}\) therefore supplies the promotion from the
    tree-current suboperad to the full mixed-HT chiral coefficient
    system, the homotopy-central Morita action, and the Ran/Weiss
    comparison needed before local Hochschild charts can glue.

  \item \emph{Hochschild-centre chart lift.}
    \(\mathcal C_{\Phi_N^{\mathrm{Hoch}}}\) is the Maurer--Cartan element
    of
    \(\mathcal K^{\mathrm{cent}}_{\partial,N}\) in
    Definition~\ref{defn:typed-hochschild-centre-lift-datum}.  This row
    separates the prescribed first Taylor coefficient from the full
    Hochschild twisting cochain.  It requires a continuous filtered
    twisting cochain
    \[
      \tau_N=\mathcal C_{\Phi_N^{\mathrm{Hoch}}}
      =
      \tau_{N,1}
      +H^{\mathrm{Lie}}+H^{\mathrm{act}}
      +H^{\mathrm{prod}}+H^{P_0}+H^3+\sum_{m\ge4}H^m
      \in
      \mathcal K^{\mathrm{cent}}_{\partial,N},
    \]
    where
    \[
      \mathcal K^{\mathrm{cent}}_{\partial,N}
      \subset
      \underline{\operatorname{Hom}}_{\mathcal B^{\mathrm{adm}}}
      \!\left(
	        \overline C^{\mathrm{CE},\mathrm{cont}}_\bullet
	          (\mathfrak g_{\mathrm{src}}),\,
        C^\bullet_{\mathrm{Hoch},\mathrm{adm}}(A,A)[1]\right)
    \]
    preserves
    \(d_{\mathrm{Hoch}},\smile,[\, ,\,]_G,M_{(-)},\rho,\rho^{P_0}\),
    finite-window projections, and open-input restrictions.  The full
    lift is the \(L_\infty\) Maurer--Cartan equation
    \[
      \sum_{\ell\ge1}\frac1{\ell!}
      \ell_\ell\!\left(
        \mathcal C_{\Phi_N^{\mathrm{Hoch}}}^{\otimes \ell}\right)=0.
    \]
    For \(x,y\in\mathfrak g_{\mathrm{src}}\) and an open observable \(a\), the
    primary defects are
    \[
      \begin{aligned}
      \delta^{\mathrm{Lie}}_{x,y}
	        &=[\tau_{N,1}(x),\tau_{N,1}(y)]_G
	          -\tau_{N,1}([x,y]),\\
	      \delta^{\mathrm{act}}_{x|y}
	        &=[\tau_{N,1}(x),M_{O_y}]_G-M_{\rho_x(O_y)},\\
	      \delta^{\mathrm{prod}}_{x,a}
	        &=[\tau_{N,1}(x),M_a]_G-M_{\rho_x(a)},\\
	      \delta^{P_0}_{x,a}
	        &=\{\tau_{N,1}(x),M_a\}_{P_0}
          -M_{\rho_x^{P_0}(a)} .
      \end{aligned}
    \]
    The first centrality equations are
    \[
      d_{\mathrm{cent}}H^{\mathrm{Lie}}=-\delta^{\mathrm{Lie}},\quad
      d_{\mathrm{cent}}H^{\mathrm{act}}=-\delta^{\mathrm{act}},\quad
      d_{\mathrm{cent}}H^{\mathrm{prod}}=-\delta^{\mathrm{prod}},\quad
      d_{\mathrm{cent}}H^{P_0}=-\delta^{P_0},
    \]
    and the signed Jacobiator/associator defect satisfies
    \(d_{\mathrm{cent}}H^3=-\delta^3\).  At arity \(m+1\), the
    obstruction representative is the associated-graded component
    \[
      \mathcal R^{\mathrm{cent}}_{m+1}
      =
      \left(
      \sum_{\ell\ge1}\frac1{\ell!}
      \ell_\ell\!\left(
        \mathcal C_{\le m}^{\otimes \ell}\right)\right)_{m+1}
      \in
      \operatorname{gr}^{m+1}_{F}
      \mathcal K^{\mathrm{cent}}_{\partial,N},
    \]
    where \(F^\bullet\) is the centrality arity filtration and
    \(\mathcal C_{\le m}\) is the truncation of
    \(\mathcal C_{\Phi_N^{\mathrm{Hoch}}}\) through arity \(m\).  The
    next primitive must satisfy
    \(d_{\mathrm{cent}}H^{m+1}=-\mathcal R^{\mathrm{cent}}_{m+1}\).
    Equivalently, after fixing the first Taylor boundary value
    \(\tau_{N,1}\), the relative class
    \[
      \operatorname{Ob}^{\mathrm{cent}/\mathrm{gen}}_{m+1}
      \in
      H^2\!\left(
      \operatorname{gr}^{m+1}_{F}
      \mathfrak D_{\mathrm{Hoch}/\mathrm{gen}}\right)
    \]
    of Definition~\ref{defn:relative-hochschild-centre-extension-obstruction}
    vanishes.  Thus the Hochschild-centre lift extends the fixed
    trace-sector first Taylor coefficient.
    In unsuspended Hochschild degrees \(H^m\) has degree \(1-m\).  The
    finite-window, open-input, and Roos transition defects of the
    primitives are part of the secondary compatibility row
    \textup{(D9)}.  Only after this tower is solved does the induced
    twisting cochain receive the name \(\Phi_N^{\mathrm{Hoch}}\).  The
    displayed formula \(c_f\mapsto\theta_f,\ u_f\mapsto M_{J(f)}\) is
    obtained afterward as the coordinate-dual expression
    \(\Phi_N^{\mathrm{gen}}\) of \(\tau_{N,1}^{\mathrm{tr}}\).

  \item \emph{Darboux-torsor transport.}
    \(\nabla_{\mathrm{Darb}}\) is the flat transport datum of
    Definition~\ref{defn:hol-bd-datum-HT}; starting from a preconnection
    \(\nabla_0\) on the bar-centre family over
    \(\mathcal T_{\mathrm{Darb}}(b)\), it consists of primitives in
    \(\mathcal E^\bullet_{\mathrm{Darb}}\) for
    \[
      \Omega_{\nabla}=\nabla_0^2,\qquad
      \mathcal R_{\nabla,\Theta}=\nabla_0\Theta ,
    \]
    equivalently
    \[
      \operatorname{Ob}^{\mathrm{Darb}}
      =
      \bigl([\Omega_{\nabla}],
      [\mathcal R_{\nabla,\Theta}]\bigr)=0,
      \qquad
      d_{\mathrm{Darb}}h_{\nabla}=-\Omega_{\nabla},\quad
      d_{\mathrm{Darb}}h_{\nabla,\Theta}
      =-\mathcal R_{\nabla,\Theta}.
    \]
    This is a datum over the Darboux-chart torsor, not a statement in
    one coordinate system.  When this row is required to extend the
    strict trace-sector transport of
    Proposition~\ref{prop:trace-sector-darboux-flat-transport}, it also
    includes a trace receiver \(r_{\mathrm{tr}}\) and the vanishing of the
    relative class
    \[
      \operatorname{Ob}^{\mathrm{Darb}/\mathrm{tr}}
      =
      \bigl([(\Omega_\nabla,0)],
            [(\mathcal R_{\nabla,\Theta},0)]\bigr)
    \]
    in the homotopy fibre complex
    \(\mathcal E^\bullet_{\mathrm{Darb}/\mathrm{tr}}\) of
    Definition~\ref{defn:relative-darboux-extension-obstruction}.  The
    full primitives must restrict to zero trace-sector primitives for a
    strict extension, or to trace-exact primitives for extension up to
    trace-sector gauge homotopy.  Let
    \[
      \mathsf G_{\mathrm{Darb}}
      =
      \mathcal T_{\mathrm{Darb}}(b)\times
      G_{\mathrm{Darb}}
      \rightrightarrows
      \mathcal T_{\mathrm{Darb}}(b)
    \]
    be the formal symplectomorphism action groupoid.  A chart
    \(\alpha\in\mathcal T_{\mathrm{Darb}}(b)\) carries a local complex
    \(\mathcal E^\bullet_{\mathrm{Darb},\alpha}\), a local
    preconnection \(\nabla_{\alpha,0}\), and local primitives
    \(h_{\nabla,\alpha},h_{\nabla,\Theta,\alpha}\).  A transition
    \(g_{\alpha\beta}\in G_{\mathrm{Darb}}\) identifies
    \(g_{\alpha\beta}^*\mathcal E^\bullet_{\mathrm{Darb},\alpha}\) with
    \(\mathcal E^\bullet_{\mathrm{Darb},\beta}\).  The Darboux descent
    complex is the total nerve/Roos complex
    \[
      \mathcal C^\bullet_{\mathrm{Darb}}
      =
      \operatorname{Tot}\,
      C^\bullet\!\left(
      N_\bullet\mathsf G_{\mathrm{Darb}},
      \mathcal E^\bullet_{\mathrm{Darb}}\right),
      \qquad
      D_{\mathrm{Darb}}=d_{\mathrm{Darb}}
      +(-1)^{|\cdot|}\delta_{\mathsf G}.
    \]
    The overlap part of \(\nabla_{\mathrm{Darb}}\) consists of
    degree-zero secondary primitives
    \[
      K^{\mathrm{Darb}}_{\alpha\beta}
      =
      (K^\nabla_{\alpha\beta},K^{\nabla,\Theta}_{\alpha\beta})
      \in
      \mathcal E^0_{\mathrm{Darb},\beta}
      \oplus
      \mathcal E^0_{\mathrm{Darb},\beta}
    \]
    satisfying
    \[
      d_{\mathrm{Darb}}K^\nabla_{\alpha\beta}
      =
      -\bigl(g_{\alpha\beta}^*h_{\nabla,\alpha}
        -h_{\nabla,\beta}\bigr),
      \qquad
      d_{\mathrm{Darb}}K^{\nabla,\Theta}_{\alpha\beta}
      =
      -\bigl(g_{\alpha\beta}^*h_{\nabla,\Theta,\alpha}
        -h_{\nabla,\Theta,\beta}\bigr).
    \]
    On a triple overlap the first Darboux descent obstruction is the
    cocycle
    \[
      \mathcal R^{\mathrm{Darb}}_{\alpha\beta\gamma}
      =
      g_{\beta\gamma}^*K^{\mathrm{Darb}}_{\alpha\beta}
      +K^{\mathrm{Darb}}_{\beta\gamma}
      -K^{\mathrm{Darb}}_{\alpha\gamma}
      \in
      \mathcal C^2_{\mathrm{Darb}},
    \]
    and it is part of the D4 datum to choose a degree-\(-1\) tertiary
    primitive
    \[
      d_{\mathrm{Darb}}L^{\mathrm{Darb}}_{\alpha\beta\gamma}
      =
      -\mathcal R^{\mathrm{Darb}}_{\alpha\beta\gamma}.
    \]
    Higher simplex defects in
    \(N_\bullet\mathsf G_{\mathrm{Darb}}\) are killed by the
    corresponding components of the global primitive-coherence cochain
    \(P_{\mathrm{glob}}\) in row \textup{(D9)}.  Thus Darboux
    naturality is the vanishing of the class of this total cocycle in
    \(H^\bullet(\mathcal C^\bullet_{\mathrm{Darb}},D_{\mathrm{Darb}})\),
    with the displayed primitives chosen.

  \item \emph{Unreduced cotangent-current realization.}
    For every brane interval \(I\), \(\kappa_I^{\mathrm{cot}}\) is the
    continuous filtered Maurer--Cartan morphism of
    Definition~\ref{defn:typed-cotangent-current-lift-datum}.  Its
    source is the interval shifted-cotangent current Lie algebra
    \[
      \mathfrak g_I^{\mathrm{cur}}
      =
      \bigl(\Omega^\bullet_c(I)\widehat\otimes\mathfrak h\bigr)
      \ltimes
      \bigl(\mathcal C^{\mathrm{dR}}_{c,\mathrm{cur}}(I)
      \widehat\otimes P\bigr)[1],
      \qquad P=\mathfrak h^\vee_{\mathrm{cont}},
    \]
    and its target is the unreduced principal-part boundary algebra
    \[
      \widetilde{\mathrm{Obs}}^{\mathrm{cl,pp}}_{\partial,N}(I)
      =
      \mathrm{Obs}^{\mathrm{cl}}_{\mathrm{open},N}(I)
      \widehat\otimes
      \widehat{\mathbf S}\!\left(
        \mathcal C^{\mathrm{dR}}_{c,\mathrm{cur}}(I)
        \widehat\otimes P[1]\right).
    \]
    The reduced image is fixed on generators by
    \[
      u_{a(t)dt\otimes f}\longmapsto B_f(a),
      \qquad
      c_{B,\rho}\longmapsto\Theta_\rho(B),
    \]
    where \(B_f(a)\) is the smeared Hamiltonian boundary observable and
    \(\Theta_\rho(B)\) is the odd principal-part current.  The reduced
    bracket equations are
    \[
      \{B_f(a),\Theta_\rho(B)\}
      =
      \Theta_{f\cdot\rho}(aB),
      \qquad
      \{\Theta_\rho(B),\Theta_\sigma(C)\}=0,
    \]
    equivalently, for \(f=z_1^pz_2^q\) and
    \(\rho_{i,j}=z_1^{-i-1}z_2^{-j-1}dz_1dz_2\),
    \[
      \{B_{z_1^pz_2^q}(a),\Theta_{\rho_{i,j}}(B)\}
      =
      -(pj-qi+p-q)
      \Theta_{\rho_{i-p+1,j-q+1}}(aB),
    \]
    with negative-index principal parts set to zero.  The unreduced
    distribution-product convention is the same as in
    Definition~\ref{defn:typed-cotangent-current-lift-datum}: \(aB\)
    means \((aB)(\varphi)=B(a\varphi)\), no product \(BC\) is formed, and
    every graph-level use of point-mass currents is routed through the
    wavefront-admissible or renormalized Costello graph complexes.  The
    equal-time \(\delta(t-t')\) is only the diagonal contact distribution
    whose smooth smearing gives \(ab\).
    The unreduced
    lift contains a primitive projection
    \[
      \pi_{\mathrm{prim}}\colon
      \mathrm{Obs}^{\mathrm{cl}}_{\mathrm{open},N}(I)
      \longrightarrow
      \Omega^\bullet_c(I)\widehat\otimes \bar A
    \]
	    and centrality homotopies against every unreduced open observable
	    \(A\):
	    \[
	      d_{\mathrm{cot}}H^{\mathrm{mix}}_{f,\rho,a,B}
      =
      -\bigl(\{B_f(a),\Theta_\rho(B)\}_{P_0}
        -\Theta_{f\cdot\rho}(aB)\bigr),
      \qquad
      d_{\mathrm{cot}}H^{\mathrm{open}}_{\rho,B;A}
	      =
	      -\bigl(\{\Theta_\rho(B),A\}_{P_0}
	        -\operatorname{act}^{P_0}_{\rho,B}(A)\bigr).
	    \]
	    Equivalently,
	    \(d_{\mathrm{cot}}H^{\mathrm{open}}_{\rho,B;A}
	    =-\Delta^{\mathrm{open}}_{\widetilde\kappa_I}(c_{B,\rho};A)\)
	    in
	    \(\mathcal H^\bullet_{\mathrm{cot,open}}(I)\).  Here
	    \(\operatorname{act}^{P_0}_{\rho,B}\) is the chosen \(P_0\)-action of
	    the cotangent current \(c_{B,\rho}\) on unreduced open observables in
	    the source current complex.
    Thus \(P=\mathfrak h^\vee_{\mathrm{cont}}\) is realized by boundary
    currents, not by the trace map \(J\).

    The analytic Bochner--Martinelli--Koppelman representative of this
    row carries the finite-window obstruction vector
    \[
      \operatorname{Ob}^{\mathrm{BMK\mbox{-}curr}}_M
      =
      \bigl(\operatorname{ob}_{\mathrm{diag}},
      \operatorname{ob}_{\mathrm{an}/\mathrm{fin},M},
      \operatorname{ob}_{\mathfrak h,M},
      \operatorname{ob}_{\mathrm{collar\mbox{-}fact},M}\bigr)
    \]
    of Theorem~\ref{thm:app-finite-window-bmk-current-transfer}.  The
    Hamiltonian-window component is nonzero for the strict finite
    support-local quotient because
    \(z_1^{M+2}\cdot\rho_{M+1,0}=-(M+2)\rho_{0,1}\).  The D5 datum
    therefore consists of either a support-local quotient killing this
    vector or the one-pair analytic pro-Matlis replacement
    \(\widehat{\mathcal P}_{\Pi}\) with compatible projections
    \(\pi_M P_{\Pi}^{\mathrm{an}}=p_M\).  The obstruction class
    \[
      \operatorname{Ob}^{\mathrm{cot}}_\partial(I)
      \in H^2(\mathcal K^{\mathrm{cot}}_\partial(I))
    \]
    is represented by
    \(\mathcal R^{\mathrm{cot}}_{\partial,I}\) in the relative homotopy
    fibre
    \(\operatorname{hofib}_{\Theta_I^{\mathrm{red}}}
    (\rho_{\mathrm{cot},I})\) and is killed by a degree-one primitive
    \[
      d_{\mathrm{cot}}h_{\mathrm{cot},I}
      =
      -\mathcal R^{\mathrm{cot}}_{\partial,I},
    \]
    compatible with interval inclusions and factorization products.

  \item \emph{Brane-link descent.}
    \(\mathfrak F_{\mathrm{desc}}\) is the total degree-zero
    \v{C}ech--Weiss cocycle of
    Definition~\ref{defn:typed-brane-link-descent-cone-datum},
    \[
      d_{\mathrm{desc}}\mathfrak F_{\mathrm{desc}}=0,
    \]
    whose local components are maps
    \(\mathrm{OCA}_{N,a}\colon
    \mathcal A|_{U_a}\to Z_\partial|_{U_a}\).  Equivalently, each next
    descent curvature represents a class
    \[
      \operatorname{Ob}^{\mathrm{desc}}_{k+1}
      =
      [\mathcal R^{\mathrm{desc}}_{k+1}]
      \in
      H^1\!\left(
      \operatorname{gr}^{k+1}_{F_{\check C}}
      \operatorname{Tot}\check C^\bullet_{\mathrm{Weiss}}
      (\mathfrak U;\mathcal E^\bullet_{\mathrm{desc}})
      \right)
    \]
    and is killed by a Čech-degree-\((k+1)\) primitive
    \(d_{\mathrm{desc}}H^{(k+1)}
      =-\mathcal R^{\mathrm{desc}}_{k+1}\).

  \item \emph{All-order Costello QME tower.}
    The data
    \(I^w_{n,M}[L]\), \(C_{n,w,r,M,L}\),
    \(K^{\mathrm{RG}}_{n,w,r,M;L_1,L_2}\), and
    \(K^{\mathrm{win}}_{n,w,r;M,M',L}\) are those of
    Definition~\ref{defn:typed-brane-defect-qme-tower}.  They satisfy
    \[
      d^{L}_{\mathrm{QME},n,M}C_{n,w,r,M,L}
      =
      -\mathcal R^{\mathrm{QME}}_{n,w,r,M,L},
    \]
    together with the scale and window equations
    \[
      d^{L_2}_{\mathrm{QME},n,M}K^{\mathrm{RG}}
      =-\bigl(W(P_{L_1,L_2},C_{n,w,r,M,L_1})
        -C_{n,w,r,M,L_2}\bigr),
    \]
    \[
      d^{L}_{\mathrm{QME},n,M}K^{\mathrm{win}}
      =-\bigl(\pi_{M,M'}C_{n,w,r,M,L}
        -C_{n,w,r,M',L}\pi^{\mathrm{CE}}_{M,M'}\bigr).
    \]

  \item \emph{Anomaly and locality.}
    The anomaly and locality primitives are those of
    Definition~\ref{defn:typed-anomaly-locality-datum}.  The anomaly
    complex is the completed operadic convolution \(L_\infty\)-algebra
    \[
      \mathcal K^{\mathrm{anom}}
      =
      \mathrm{Conv}_{\mathrm{SC}}\!\left(
      \mathcal A^{\mathrm{cl}/\hbar}_{\mathrm{bulk}},
      Z_\partial\right),
    \]
    with differential \(d_{\mathrm{Conv}}\), Swiss-cheese brackets, and
    the diagonal/contact support filtration inherited from the
    Costello local-functional and factorization structures.  The
    curvature of \(\mathrm{OCA}_N\) is
    \[
      \mathcal R^{\mathrm{anom}}
      =
      d_{\mathrm{Conv}}\mathrm{OCA}_N
      +\frac12[\mathrm{OCA}_N,\mathrm{OCA}_N]_{\mathrm{SC}}
      +\sum_{k\ge3}\frac1{k!}\ell_k(\mathrm{OCA}_N^{\otimes k})
      \in(\mathcal K^{\mathrm{anom}})^2 .
    \]
    The scalar-contact projection
    \[
      \sigma_{\mathrm{Cap}}\colon
      (\mathcal K^{\mathrm{anom}})^2
      \longrightarrow
      \C[[\hbar_{\mathrm{W}}]]\,\gamma_{\mathbf1}
    \]
    is the projection to the central identity-contact component of the
    trace representation.  Strict anomaly cancellation means
    \[
      \sigma_{\mathrm{Cap}}(\mathcal R^{\mathrm{anom}})=0,\qquad
      d_{\mathrm{Conv}}h^{\mathrm{str}}_{\mathrm{anom}}
      =
      -\mathcal R^{\mathrm{anom}}.
    \]
    The projective anomaly datum instead fixes the target central
    extension obtained by taking the raw Rees/Weyl commutator of the
    canonical unscaled constant-Hamiltonian extension and pushing out
    along \(\chi_N(K)=N\):
    \[
      0\to \C[[\hbar_{\mathrm{W}}]]\,\gamma_{\mathbf1}
      \to Z_\partial^{\mathrm{proj}}
      \to Z_\partial\to0,
      \qquad
      \sigma_{\mathrm{Cap}}(\mathcal R^{\mathrm{anom}})
      =
      \hbar_{\mathrm{W}}\chi_N(K)\bar c\,\gamma_{\mathbf1}
      =\hbar_{\mathrm{W}}N\bar c\,\gamma_{\mathbf1},
    \]
    and trivializes only the reduced curvature
    \[
      (\mathcal R^{\mathrm{anom}})_{\mathrm{red}}
      =
      \mathcal R^{\mathrm{anom}}
      -\hbar_{\mathrm{W}}N\bar c\,\gamma_{\mathbf1},
      \qquad
      d_{\mathrm{Conv}}h_{\mathrm{anom}}
      =
      -(\mathcal R^{\mathrm{anom}})_{\mathrm{red}}.
    \]
    Thus the raw Capelli trace scalar is retained in the projective case
    as the scalar commutator defect of \(Z_\partial^{\mathrm{proj}}\).
    The source extension still has cocycle \(\bar cK\).

    Locality is a separate factorization condition.  For
    \(\underline U=(U_1,\ldots,U_m)\in
    \operatorname{Ran}_{\mathrm{disj}}(\partial X)\), put
    \[
      \mathcal E^\bullet_{\mathrm{loc}}(\underline U)
      =
      \operatorname{Cone}\!\left(
      \mathcal E^\bullet_{\mathrm{desc}}(\coprod_iU_i)
      \longrightarrow
      \widehat\bigotimes_i
      \mathcal E^\bullet_{\mathrm{desc}}(U_i)\right)[-1],
    \]
    and
    \[
      \mathcal K^{\mathrm{loc}}
      =
      \operatorname{Tot}\check C^\bullet_{\mathrm{Weiss,Ran}}
      (\mathfrak U;\mathcal E^\bullet_{\mathrm{loc}}).
    \]
    The locality curvature and primitive are
    \[
      \mathcal R^{\mathrm{loc}}
      =
      \partial_{\mathrm{loc}}\mathrm{OCA}_N\in
      (\mathcal K^{\mathrm{loc}})^1,\qquad
      \partial_{\mathrm{loc}}h_{\mathrm{loc}}
      =
      -\mathcal R^{\mathrm{loc}}.
    \]
    The primitive \(h_{\mathrm{loc}}\) is required to be continuous for
    Weiss restriction, symmetric in the disjoint opens \(U_i\), and
    compatible with Ran refinement; secondary and higher refinement
    defects are killed by the global primitive-coherence row
    \textup{(D9)}.  The unreduced scalar projection of
    \(\mathcal R^{\mathrm{anom}}\) is the raw Rees/Weyl trace class
    \(\hbar_{\mathrm{W}}\chi_N(K)\bar c\,\gamma_{\mathbf1}
    =\hbar_{\mathrm{W}}N\bar c\,\gamma_{\mathbf1}\).  A
    determinant-line or modular-line curvature statement requires a
    supplied line, connection, Atiyah-class pullback, and
    \(\operatorname{Ob}^{\det}_J\)-killing lift.

  \item \emph{Common deformation complex and secondary compatibility.}
    The comparison-unifying datum is the common filtered
    \(L_\infty\)-algebra
    \[
      \mathfrak K_{\mathrm{HT}}
      =
      \operatorname{Conv}_{\mathrm{SC}}\!\left(
      \mathcal A^{\mathrm{cl}/\hbar}_{\mathrm{bulk}},
      Z^{\mathrm{der},\mathrm{Mor}}_{E_2^{\mathrm{HT}}}
      (\mathcal C^{\mathrm{op}}_\partial)\right)
    \]
    of Definition~\ref{defn:common-ht-deformation-complex}, together
    with the filtered comparison morphisms
    \[
      \chi_{\mathrm{nat}}
      =
      (\chi_{\mathrm{rad}},\chi_{\mathrm{Lie}},\chi_{\mathrm{QME}},
       \chi_{\mathrm{cent}},\chi_{\mathrm{BM}},\chi_{\mathrm{glob}})
    \]
    from the six native complexes of
    \(\S\ref{ssec:master-deformation-local-datum}\) to
    \(\mathfrak K_{\mathrm{HT}}\).  No inverse or quasi-inverse is part of
    this datum.  A reverse comparison is a separate filtered
    \(L_\infty\)-morphism
    \(\psi_i\colon\mathfrak K_{\mathrm{HT}}\to\mathfrak K_i\) together
    with its own comparison cone and homotopies; it is not inferred from
    \(\chi_i\).  For a native complex
    \((\mathfrak K_i,d_i,\ell^i_\bullet)\), the corresponding component
    \(\chi_i=(\chi_{i,k})_{k\ge1}\) is a filtered \(L_\infty\)-morphism:
    for every \(k\ge1\),
    \[
    \begin{aligned}
      \mathcal R_{\chi_i,k}(x_1,\ldots,x_k)
      &=
      \sum_{r\ge1}\frac1{r!}
      \sum_{I_1\sqcup\cdots\sqcup I_r=\{1,\ldots,k\}}
      \epsilon(I_\bullet)\,
      \ell^{\mathrm{HT}}_{r}\!\left(
        \chi_{i,|I_1|}(x_{I_1}),\ldots,
        \chi_{i,|I_r|}(x_{I_r})\right)  \\
      &\quad
      -
      \sum_{j=1}^{k}
      \sum_{\sigma\in\operatorname{Sh}(j,k-j)}
      \epsilon(\sigma)\,
      \chi_{i,k-j+1}\!\left(
        \ell^i_j(x_{\sigma(1)},\ldots,x_{\sigma(j)}),
        x_{\sigma(j+1)},\ldots,x_{\sigma(k)}\right)
      =
      0 .
    \end{aligned}
    \]
    Here \(I_\bullet\) ranges over ordered set partitions into nonempty
    blocks, \(x_{I_a}\) denotes the subsequence indexed by \(I_a\) in the
    inherited order, and \(\epsilon(I_\bullet),\epsilon(\sigma)\) are the
    Koszul signs fixed by Appendix~\ref{app:sign-conventions}.  The
    comparison-map deformation complex is
    \(\mathcal C^\bullet_{\chi_i}\) of
    Definition~\ref{defn:comparison-map-deformation-complex}.  By
    Lemma~\ref{lem:comparison-map-defect-cocycle}, after all lower
    filtration comparison equations are solved, the next defect is a
    cocycle.  Its obstruction is the class
    \[
      \operatorname{Ob}^{\mathrm{unif}}_{\mathrm{HT}}
      =
      \left\{[\mathcal R_{\chi,m+1}]\right\}_{m\ge0}
      \in
      \prod_{m\ge0}
      H^1\!\left(
        \operatorname{gr}^{m+1}\mathcal C^\bullet_\chi,
        D^{(m)}_\chi\right),
    \]
    where \(\mathcal R_{\chi,m+1}\) is the order-\((m+1)\) collection
    of the six comparison defects.  When
    \(\operatorname{Ob}^{\mathrm{unif}}_{\mathrm{HT}}=0\), the transported
    element \(\Theta_{\mathrm{OCA}}\in F^1\mathfrak K_{\mathrm{HT}}^1\)
    has curvature
    \[
      F_{\mathrm{HT}}(\Theta_{\mathrm{OCA}})
      =
      d_{\mathrm{HT}}\Theta_{\mathrm{OCA}}
      +\frac12\ell_2(\Theta_{\mathrm{OCA}},\Theta_{\mathrm{OCA}})
      +\sum_{k\ge3}\frac1{k!}
       \ell_k(\Theta_{\mathrm{OCA}}^{\otimes k}),
    \]
    and Proposition~\ref{prop:six-projections} identifies the radial,
    scalar Lie, QME, centrality/descent, Bochner--Martinelli--Koppelman,
    and Hamiltonian-period projections of this one curvature with the six
    native curvature representatives.  Without
    \((\mathfrak K_{\mathrm{HT}},\chi_{\mathrm{nat}})\), the obstruction
    atlas remains a tuple of native cohomology classes and is not a single
    curvature.

    Let \(\mathsf I_{\mathrm{glob}}\) be the transition category
    generated by arity truncations, Weiss-cover refinements, Costello
    scale changes, finite-window projections, admissible-weight
    transports, Darboux-chart changes, and brane-interval inclusions.
    For an object \(\xi\in\mathsf I_{\mathrm{glob}}\), let
    \(\mathcal P^\bullet(\xi)\) be the native primitive complex for the
    row of the datum over \(\xi\), with differential \(d_\xi\) equal to
    \(d_{\Theta_{\le n}},d_{\mathrm{Mor}},d_{\mathrm{desc}},
    d^L_{\mathrm{QME},n,M},d_{\mathrm{Conv}},\partial_{\mathrm{loc}},
    d_{\mathrm{Darb}}\), or \(d_{\mathrm{cot}}\) on the corresponding
    factor.  The primitive-compatibility complex is the total nerve/Roos
    complex
    \[
      \mathcal C^\bullet_{\mathrm{prim,glob}}
      =
      \operatorname{Tot}\,
      C^\bullet\!\left(
        N_\bullet\mathsf I_{\mathrm{glob}},\mathcal P^\bullet\right),
      \qquad
      D_{\mathrm{prim}}=d_\xi+(-1)^{|\cdot|}\delta_{\mathsf I}.
    \]
    The datum \(\mathfrak C_{\mathrm{prim}}\) is a cochain
    \(P_{\mathrm{glob}}\in\mathcal C^0_{\mathrm{prim,glob}}\) satisfying
    \[
      D_{\mathrm{prim}}P_{\mathrm{glob}}=-\mathcal R_{\mathrm{glob}},
    \]
    where \(\mathcal R_{\mathrm{glob}}\) has components
    \[
      \mathcal R^{\mathrm{SC}},\quad
      \mathcal R^{\mathrm{ch}\,E_2}_{\mathrm{Deligne}},\quad
      \mathcal R^{\mathrm{desc}},\quad
      \mathcal R^{\mathrm{QME}},\quad
      (\mathcal R^{\mathrm{anom}})_{\mathrm{red}},\quad
      \mathcal R^{\mathrm{loc}},\quad
      \Omega_\nabla,\quad
      \mathcal R_{\nabla,\Theta},\quad
      \mathcal R^{\mathrm{cot}} .
    \]
    Its \(0\)-simplex components are the comparison maps above and the
    primary primitives of \textup{(D1)}--\textup{(D8)}.  Its \(1\)- and
    \(2\)-simplex
    components include
    \[
      (K^{\mathrm{SC}},K^{\mathrm{cov}}_{\mathrm{desc}},
       K^{\mathrm{cov}}_{\mathrm{loc}},K^{\mathrm{RG}},K^{\mathrm{win}},
       K^{\mathrm{wt}},K^{\mathrm{Darb}},K^{\mathrm{cot}};\,
       L^{\mathrm{SC}},L^{\mathrm{cov}}_{\mathrm{desc}},
       L^{\mathrm{cov}}_{\mathrm{loc}},L^{\mathrm{RG}},L^{\mathrm{win}},
       L^{\mathrm{wt}},L^{\mathrm{Darb}},L^{\mathrm{cot}})
    \]
    with the secondary equations
    \[
    \begin{gathered}
      dK^{\mathrm{SC}}_{m,n}
      =-\bigl(\mathrm{res}_{m,n}H^{\mathrm{SC}}_{m+1}
        -H^{\mathrm{SC}}_{n+1}\bigr),\\
      d_{\mathrm{desc}}K^{\mathrm{cov}}_{r,\mathrm{desc}}
      =-\bigl(r^*h_{\mathrm{desc},\mathfrak U}
        -h_{\mathrm{desc},\mathfrak V}\bigr),\qquad
      \partial_{\mathrm{loc}}K^{\mathrm{cov}}_{r,\mathrm{loc}}
      =-\bigl(r^*h_{\mathrm{loc},\mathfrak U}
        -h_{\mathrm{loc},\mathfrak V}\bigr),\\
      d^{L_2}_{\mathrm{QME},n,M}K^{\mathrm{RG}}_{n,w,r,M;L_1,L_2}
      =-\bigl(W(P_{L_1,L_2},C_{n,w,r,M,L_1})
        -C_{n,w,r,M,L_2}\bigr),\\
      d^L_{\mathrm{QME},n,M}K^{\mathrm{win}}_{n,w,r;M,M',L}
      =-\bigl(\pi_{M,M'}C_{n,w,r,M,L}
        -C_{n,w,r,M',L}\pi^{\mathrm{CE}}_{M,M'}\bigr),\\
      d^L_{\mathrm{QME},n,M}K^{\mathrm{wt}}_{n,w,w';r,M,L}
      =-\bigl(R^M_{w,w'}C_{n,w,r,M,L}-C_{n,w',r,M,L}\bigr),\\
      d_{\mathrm{Darb}}K^{\mathrm{Darb}}_{\alpha\beta}
      =-\bigl(g_{\alpha\beta}^*h_{\nabla,\alpha}
        -h_{\nabla,\beta}\bigr),\qquad
      d_{\mathrm{cot}}K^{\mathrm{cot}}_{J,I}
      =-\bigl(\rho_{J,I}h_{\mathrm{cot},J}-h_{\mathrm{cot},I}\bigr).
    \end{gathered}
    \]
    The tertiary equations are the corresponding alternating-face
    equations; for example
    \[
    \begin{gathered}
      dL^{\mathrm{SC}}_{\ell,m,n}
      =
      -\bigl(\mathrm{res}_{m,n}K^{\mathrm{SC}}_{\ell,m}
        +K^{\mathrm{SC}}_{m,n}-K^{\mathrm{SC}}_{\ell,n}\bigr),\\
      d_{\mathrm{desc}}L^{\mathrm{cov}}_{s,r,\mathrm{desc}}
      =
      -\bigl(s^*K^{\mathrm{cov}}_{r,\mathrm{desc}}
        +K^{\mathrm{cov}}_{s,\mathrm{desc}}
        -K^{\mathrm{cov}}_{rs,\mathrm{desc}}\bigr),\\
      \partial_{\mathrm{loc}}L^{\mathrm{cov}}_{s,r,\mathrm{loc}}
      =
      -\bigl(s^*K^{\mathrm{cov}}_{r,\mathrm{loc}}
        +K^{\mathrm{cov}}_{s,\mathrm{loc}}
        -K^{\mathrm{cov}}_{rs,\mathrm{loc}}\bigr),\\
      d^{L_3}_{\mathrm{QME},n,M}L^{\mathrm{RG}}_{n,w,r,M;L_1,L_2,L_3}
      =
      -\bigl(W(P_{L_2,L_3},K^{\mathrm{RG}}_{n,w,r,M;L_1,L_2})
        +K^{\mathrm{RG}}_{n,w,r,M;L_2,L_3}
        -K^{\mathrm{RG}}_{n,w,r,M;L_1,L_3}\bigr),\\
      d_{\mathrm{Darb}}L^{\mathrm{Darb}}_{\alpha\beta\gamma}
      =
      -\bigl(g_{\beta\gamma}^*K^{\mathrm{Darb}}_{\alpha\beta}
        +K^{\mathrm{Darb}}_{\beta\gamma}
        -K^{\mathrm{Darb}}_{\alpha\gamma}\bigr),\\
      d_{\mathrm{cot}}L^{\mathrm{cot}}_{K,J,I}
      =
      -\bigl(\rho_{J,I}K^{\mathrm{cot}}_{K,J}
        +K^{\mathrm{cot}}_{J,I}
        -K^{\mathrm{cot}}_{K,I}\bigr).
    \end{gathered}
    \]
    For every \(p\ge3\), the \(p\)-simplex component of
    \(P_{\mathrm{glob}}\) kills the alternating face defect of the
    \((p-1)\)-simplex components.  The obstruction to extending the
    displayed low skeleton to all simplices is the corresponding class
    in \(H^1(\mathcal C^\bullet_{\mathrm{prim,glob}})\).

  \item \emph{Comparison cone.}
    After descent constructs \(\mathrm{OCA}_N\), the cone
    \(C_{\mathrm{OCA}}=\operatorname{Cone}(\mathrm{OCA}_N)\) carries a
    degree-\(-1\) contracting homotopy
    \[
      d_{\operatorname{Cone}}\kappa_{\operatorname{Cone}}
      +\kappa_{\operatorname{Cone}}d_{\operatorname{Cone}}
      =
      \operatorname{id}_{C_{\mathrm{OCA}}},
    \]
    or the following filtered associated-graded replacement.  Choose a
    complete separated descending filtration \(F^\bullet C_{\mathrm{OCA}}\)
    by subcomplexes, identify
    \[
      \operatorname{gr}C_{\mathrm{OCA}}
      =
      \operatorname{Cone}\!\left(
      \operatorname{gr}\mathrm{OCA}_N\colon
      \operatorname{gr}\mathcal A^{\mathrm{cl}/\hbar}_{\mathrm{bulk}}
      \to\operatorname{gr}Z_\partial\right),
    \]
    and supply a contraction \(\kappa_0\) of this associated-graded cone.
    A continuous lift \(\widetilde\kappa_0\) and filtration-raising error
    \[
      E=\operatorname{id}-
      (d_{\operatorname{Cone}}\widetilde\kappa_0
      +\widetilde\kappa_0d_{\operatorname{Cone}}),
      \qquad E(F^r)\subset F^{r+1},
    \]
    determine the actual contraction
    \[
      \kappa_{\operatorname{Cone}}
      =
      \widetilde\kappa_0\sum_{n\ge0}E^n
    \]
    in the complete filtered topology by
    Lemma~\ref{lem:complete-filtered-contraction-lift}.  This is the
    quasi-isomorphism datum; the existence of \(\mathrm{OCA}_N\) alone is
    only a morphism.
\end{enumerate}
\end{defn}

\begin{defn}[Relative Swiss-cheese extension obstruction]
\label{defn:relative-swiss-cheese-extension-obstruction}
With the notation of
Definition~\ref{defn:typed-mixed-ht-deligne-equivalence-datum}, put
\[
  \rho_{\mathrm{SC}}
  =
  (\rho_{\mathrm{tr},\le2},\rho_{\mathsf Q})
  \colon
  F^1\mathfrak{Def}^{\mathrm{HT}}_{\mathrm{SC}}
  (\mathcal A^{\mathrm{cl}/\hbar}_{\mathrm{bulk}},
   \mathcal C^{\mathrm{op}}_\partial)
  \longrightarrow
  \mathfrak{Def}_{\mathrm{tr},\le2}
  \oplus
  \mathrm{Conv}\bigl(\mathrm B\mathsf Q,
    \operatorname{End}_{\mathcal A^{\mathrm{Ham\mbox{-}cur}}_{\mathrm{bulk}}}
  \bigr).
\]
The target carries the fixed Maurer--Cartan element
\((\Theta^{\mathrm{tr}}_{\le2},\Theta_{\mathsf Q}^{\mathrm{cur}})\).
For an \(n\)-truncated lift
\(\Theta_{\le n}\) whose two restrictions agree with this fixed target
modulo \(F^{n+1}\), the arity-\((n+1)\) curvature is
\[
  \mathcal R^{\mathrm{SC}}_{n+1}(\Theta_{\le n})
  =
  \left(
  d\Theta_{\le n}
  +\frac12[\Theta_{\le n},\Theta_{\le n}]
  +\sum_{\ell\ge3}\frac1{\ell!}
    \ell_\ell(\Theta_{\le n}^{\otimes\ell})
  \right)_{F^{n+1}/F^{n+2}} .
\]
The target faces are already Maurer--Cartan at this stage: the trace
face is fixed by \(\Theta^{\mathrm{tr}}_{\le2}\), and the closed
tree-current face is the constructed action
\(\Theta_{\mathsf Q}^{\mathrm{cur}}\).  Hence
\[
  \rho_{\mathrm{SC}}
  \bigl(\mathcal R^{\mathrm{SC}}_{n+1}(\Theta_{\le n})\bigr)=0
  \quad\text{in}\quad F^{n+1}/F^{n+2}.
\]
Thus the obstruction to extending the fixed trace and tree-current faces
to arity \(n+1\) is the relative class
\[
  \operatorname{Ob}^{\mathrm{SC}/(\mathrm{tr},\mathsf Q)}_{n+1}
  =
  [(\mathcal R^{\mathrm{SC}}_{n+1}(\Theta_{\le n}),0)]
  \in
  H^2\!\left(
    \operatorname{gr}^{n+1}_{F}
    \mathfrak{Def}^{\mathrm{HT}}_{\mathrm{SC,rel}}
  \right).
\]
\end{defn}

\begin{prop}[Relative criterion for the all-arity Swiss-cheese lift]
\label{prop:relative-swiss-cheese-extension-criterion}
Assume the arity filtration on
\(\mathfrak{Def}^{\mathrm{HT}}_{\mathrm{SC,rel}}\) is complete and
pronilpotent.  An \(n\)-truncated lift \(\Theta_{\le n}\) extending
\(\Theta^{\mathrm{tr}}_{\le2}\) and
\(\Theta_{\mathsf Q}^{\mathrm{cur}}\) extends to arity \(n+1\) if and
only if
\[
  \operatorname{Ob}^{\mathrm{SC}/(\mathrm{tr},\mathsf Q)}_{n+1}=0 .
\]
A primitive is a degree-one element
\[
  H^{\mathrm{SC}}_{n+1}
  \in
  \operatorname{gr}^{n+1}_{F}
  \mathfrak{Def}^{\mathrm{HT}}_{\mathrm{SC,rel}}
\]
with
\[
  d_{\Theta_{\le n}}H^{\mathrm{SC}}_{n+1}
  =
  -\mathcal R^{\mathrm{SC}}_{n+1}(\Theta_{\le n}).
\]
Strict extension of the fixed faces means
\(\rho_{\mathrm{SC}}(H^{\mathrm{SC}}_{n+1})=0\).  Extension up to
trace/tree-current gauge homotopy replaces this equality by a specified
degree-zero primitive of \(\rho_{\mathrm{SC}}(H^{\mathrm{SC}}_{n+1})\)
in the target restriction complex.  A compatible tower of such
primitives for all \(n\ge2\) is equivalent to a Maurer--Cartan point of
\(\mathfrak{Def}^{\mathrm{HT}}_{\mathrm{SC,rel}}\).
\end{prop}

\begin{proof}
The obstruction theory is the arity-filtration obstruction theory for a
complete \(L_\infty\)-algebra, applied to the homotopy fibre of
\(\rho_{\mathrm{SC}}\) over the fixed Maurer--Cartan element.  The
Bianchi identity for the \(L_\infty\) curvature gives
\[
  d_{\Theta_{\le n}}
  \mathcal R^{\mathrm{SC}}_{n+1}(\Theta_{\le n})=0
  \quad\text{in}\quad
  \operatorname{gr}^{n+1}_{F}.
\]
Changing the arity-\((n+1)\) component by
\(H^{\mathrm{SC}}_{n+1}\) changes the curvature by
\(d_{\Theta_{\le n}}H^{\mathrm{SC}}_{n+1}\).  Since the two target faces
are fixed Maurer--Cartan faces, the curvature lies in the relative cone.
Its relative cohomology class vanishes exactly when a primitive with the
displayed differential exists.  Completeness and pronilpotence identify
the inverse system of arity extensions with a Maurer--Cartan element in
the complete relative complex.
\end{proof}

\begin{defn}[Relative Morita \(E_2^{\mathrm{HT}}\)-extension obstruction]
\label{defn:relative-morita-e2-extension-obstruction}
Let
\[
  \mathfrak D_{\mathrm{Mor}}
  =
  \mathfrak{Def}^{\mathrm{HT}}_{\mathrm{Mor}\mbox{-}E_2}
  (\mathcal C^{\mathrm{op}}_\partial,b_p)
  =
  \operatorname{fib}(\operatorname{res}_{b_p})[-1]
\]
be the complete filtered \(L_\infty\)-algebra of
Definition~\ref{app:defn:morita-chiral-E2-extension-complex}, twisted
by the fixed tree-current Maurer--Cartan element
\(\Theta_{\mathsf Q}^{\mathrm{cur}}\).  For an \(r\)-truncated Morita
lift \(\Theta_{\mathrm{Mor},\le r}\), the order-\((r+1)\) curvature is
\[
  \mathcal R^{\mathrm{Mor}\mbox{-}E_2}_{r+1}
  (\Theta_{\mathrm{Mor},\le r})
  =
  \left(
  d\Theta_{\mathrm{Mor},\le r}
  +\frac12[\Theta_{\mathrm{Mor},\le r},
           \Theta_{\mathrm{Mor},\le r}]
  +\sum_{\ell\ge3}\frac1{\ell!}
    \ell_\ell(\Theta_{\mathrm{Mor},\le r}^{\otimes\ell})
  \right)_{F^{r+1}/F^{r+2}} .
\]
The four projections of this curvature are
\[
  \mathcal R^{\mathrm{Mor}\mbox{-}E_2}_{r+1}
  =
  (\mathcal R_F,\mathcal R_C,\mathcal R_R,\mathcal R_{\mathrm{Mor}})_{r+1},
\]
where the four entries are the formality, centrality, Ran/Weiss, and
Morita-fibre faces appearing in row \textup{(D2)} of
Definition~\ref{defn:typed-mixed-ht-deligne-equivalence-datum}.  The
relative Morita obstruction class is
\[
  \operatorname{Ob}^{\mathrm{Mor}\mbox{-}E_2}_{r+1}
  =
  \left[
  \mathcal R^{\mathrm{Mor}\mbox{-}E_2}_{r+1}
  (\Theta_{\mathrm{Mor},\le r})
  \right]
  \in
  H^2\!\left(
    \operatorname{gr}^{r+1}_{F}\mathfrak D_{\mathrm{Mor}}
  \right).
\]
This is the operadic Morita-centre lift obstruction.  It is distinct
from the categorical Morita-density quotient
\(\mathcal Q^{\mathrm{Mor}}_{b_p}\) and the Hochschild-chart cone
\(\operatorname{Cone}(\operatorname{ev}_{b_p})\) of
Definition~\ref{defn:morita-density-quotient-chart-cone}.
\end{defn}

\begin{prop}[Relative criterion for the Morita centre lift]
\label{prop:relative-morita-e2-extension-criterion}
Assume the filtration on \(\mathfrak D_{\mathrm{Mor}}\) is complete and
pronilpotent.  An \(r\)-truncated Morita mixed-HT
\(E_2^{\mathrm{HT}}\)-centre lift extending
\(\Theta_{\mathsf Q}^{\mathrm{cur}}\) extends to order \(r+1\) if and
only if
\[
  \operatorname{Ob}^{\mathrm{Mor}\mbox{-}E_2}_{r+1}=0 .
\]
A primitive is a degree-one element
\[
  H^{\mathrm{Mor}\mbox{-}E_2}_{r+1}
  =
  (h_F,h_C,h_R,h_{\mathrm{Mor}})_{r+1}
  \in
  \operatorname{gr}^{r+1}_{F}\mathfrak D_{\mathrm{Mor}}
\]
with
\[
  d_{\Theta_{\mathrm{Mor},\le r}}
  H^{\mathrm{Mor}\mbox{-}E_2}_{r+1}
  =
  -\mathcal R^{\mathrm{Mor}\mbox{-}E_2}_{r+1}
  (\Theta_{\mathrm{Mor},\le r}).
\]
Strict preservation of the fixed formal-Darboux trace/current chart
means
\[
  \operatorname{res}_{b_p}
  \bigl(H^{\mathrm{Mor}\mbox{-}E_2}_{r+1}\bigr)=0 .
\]
Preservation up to chart gauge homotopy replaces this equality by a
specified degree-zero primitive in the restriction complex.  A
compatible tower of such primitives for all \(r\) is equivalent to a
Maurer--Cartan lift in \(\mathfrak D_{\mathrm{Mor}}\).
\end{prop}

\begin{proof}
This is the obstruction theory of the filtered homotopy fibre of
\(\operatorname{res}_{b_p}\).  The \(L_\infty\) Bianchi identity gives
\[
  d_{\Theta_{\mathrm{Mor},\le r}}
  \mathcal R^{\mathrm{Mor}\mbox{-}E_2}_{r+1}
  (\Theta_{\mathrm{Mor},\le r})=0
  \quad\text{in}\quad
  \operatorname{gr}^{r+1}_{F}\mathfrak D_{\mathrm{Mor}}.
\]
Changing the order-\((r+1)\) component by
\(H^{\mathrm{Mor}\mbox{-}E_2}_{r+1}\) changes the curvature by
\(d_{\Theta_{\mathrm{Mor},\le r}}
H^{\mathrm{Mor}\mbox{-}E_2}_{r+1}\).  Hence the displayed cohomology
class vanishes exactly when the next correction exists.  The homotopy
fibre condition is precisely the restriction equation along
\(\operatorname{res}_{b_p}\).  Completeness and pronilpotence identify
the inverse limit of finite-stage corrections with a Maurer--Cartan
element of \(\mathfrak D_{\mathrm{Mor}}\).
\end{proof}

\begin{defn}[Relative Hochschild-centre extension obstruction]
\label{defn:relative-hochschild-centre-extension-obstruction}
Let
\[
  \mathcal K^{\mathrm{gen}}_{\partial,N}
  =
  \underline{\operatorname{Hom}}_{\mathrm{cont}}\!\left(
    \mathfrak g_{\mathrm{src}},\,
    C^\bullet_{\mathrm{Hoch},\mathrm{adm}}(A,A)[1]\right)
\]
be the first-Taylor-coefficient target of
Definition~\ref{defn:typed-hochschild-centre-lift-datum}, and let
\[
  r_{\mathrm{gen}}\colon
  \mathcal K^{\mathrm{cent}}_{\partial,N}
  \longrightarrow
  \mathcal K^{\mathrm{gen}}_{\partial,N}
\]
be restriction of a twisting cochain to the cogenerators of the
Chevalley--Eilenberg chain coalgebra, followed by desuspension.  Fix the
formal-Darboux first Taylor face \(\tau_{N,1}\), whose trace-sector
coordinate dual is
\[
  \Phi_N^{\mathrm{gen}}(c_f)=\theta_f,\qquad
  \Phi_N^{\mathrm{gen}}(u_f)=M_{J(f)} .
\]
The relative Hochschild-centre deformation complex is the complete
filtered \(L_\infty\)-algebra
\[
  \mathfrak D_{\mathrm{Hoch}/\mathrm{gen}}
  =
  \operatorname{fib}(r_{\mathrm{gen}})[-1],
\]
twisted by the fixed face \(\tau_{N,1}\).  Its elements are centrality
corrections whose first Taylor projection is zero.

For an \(m\)-truncated lift
\[
  \mathcal C_{\le m}
  =
  \tau_{N,1}
  +H^{\mathrm{Lie}}+H^{\mathrm{act}}
  +H^{\mathrm{prod}}+H^{P_0}+H^3+\cdots+H^m
\]
with \(r_{\mathrm{gen}}(\mathcal C_{\le m})=\tau_{N,1}\),
the order-\((m+1)\) relative curvature is
\[
  \mathcal R^{\mathrm{cent}/\mathrm{gen}}_{m+1}
  (\mathcal C_{\le m})
  =
  \left(
  \sum_{\ell\ge1}\frac1{\ell!}
  \ell_\ell(\mathcal C_{\le m}^{\otimes\ell})
  \right)_{F^{m+1}/F^{m+2}}
  \in
  \operatorname{gr}^{m+1}_{F}
  \mathfrak D_{\mathrm{Hoch}/\mathrm{gen}} .
\]
For \(m=2\), this representative has the primary and first Stasheff
faces
\[
  (\delta^{\mathrm{Lie}},\delta^{\mathrm{act}},
    \delta^{\mathrm{prod}},\delta^{P_0},\delta^3),
\]
with the signs and degrees of
Definition~\ref{defn:typed-hochschild-centre-lift-datum}.  The relative
Hochschild-centre obstruction class is
\[
  \operatorname{Ob}^{\mathrm{cent}/\mathrm{gen}}_{m+1}
  =
  \left[
  \mathcal R^{\mathrm{cent}/\mathrm{gen}}_{m+1}
  (\mathcal C_{\le m})
  \right]
  \in
  H^2\!\left(
    \operatorname{gr}^{m+1}_{F}
    \mathfrak D_{\mathrm{Hoch}/\mathrm{gen}}\right).
\]
The absolute tower
\(\operatorname{Ob}^{\mathrm{cent}}_{\mathrm{Hoch}}\) is the same
centrality curvature before imposing the fixed first Taylor face; the
relative class above is the obstruction to extending the constructed
first Taylor coefficient to a continuous Hochschild-centre twisting
cochain.
\end{defn}

\begin{prop}[Relative criterion for the Hochschild-centre lift]
\label{prop:relative-hochschild-centre-extension-criterion}
Assume the centrality filtration on
\(\mathfrak D_{\mathrm{Hoch}/\mathrm{gen}}\) is complete and
pronilpotent.  An \(m\)-truncated Hochschild-centre lift extending
\(\tau_{N,1}\) extends to order \(m+1\) if and only if
\[
  \operatorname{Ob}^{\mathrm{cent}/\mathrm{gen}}_{m+1}=0 .
\]
A primitive is a degree-one suspended element
\[
  H^{\mathrm{cent}}_{m+1}
  =
  (H^{\mathrm{Lie}},H^{\mathrm{act}},H^{\mathrm{prod}},
   H^{P_0},H^3,\ldots,H^{m+1})_{F^{m+1}}
  \in
  \operatorname{gr}^{m+1}_{F}
  \mathfrak D_{\mathrm{Hoch}/\mathrm{gen}}
\]
whose order-\((m+1)\) component satisfies
\[
  d_{\mathcal C_{\le m}}H^{\mathrm{cent}}_{m+1}
  =
  -\mathcal R^{\mathrm{cent}/\mathrm{gen}}_{m+1}
  (\mathcal C_{\le m}).
\]
Strict preservation of the formal-Darboux first Taylor coefficient means
\[
  r_{\mathrm{gen}}(H^{\mathrm{cent}}_{m+1})=0.
\]
Preservation up to first-Taylor gauge homotopy replaces this equality by a
specified degree-zero primitive in
\(\mathcal K^{\mathrm{gen}}_{\partial,N}\).  A compatible tower of such
primitives for all \(m\) is equivalent to a Maurer--Cartan element
\(\tau_N=\mathcal C_{\Phi_N^{\mathrm{Hoch}}}\) whose projection is
\(\tau_{N,1}\).
\end{prop}

\begin{proof}
This is the obstruction theory of the filtered homotopy fibre of
\(r_{\mathrm{gen}}\), based at \(\tau_{N,1}\).  The
\(L_\infty\) Bianchi identity gives
\[
  d_{\mathcal C_{\le m}}
  \mathcal R^{\mathrm{cent}/\mathrm{gen}}_{m+1}
  (\mathcal C_{\le m})=0
  \quad\text{in}\quad
  \operatorname{gr}^{m+1}_{F}
  \mathfrak D_{\mathrm{Hoch}/\mathrm{gen}} .
\]
Changing the order-\((m+1)\) component by
\(H^{\mathrm{cent}}_{m+1}\) changes the curvature by
\(d_{\mathcal C_{\le m}}H^{\mathrm{cent}}_{m+1}\).  Hence the displayed
cohomology class vanishes exactly when the next correction exists.  The
homotopy-fibre condition is precisely the restriction equation along
\(r_{\mathrm{gen}}\).  Completeness and pronilpotence identify the
inverse limit of finite centrality corrections with the Maurer--Cartan
element defining \(\Phi_N^{\mathrm{Hoch}}\).
\end{proof}

\begin{lem}[Acyclic dependency order for the typed Deligne datum]
\label{lem:typed-deligne-dependency-dag}
The rows of Definition~\ref{defn:typed-mixed-ht-deligne-equivalence-datum}
are used in the following partial order.  Put
\[
  T_0=\Theta^{\mathrm{tr}}_{\leq2}.
\]
Then
\[
\begin{gathered}
  T_0
  \longrightarrow
  (\mathrm{D1},\mathrm{D2},\mathrm{D3},\mathrm{D4},\mathrm{D5}),\\
  (\mathrm{D1},\mathrm{D2},\mathrm{D3},\mathrm{D4},\mathrm{D5})
  \longrightarrow \mathrm{D6},\qquad
  (\mathrm{D1},\mathrm{D2},\mathrm{D3},\mathrm{D5})
  \longrightarrow \mathrm{D7},\\
  (\mathrm{D6},\mathrm{D7})\longrightarrow\mathrm{D8},\qquad
  (\mathrm{D1},\ldots,\mathrm{D8})\longrightarrow\mathrm{D9},\qquad
  (\mathrm{D6},\mathrm{D8},\mathrm{D9})\longrightarrow\mathrm{D10}.
\end{gathered}
\]
The optional topologisation row is external to this order unless the
target statement includes the topological closed colour.  No proof of a
row in this order may use a primitive, null-homotopy, counterterm, or
cone contraction from a later row as input.
\end{lem}

\begin{proof}
The formal-Darboux trace fibre \(T_0\) fixes the first Taylor boundary
coefficient, its coordinate-dual values, and the binary diagonal kernel.
Theorem~\ref{app:thm:formal-polydisk-all-arity-current-operad}
extends the closed Hamiltonian-current part to the
\(\mathsf Q\)-restricted tree-kernel action
\(\Theta_{\mathsf Q}^{\mathrm{cur}}\).  Rows
\(\mathrm{D1}\)--\(\mathrm{D5}\) are the local habitats which can then be
posed over this fibre and this current subaction: the residual
two-coloured Swiss-cheese action, Morita centre action, Hochschild-centre
lift, Darboux torsor transport, and unreduced cotangent-current lift.  The
descent row \(\mathrm{D6}\) requires these local maps and local homotopies
before it can form the Cech--Weiss cocycle.  The Costello row
\(\mathrm{D7}\) requires the local action, centre, generator, and
cotangent-current data in the finite-scale local-functional complex, but
not the global cone contraction.  The anomaly and locality row
\(\mathrm{D8}\) is the curvature of the descended and quantum-corrected
open--closed comparison, so it follows \(\mathrm{D6}\) and
  \(\mathrm{D7}\).  Row \(\mathrm{D9}\) first places the native radial,
  scalar Lie, QME, centrality/descent, Bochner--Martinelli--Koppelman, and
  Hamiltonian-period complexes in the common deformation complex
  \(\mathfrak K_{\mathrm{HT}}\), then compares all primary primitives under
  arity, cover, scale, window, weight, Darboux, and interval transitions.
  The comparison cone \(\mathrm{D10}\) is formed
only after the descended comparison morphism and its compatible
curvature-killing data exist.  Every arrow therefore points from a datum
to a later construction using it, and the displayed relation is acyclic.
\end{proof}

\begin{prop}[Equivalence-datum criterion]
\label{prop:typed-mixed-ht-deligne-equivalence-datum}
The data of Definition~\ref{defn:typed-mixed-ht-deligne-equivalence-datum}
construct a mixed-HT open--closed morphism
\[
  \mathrm{OCA}_N\colon
  \mathcal A^{\mathrm{cl}/\hbar}_{\mathrm{bulk}}
  \longrightarrow
  Z^{\mathrm{der},\mathrm{Mor}}_{E_2^{\mathrm{HT}}}
  (\mathcal C^{\mathrm{op}}_\partial)
\]
whose formal-Darboux fibre has first Taylor coefficient
\(\tau_{N,1}^{\mathrm{tr}}\).  Its coordinate-dual expression is
\(\Phi_N^{\mathrm{gen}}(c_f)=\theta_f\) and
\(\Phi_N^{\mathrm{gen}}(u_f)=M_{J(f)}\).  The cone component of the
datum makes this morphism a quasi-isomorphism in
\(\mathcal B^{\mathrm{adm}}\); the operadic homotopy-inverse data in
\((\mathrm{D1})\)--\((\mathrm{D10})\) make it an equivalence of
mixed-HT Swiss-cheese algebras.  If any one of the ten displayed rows is
absent, the corresponding construction is undefined: the result is then
only a formal-Darboux trace-sector fibre, a local morphism, a projective
anomaly class, or a comparison map without a quasi-isomorphism proof.
\end{prop}

\begin{proof}
The dependency order is Lemma~\ref{lem:typed-deligne-dependency-dag}.
The constructed input inside rows \textup{(D1)}--\textup{(D2)} is the
\(\mathsf Q\)-restricted current action
\(\Theta_{\mathsf Q}^{\mathrm{cur}}\) of
Theorem~\ref{app:thm:formal-polydisk-all-arity-current-operad}.  Row
\textup{(D1)} supplies the residual mixed-HT Swiss-cheese source action
extending that subaction, and row \textup{(D2)} supplies its Morita centre
target lift along \(\operatorname{res}_{b_p}\).  Row
\textup{(D3)} extends the first Taylor coefficient to a continuous
Hochschild-centre twisting cochain.  Rows \textup{(D4)}--\textup{(D5)} make the construction
functorial over formal Darboux coordinates and include the full shifted
cotangent BF source.  Row \textup{(D6)} descends the local maps to
\(\partial X\).  Row \textup{(D7)} solves the Costello QME at every
  order, scale, finite window, and admissible weight.  Row \textup{(D8)}
  removes the non-scalar anomaly and locality curvatures, retaining the
  raw Rees/Weyl trace class
  \(\hbar_{\mathrm W}\chi_N(K)[\bar c]\) only in the projective target
  extension.  Row \textup{(D9)} supplies
  \((\mathfrak K_{\mathrm{HT}},\chi_{\mathrm{nat}})\), so the native
  obstruction classes become projections of one curvature, and makes the
  primary primitives compatible under all transition maps.  Row
\textup{(D10)} contracts the comparison cone.  The composition of these
constructions is exactly the typed open--closed morphism of
Definition~\ref{defn:oca-comparison}; the cone contraction is exactly the
criterion of Definition~\ref{defn:oca-comparison-cone} and
Theorem~\ref{thm:filtered-oca-cone-criterion}.
\end{proof}

\begin{thm}[Typed global mixed-HT Deligne criterion]
\label{thm:typed-global-mixed-ht-deligne}
In $\mathcal B^{\mathrm{adm}}$ on the brane-link boundary $\partial X$
(with possible refinement to $\mathcal C_\rho$ for topologisation), the
formal-Darboux trace theorem supplies the arity-\(\le2\)
trace-sector fibre of the pre-open--closed comparison datum
\(\mathrm{preOCA}_N\) of Definition~\ref{defn:oca-comparison}.  A typed
global mixed-HT Deligne datum adds the global open--closed morphism,
Morita centre lift, descent datum, all-order Costello quantum
master-equation primitives, and comparison-cone contraction.  It is the
equivalence datum of
Definition~\ref{defn:typed-mixed-ht-deligne-equivalence-datum}: a
supplied-data lift of \(\mathrm{preOCA}_N\) to an open--closed
comparison morphism
\[
  \mathrm{OCA}_N\colon\mathcal A^{\mathrm{cl}}_{\mathrm{bulk}}\to
  Z^{\mathrm{der},\mathrm{Mor}}_{E_2^{\mathrm{HT}}}
  (\mathcal C^{\mathrm{op}}_\partial)
\]
of supplied \(\mathrm{SC}^{\mathrm{HT}}_{(1,2)}\)-algebras, together
with the primitive choices below.  The all-arity action obstruction
tower of Theorem~\ref{thm:sc-ht-deligne-swiss-cheese} is sharpened here
to the relative classes
\(\operatorname{Ob}^{\mathrm{SC}/(\mathrm{tr},\mathsf Q)}_{n+1}\) of
Definition~\ref{defn:relative-swiss-cheese-extension-obstruction},
represented by
\(\mathcal R^{\mathrm{SC}}_{n+1}(\Theta_{\le n})\) and killed by
primitives \(H^{\mathrm{SC}}_{n+1}\) compatible with the fixed
trace-sector and tree-current faces; the Morita centre lift is controlled
by the four-face relative classes
  \(\operatorname{Ob}^{\mathrm{Mor}\mbox{-}E_2}_{r+1}\) of
Definition~\ref{defn:relative-morita-e2-extension-obstruction}, with
primitive
\((h_F,h_C,h_R,h_{\mathrm{Mor}})_{r+1}\) compatible with
\(\operatorname{res}_{b_p}\); the Hochschild-centre chart lift is
controlled by the relative centrality classes
\(\operatorname{Ob}^{\mathrm{cent}/\mathrm{gen}}_{m+1}\) of
Definition~\ref{defn:relative-hochschild-centre-extension-obstruction},
with centrality primitives preserving the fixed
\(\tau_{N,1}\) face; its coordinate-dual expression is
\(\Phi_N^{\mathrm{gen}}\).  The five native morphism-stage obstruction
classes of \eqref{eq:obs-five-component} are represented by
\(\mathcal R^{\mathrm{global},5}\) with compatible primitive equations
\eqref{eq:global-five-primitive-equations} and the secondary primitive
compatibility datum
\(\mathfrak C_{\mathrm{prim}}\) of
\eqref{eq:global-primitive-compatibility-datum} and
\eqref{eq:global-tertiary-primitive-compatibility-datum}, extended to the
all-simplex coherence datum
\eqref{eq:global-all-simplex-primitive-compatibility-datum}.  If the boundary datum is
required to vary over the formal Darboux coordinate torsor, the flat
transport obstruction pair
\[
  \operatorname{Ob}^{\mathrm{Darb}}
  =
  \bigl([\Omega_{\nabla}],
        [\mathcal R_{\nabla,\Theta}]\bigr)
\]
of Definition~\ref{defn:hol-bd-datum-HT} is an additional typed
boundary datum and must be killed by its own primitives.  If the datum is
required to extend the constructed trace-sector transport, the sharper
condition is the relative class
\(\operatorname{Ob}^{\mathrm{Darb}/\mathrm{tr}}\) of
Definition~\ref{defn:relative-darboux-extension-obstruction}, together
with primitives whose restrictions are the zero trace-sector primitives
or trace-exact gauge homotopies.  This full bar-centre extension is not
constructed by the formal-Darboux trace theorem.  If the target statement
uses the full shifted-cotangent BF source rather than only the stable
Hamiltonian trace sector, then the cotangent-current obstruction
\(\operatorname{Ob}^{\mathrm{cot}}_\partial(I)\) of
Definition~\ref{defn:typed-cotangent-current-lift-datum} is an
additional typed boundary datum on every brane interval \(I\), with
primitives \(h_{\mathrm{cot},I}\) compatible under interval extension
and factorization products.  The QME component
\(\operatorname{Ob}^{\mathrm{QME}}_{\mathrm{all\mbox{-}order}}\) is
the brane-defect tower of
Definition~\ref{defn:typed-brane-defect-qme-tower}: an order-indexed,
scale-indexed, and finite-window-indexed Costello problem, not a single
counterterm equation.  The descent component and the comparison-cone
obstruction are the two stages of
Definition~\ref{defn:typed-brane-link-descent-cone-datum}: first a
total \v{C}ech--Weiss cocycle descends the local maps, then a filtered
cone contraction supplies the quasi-isomorphism datum for the descended
map.  The
anomaly and locality components are the strict/projective curvature and
disjoint-union factorization data of
Definition~\ref{defn:typed-anomaly-locality-datum}; the raw Rees/Weyl
trace class \(\hbar_{\mathrm{W}}\chi_N(K)[\bar c]
=\hbar_{\mathrm{W}}N[\bar c]\) is retained as a projective anomaly class
unless a strict primitive is supplied.  The
open--closed comparison carries a quasi-isomorphism datum exactly when
  \(C_{\mathrm{OCA}}=\operatorname{Cone}(\mathrm{OCA}_N)\) is equipped with
  a contracting homotopy in \(\mathcal B^{\mathrm{adm}}\) in the sense of
  Definition~\ref{defn:oca-comparison-cone}, equivalently when the
  filtered cone datum
  \((F^\bullet C_{\mathrm{OCA}},\kappa_0,\widetilde\kappa_0,E)\) of
  Theorem~\ref{thm:filtered-oca-cone-criterion} has been supplied, with
\(E\) filtration-raising and \(\sum_{n\ge0}E^n\) convergent.  A
nonzero typed component, missing primitive, missing null-homotopy, or
non-acyclic comparison cone obstructs the typed OCA quasi-isomorphism.
If topologisation of the holomorphic closed colour is included, the
translation class \([\tau]\), the tower
\(\mathfrak o^{\mathrm{top}}_{n+1}\), the six hypotheses
\((\mathrm{T1})\)--\((\mathrm{T6})\), and the secondary coherence class
of Theorem~\ref{thm:topologisation-by-stress-tensor} are part of the
same obstruction datum.

The typed obstruction ledger is
\begin{equation}\label{eq:typed-total-obstruction-ledger}
  \operatorname{Ob}^{\mathrm{typed}}
  =
  \left(
    \{\mathfrak o^{\mathrm{SC}}_{n+1}\}_{n\ge2},
    \operatorname{Ob}^{\mathrm{cent}}_{\mathrm{Hoch}},
    \operatorname{Ob}^{\mathrm{global},5},
    \operatorname{Ob}^{\mathrm{Darb}},
    \operatorname{Ob}^{\mathrm{cot}}_\partial,
    \operatorname{Ob}^{\mathrm{unif}}_{\mathrm{HT}},
    \operatorname{Ob}^{\mathrm{prim}}_{\mathrm{glob}},
    H^\bullet\!\left(\operatorname{Cone}(\mathrm{OCA}_N)\right),
    \operatorname{Ob}^{\mathrm{top}}_\rho
  \right),
\end{equation}
where
\[
  \operatorname{Ob}^{\mathrm{prim}}_{\mathrm{glob}}
  =
  [d_{\mathsf I}P_{\mathrm{glob}}]
  \in H^1(\mathcal C^\bullet_{\mathrm{prim,glob}})
\]
and
\[
  \operatorname{Ob}^{\mathrm{unif}}_{\mathrm{HT}}
  =
  \left\{[\mathcal R_{\chi,m+1}]\right\}_{m\ge0}
  \in
  \prod_{m\ge0}
  H^1\!\left(
    \operatorname{gr}^{m+1}\mathcal C^\bullet_\chi,
    D^{(m)}_\chi\right)
\]
is the obstruction to the comparison maps from the six native complexes
being filtered \(L_\infty\)-morphisms to the common deformation complex of
Definition~\ref{defn:common-ht-deformation-complex}; the differential
\(D^{(m)}_\chi\) is the linearized \(L_\infty\)-morphism defect differential of
Definition~\ref{defn:comparison-map-deformation-complex}.  It is not an
inverse-comparison obstruction: a reverse map out of
\(\mathfrak K_{\mathrm{HT}}\) is separate filtered \(L_\infty\)-data
with its own cone.  The class
\(\operatorname{Ob}^{\mathrm{prim}}_{\mathrm{glob}}\) uses
\(d_{\mathsf I}\), the alternating-face differential on
\(\mathcal C^\bullet_{\mathrm{prim,glob}}\).  This class is the
obstruction to the all-simplex primitive-coherence equations in the
transition category \(\mathsf I_{\mathrm{glob}}\), and
\[
  \operatorname{Ob}^{\mathrm{top}}_\rho
  =
  \bigl([\tau],\{\mathfrak o^{\mathrm{top}}_{n+1}\}_{n\ge1}\bigr)
\]
is omitted unless the \(\mathcal C_\rho\)-topologisation is part of the
target statement.  The trace-sector formal-Darboux theorem supplies the
arity-\(\le2\) boundary value and the raw Rees/Weyl trace projection
\(\hbar_{\mathrm W}\chi_N(K)[\bar c]\).  The
representatives and primitives in
\eqref{eq:typed-total-obstruction-ledger} are the separate
morphism-stage lift data.

The native obstruction tuple for constructing the morphism-stage lift of
\(\mathrm{preOCA}_N\) is
\begin{equation}\label{eq:obs-five-component}
  \operatorname{Ob}^{\mathrm{global},5}
   \;:=\;
  \Bigl(
   \operatorname{Ob}^{\mathrm{ch}\,E_2}_{\mathrm{Deligne}}
   ,
   \operatorname{Ob}^{\mathrm{desc}}_{\partial}
   ,
   \operatorname{Ob}^{\mathrm{QME}}_{\mathrm{all\mbox{-}order}}
   ,
   \operatorname{Ob}^{\mathrm{anom}}
   ,
   \operatorname{Ob}^{\mathrm{loc}}
  \Bigr),
\end{equation}
where the first three components are those of
Theorem~\ref{thm:main-global-deligne-criterion} and live in their own
cohomology complexes, and the additional two are:
\begin{itemize}
  \item \(\operatorname{Ob}^{\mathrm{anom}}\in
        H^2\bigl(\mathrm{Conv}_{\mathrm{SC}}\!\bigl(
          \mathcal A^{\mathrm{cl}}_{\mathrm{bulk}},
          Z^{\mathrm{der},\mathrm{Mor}}_{E_2^{\mathrm{HT}}}\!\bigl(
            \mathcal C^{\mathrm{op}}_\partial\bigr)\bigr)\bigr)\),
        the anomaly-cancellation class in the operadic convolution complex of
        Definition~\ref{defn:convolution-complex}, whose trace-sector
        projective anomaly is the raw Rees/Weyl class
        \(\hbar_{\mathrm{W}}\chi_N(K)[\bar c]
        =\hbar_{\mathrm{W}}N[\bar c]\);
  \item \(\operatorname{Ob}^{\mathrm{loc}}=[\mathcal R^{\mathrm{loc}}]\in
        H^1\!\bigl(
        \operatorname{Tot}\check C^\bullet_{\mathrm{Weiss,Ran}}
        (\mathfrak U;\mathcal E^\bullet_{\mathrm{loc}})\bigr)\), the
        locality class of the disjoint-union extension of
        \(\mathrm{OCA}_N\) in the factorisation
        operad \(\mathrm{Fact}_{\mathrm{SC}^{\mathrm{HT}}_{(1,2)}}\).
        The complex \(\mathcal E^\bullet_{\mathrm{loc}}\) is the cone
        comparing the value on a Ran configuration with the tensor product
        over its pairwise disjoint components.
\end{itemize}
Choose representatives
\begin{equation}\label{eq:global-five-representatives}
  \mathcal R^{\mathrm{global},5}
  =
  \left(
  \mathcal R^{\mathrm{ch}\,E_2}_{\mathrm{Deligne}},
  \mathcal R^{\mathrm{desc}}_{\partial},
  \{\mathcal R^{\mathrm{QME}}_{n,w,r,M,L}\}_{n,w,r,M,L},
  \mathcal R^{\mathrm{anom}},
  \mathcal R^{\mathrm{loc}}
  \right)
\end{equation}
whose cohomology classes are the five entries of
\eqref{eq:obs-five-component}.  The trace-sector scalar projection of
\(\mathcal R^{\mathrm{anom}}\) is the raw Rees/Weyl pushout
\(\hbar_{\mathrm{W}}\chi_N(K)\bar c
=\hbar_{\mathrm{W}}N\bar c\).  The reduced anomaly
representative \((\mathcal R^{\mathrm{anom}})_{\mathrm{red}}\) is the
component left after this projective scalar line is recorded.
The differentials in the primitive equations are fixed as follows:
\(d_{\mathrm{Mor}}\) is the twisted differential in
\(\mathfrak{Def}^{\mathrm{HT}}_{\mathrm{Mor}\mbox{-}E_2}
(\mathcal C^{\mathrm{op}}_\partial,b_p)
=\operatorname{fib}(\operatorname{res}_{b_p})[-1]\);
\(d_{\mathrm{desc}}\) is the total differential
\(d_{\mathcal E}+(-1)^q\delta_{\check C}\) on
\(\operatorname{Tot}\check C^\bullet_{\mathrm{Weiss}}
(\mathfrak U;\mathcal E^\bullet_{\mathrm{desc}})\), where \(q\) is the
internal degree; \(d_{\mathrm{QME}}^{L}=Q+\{I[L],-\}_L+\hbar_{\mathrm{QME}}\Delta_L\);
\(d_{\mathrm{Conv}}\) is the operadic convolution differential of
Definition~\ref{defn:convolution-complex}; and
\(\partial_{\mathrm{loc}}\) is the total Costello--Gwilliam
\v{C}ech--Weiss--Ran locality differential on the disjoint-union
extension complex.  In the order-\(n\) QME primitive below,
\(d^{L}_{\mathrm{QME},n,M}\) denotes the
\(\hbar_{\mathrm{QME}}^n\)-coefficient differential of the finite-scale
linearised QME with lower orders fixed; at this coefficient it is
\(d^L_{\mathrm{cl},M}\) from
Definition~\ref{defn:typed-brane-defect-qme-tower}, not the completed
quantum differential attached to the final solution.
The compatible null-homotopies are the primitive equations
\begin{equation}\label{eq:global-five-primitive-equations}
\begin{gathered}
  d_{\mathrm{Mor}}h_{E_2}
  =-\mathcal R^{\mathrm{ch}\,E_2}_{\mathrm{Deligne}},\qquad
  d_{\mathrm{desc}}h_{\mathrm{desc}}
  =-\mathcal R^{\mathrm{desc}}_{\partial},\\
  d^{L}_{\mathrm{QME},n,M}C_{n,w,r,M,L}
  =-\mathcal R^{\mathrm{QME}}_{n,w,r,M,L},\qquad
  d_{\mathrm{Conv}}h_{\mathrm{anom}}
  =-(\mathcal R^{\mathrm{anom}})_{\mathrm{red}},\qquad
  \partial_{\mathrm{loc}}h_{\mathrm{loc}}
  =-\mathcal R^{\mathrm{loc}} .
\end{gathered}
\end{equation}
The QME row also contains Roos primitives
\(K^{\mathrm{RG}}_{n,w,r,M;L_1,L_2}\) and
\(K^{\mathrm{win}}_{n,w,r;M,M',L}\) satisfying
\[
  \begin{aligned}
  d^{L_2}_{\mathrm{QME},n,M}K^{\mathrm{RG}}_{n,w,r,M;L_1,L_2}
  &=
  -\bigl(W(P_{L_1,L_2},C_{n,w,r,M,L_1})-C_{n,w,r,M,L_2}\bigr),\\
  d^{L}_{\mathrm{QME},n,M}K^{\mathrm{win}}_{n,w,r;M,M',L}
  &=
  -\bigl(\pi_{M,M'}C_{n,w,r,M,L}
    -C_{n,w,r,M',L}\pi^{\mathrm{CE}}_{M,M'}\bigr),
  \end{aligned}
\]
and the analogous admissible-weight transport primitive.  The construction ledger is the
following.  Each row names an object in \(\mathcal B^{\mathrm{adm}}\),
the equation it must satisfy, and the primitive whose existence is the
vanishing statement.
\begin{center}
\scriptsize
\renewcommand{\arraystretch}{1.2}
\begin{tabular}{@{}p{0.10\textwidth}|p{0.25\textwidth}|p{0.33\textwidth}|p{0.20\textwidth}@{}}
\(\mathsf{Row}\) & \(\mathsf{Object}\) & \(\mathsf{Equation}\) &
\(\mathsf{Primitive}\)\\\hline
\(\mathrm{O0}\) &
Mixed-HT Swiss-cheese action on
\((S,\mathcal C^{\mathrm{op}}_\partial)\), with
\(S=Z_\partial\) or
\(S=\mathcal A^{\mathrm{cl}/\hbar}_{\mathrm{bulk}}\), extending the
constructed \(\mathsf Q\)-restricted current action
\(\Theta_{\mathsf Q}^{\mathrm{cur}}\) &
\(d\Theta_{\mathrm{SC}}+
\frac12[\Theta_{\mathrm{SC}},\Theta_{\mathrm{SC}}]=0\) in
\(\mathfrak{Def}^{\mathrm{HT}}_{\mathrm{SC}}
(S,\mathcal C^{\mathrm{op}}_\partial)\) &
\(H_{n+1}\) with
\(d_{\Theta_{\le n}}H_{n+1}
=-\mathcal R^{\mathrm{SC}}_{n+1}(\Theta_{\le n})\), whose residual class
is \(\mathfrak o^{\mathrm{SC}}_{n+1}\)\\
\(\mathrm{O1}\) &
Morita \(E_2^{\mathrm{HT}}\)-centre
\(Z_\partial\) with all chiral arities, relative to the restriction
\(\operatorname{res}_{b_p}\colon\mathsf P\to\mathsf Q\) &
Maurer--Cartan equation in the shifted homotopy fibre
\(\mathfrak{Def}^{\mathrm{HT}}_{\mathrm{Mor}\mbox{-}E_2}
(\mathcal C^{\mathrm{op}}_\partial,b_p)
=\operatorname{fib}(\operatorname{res}_{b_p})[-1]\) &
\(h_{E_2}\) with
\(d_{\mathrm{Mor}}h_{E_2}
=-\mathcal R^{\mathrm{ch}\,E_2}_{\mathrm{Deligne}}\), where
\([\mathcal R^{\mathrm{ch}\,E_2}_{\mathrm{Deligne}}]
=\operatorname{Ob}^{\mathrm{ch}\,E_2}_{\mathrm{Deligne}}\)\\
\(\mathrm{O}_{\mathrm{Hoch}}\) &
Admissible Hochschild-centre lift
\(\tau_N=\mathcal C_{\Phi_N^{\mathrm{Hoch}}}\) of the first Taylor
coefficient \(\tau_{N,1}\) &
\(d\mathcal C_{\Phi_N^{\mathrm{Hoch}}}
+\frac12[\mathcal C_{\Phi_N^{\mathrm{Hoch}}},
\mathcal C_{\Phi_N^{\mathrm{Hoch}}}]
+\frac1{3!}\ell_3(\mathcal C_{\Phi_N^{\mathrm{Hoch}}}^{\otimes3})
+\cdots=0\) in
\(\mathcal K^{\mathrm{cent}}_{\partial,N}\) &
\(H^{\mathrm{Lie}},H^{\mathrm{act}},H^{\mathrm{prod}},H^{P_0},H^3,\ldots\)
with \(d_{\mathrm{cent}}H=-\delta\), where
\([\delta]=\operatorname{Ob}^{\mathrm{cent}/\mathrm{gen}}_{m+1}\) in the
relative first-Taylor-fibre complex of
Definition~\ref{defn:relative-hochschild-centre-extension-obstruction}\\
\(\mathrm{O2}\) &
Weiss/Ran descent datum on a brane-link cover \(\mathfrak U\), as in
Definition~\ref{defn:typed-brane-link-descent-cone-datum} &
\(\operatorname{Ob}^{\mathrm{desc}}_{k+1}=0\) in
\(H^1(\operatorname{gr}^{k+1}_{F_{\check C}}
\operatorname{Tot}\check C^\bullet_{\mathrm{Weiss}}
(\mathfrak U;\mathcal E^\bullet_{\mathrm{desc}}))\) for all \(k\) &
\(h_{\mathrm{desc}}\) with
\(d_{\mathrm{desc}}h_{\mathrm{desc}}
=-\mathcal R^{\mathrm{desc}}_{\partial}\), where
\(\mathcal R^{\mathrm{desc}}_{\partial}\) represents
\(\operatorname{Ob}^{\mathrm{desc}}_{\partial}
=\{\operatorname{Ob}^{\mathrm{desc}}_{k+1}\}_{k\ge1}\)\\
\(\mathrm{O3}\) &
Costello local-functional quantization
\((Q^{\mathrm{GF}},K_t,P_{\varepsilon,L},\Delta_L,I[L],
W(P_{L_1,L_2},-),\pi_{M,M'})\) of
Definition~\ref{defn:typed-brane-defect-qme-tower} &
\(QI[L]+\frac12\{I[L],I[L]\}_{L}+\hbar_{\mathrm{QME}}\Delta_LI[L]=0\) in the
admissible local-functional complex &
counterterms \(C_{n,w,r,M,L}\) with
\(d^{L}_{\mathrm{QME},n,M}C_{n,w,r,M,L}
=-\mathcal R^{\mathrm{QME}}_{n,w,r,M,L}\), plus Roos primitives for scale
and window compatibility\\
\(\mathrm{O4}\) &
Operadic anomaly representative in
\(\mathrm{Conv}_{\mathrm{SC}}(\mathcal A^{\mathrm{cl}}_{\mathrm{bulk}},
Z_\partial)\), strict or projective as in
Definition~\ref{defn:typed-anomaly-locality-datum} &
\(\operatorname{Curv}_{\mathrm{anom}}(\mathrm{OCA}_N)=0\) modulo the
raw trace-sector scalar
\(\hbar_{\mathrm{W}}\chi_N(K)[\bar c]\) &
\(h_{\mathrm{anom}}\) with scalar projection fixed by the raw
Rees/Weyl trace pushout of the unscaled cocycle and
\(d_{\mathrm{Conv}}h_{\mathrm{anom}}
=-(\mathcal R^{\mathrm{anom}})_{\mathrm{red}}\)\\
\(\mathrm{O5}\) &
Factorization-locality extension of \(\mathrm{OCA}_N\) on disjoint
opens of \(\partial X\), as in
Definition~\ref{defn:typed-anomaly-locality-datum} &
\(\partial_{\mathrm{loc}}(\mathrm{OCA}_N)=0\) in the
Čech--Weiss--Ran locality complex &
\(h_{\mathrm{loc}}\) with
\(\partial_{\mathrm{loc}}h_{\mathrm{loc}}
=-\mathcal R^{\mathrm{loc}}\), where
\([\mathcal R^{\mathrm{loc}}]=\operatorname{Ob}^{\mathrm{loc}}\)\\
\(\mathrm{O}_{\nabla}\) &
Flat transport over the formal Darboux coordinate torsor for the bar
family and centre chart &
\(\operatorname{Ob}^{\mathrm{Darb}}
=([\Omega_{\nabla}],[\mathcal R_{\nabla,\Theta}])\) in
\(\mathcal E^\bullet_{\mathrm{Darb}}\) &
primitives \(h_{\nabla}\), \(h_{\nabla,\Theta}\) with
\(d_{\mathrm{Darb}}h_{\nabla}=-\Omega_{\nabla}\) and
\(d_{\mathrm{Darb}}h_{\nabla,\Theta}
=-\mathcal R_{\nabla,\Theta}\)\\
	\(\mathrm{O}_{\mathrm{cot}}\) &
	Unreduced cotangent current lift
	\(\kappa^{\mathrm{cot}}_I\) on
	\(\widetilde{\mathrm{Obs}}^{\mathrm{cl,pp}}_{\partial,N}(I)\) &
the reduced equations
\(\{B_f(a),\Theta_\rho(B)\}=\Theta_{f\cdot\rho}(aB)\) and
\(\{\Theta_\rho(B),\Theta_\sigma(C)\}=0\) lift through unreduced
	\(P_0\)-centrality against arbitrary open observables &
	\(h_{\mathrm{cot},I}\) with
	\(d_{\mathrm{cot}}h_{\mathrm{cot},I}
	=-\mathcal R^{\mathrm{cot}}_{\partial,I}\), where
	\([\mathcal R^{\mathrm{cot}}_{\partial,I}]
	=\operatorname{Ob}^{\mathrm{cot}}_\partial(I)\), and
	\(H^{\mathrm{open}}_{x,A}\) with
	\(d_{\mathrm{cot}}H^{\mathrm{open}}_{x,A}
	=-\Delta^{\mathrm{open}}_{\widetilde\kappa_I}(x;A)\) in
	\(\mathcal H^\bullet_{\mathrm{cot,open}}(I)\)\\
\(\mathrm{O6}\) &
Comparison cone in \(\mathcal B^{\mathrm{adm}}\), with filtered data
from Definition~\ref{defn:typed-brane-link-descent-cone-datum} &
\(d_{\operatorname{Cone}}\kappa_{\operatorname{Cone}}
+\kappa_{\operatorname{Cone}}d_{\operatorname{Cone}}
=\operatorname{id}\) &
\(\kappa_{\operatorname{Cone}}\), or equivalently the filtered
associated-graded contraction of
Theorem~\ref{thm:filtered-oca-cone-criterion}\\
\(\mathrm{O7}\) &
Optional topological transfer in \(\mathcal C_\rho\) &
\(D_\rho\Gamma=\tau\) and the Maurer--Cartan equation in
\(\mathfrak{Def}^{\rho}_{\mathrm{SC},\mathrm{top}}\) &
\(\Gamma\) with \(D_\rho\Gamma=\tau\), and topological arity primitives
with \(d_{\Theta_{\le n}}H^{\mathrm{top}}_{n+1}
=-\mathcal R^{\mathrm{top}}_{n+1}\)
\end{tabular}
\end{center}
The degree, support, filtration, and continuity data of these
primitives are part of the datum:
\begin{center}
\scriptsize
\renewcommand{\arraystretch}{1.2}
\begin{tabular}{@{}p{0.07\textwidth}|p{0.11\textwidth}|p{0.26\textwidth}|p{0.28\textwidth}|p{0.15\textwidth}@{}}
\(\mathsf{Row}\) & \(\mathsf{Degree}\) & \(\mathsf{Filtration/support}\) &
\(\mathsf{Continuity}\) & \(\mathsf{Status}\)\\\hline
\(\mathrm{O0}\) &
\(H_{n+1}\in(F^{n+1}/F^{n+2})^1\) &
arity filtration of
\(\mathfrak{Def}^{\mathrm{HT}}_{\mathrm{SC}}\); formal-Darboux
trace-sector boundary condition in arity \(\le2\) &
complete separated convolution topology; inverse-limit compatibility of
the arity primitives &
action criterion\\
\(\mathrm{O1}\) &
\(h_{E_2}\in C^1\) &
arity filtration of the cofibrant chiral \(E_2^{\mathrm{HT}}\)-operad;
formal Hamiltonian-current support at \(b_p\) &
continuous in \(\mathcal B^{\mathrm{adm}}\) and compatible with
\(\operatorname{res}_{b_p}\) &
criterion\\
\(\mathrm{O}_{\mathrm{Hoch}}\) &
\(H^{\mathrm{Lie}},H^{\mathrm{act}},H^{\mathrm{prod}},H^{P_0}\in
(\mathcal K^{\mathrm{cent}}_{\partial,N})^{-1}\),
\(H^3\in(\mathcal K^{\mathrm{cent}}_{\partial,N})^{-2}\), and
\(H^m\in(\mathcal K^{\mathrm{cent}}_{\partial,N})^{1-m}\) before
operadic suspension &
Hochschild centrality filtration, open-input filtration, finite
Hamiltonian windows, and \(P_0\)-bracket filtration &
continuous for the completed Hochschild brace operations, window
projections, and scalar/contact splitting &
typed boundary datum\\
\(\mathrm{O2}\) &
\(h_{\mathrm{desc}}\in\operatorname{Tot}^1\) &
Čech--Weiss degree on the collared brane-link cover \(\mathfrak U\);
support on finite intersections \(U_{a_0\cdots a_p}\) &
commutes with restrictions, refinements, and the completed Hom topology &
criterion\\
\(\mathrm{O3}\) &
\(C_{n,w,r,M,L}\in\mathcal Q^0_{w,r,M,L}(I)\) &
\(\hbar\)-adic, scalar-contact, Costello scale, and finite-window
filtrations; local functionals supported at the brane defect &
compatible with renormalisation-group flow, window projections, and
\(\varprojlim^1H^0\) primitives &
criterion\\
\(\mathrm{O4}\) &
\(h_{\mathrm{anom}}\in\mathrm{Conv}_{\mathrm{SC}}^1\) &
operadic convolution filtration; diagonal/contact support, with scalar
projection \(\hbar_{\mathrm{W}}\chi_N(K)[\bar c]\) &
continuous for the completed convolution \(L_\infty\)-structure,
target pushout extension, and scalar projection &
criterion or projective replacement\\
\(\mathrm{O5}\) &
\(h_{\mathrm{loc}}\in\check C^0(\underline{\mathrm{Hom}}^0)\) &
factorisation-locality degree on disjoint opens of \(\partial X\);
cone comparing a disjoint-union value with the completed tensor product
of component values &
compatible with Costello--Gwilliam structure maps, Čech restriction, and
Ran refinement &
criterion\\
\(\mathrm{O}_{\nabla}\) &
\(h_{\nabla},h_{\nabla,\Theta}\in
\mathcal E^1_{\mathrm{Darb}}\) &
formal Darboux coordinate-torsor direction; completed endomorphism
filtration on the bar family and centre chart &
continuous for the bar-family and centre-chart completions; compatible
with the centrality hierarchy and Darboux-coordinate changes &
typed boundary datum\\
\(\mathrm{O}_{\mathrm{cot}}\) &
\(h_{\mathrm{cot},I}\in
(\mathcal K^{\mathrm{cot}}_\partial(I))^1\) &
interval cosheaf filtration, principal-part current degree, primitive
trace projection, and unreduced open-observable filtration &
continuous for the completed current convolution complex; compatible
with extension by zero, multiplication of smooth functions with
distributions, and factorization products &
typed boundary datum\\
\(\mathrm{O6}\) &
\(\kappa_{\operatorname{Cone}}\) has degree \(-1\) &
complete separated cone filtration, bounded below in each cohomological
degree &
continuous endomorphism of the cone; filtered Neumann series converges on
the associated graded contraction &
quasi-isomorphism datum\\
\(\mathrm{O7}\) &
\(\Gamma\in C^0_{\mathrm{CE}}\), \(H^{\mathrm{top}}_{n+1}\in(F^{n+1}/F^{n+2})^1\) &
weight filtration \(\rho\), holomorphic-translation Chevalley--Eilenberg
degree, and topological-transfer arity filtration &
continuous in \(\mathcal C_\rho\); finite holomorphic propagation and
SDR-choice independence &
optional topologisation criterion
\end{tabular}
\end{center}
The formal-Darboux trace-sector theorem supplies the arity-\(\le2\)
boundary value entering \(\mathrm{O0}\), the formal fibre entering
\(\mathrm{O2}\), and the scalar anomaly value in \(\mathrm{O4}\).  It
supplies the first Taylor coefficient \(\tau_{N,1}^{\mathrm{tr}}\)
entering \(\mathrm{O}_{\mathrm{Hoch}}\); the coordinate expression
\(\Phi_N^{\mathrm{gen}}\) is obtained by continuous duality.  The
centrality primitives making \(\Phi_N^{\mathrm{Hoch}}\) a continuous
Hochschild twisting cochain are separate
data in the Hochschild-centre row.  The
remaining rows are independent construction problems in the displayed
complexes.  The flat-transport row \(\mathrm{O}_{\nabla}\) is required
when the boundary open--closed datum is functorial over the formal
Darboux coordinate torsor; it is a typed boundary datum beyond the five
native OCA components.  The
cotangent-current row \(\mathrm{O}_{\mathrm{cot}}\) is required when the
closed source is the full shifted-cotangent BF algebra
\(\mathfrak h\ltimes P[1]\); the trace-sector theorem accounts only for
its reduced principal-part image after primitive projection.
The vanishing of the action tower, the Hochschild-centre lift tower,
and the five global components, together with chosen null-homotopies, is
exactly the condition that
\(\mathrm{OCA}_N\) be an
\(\mathrm{SC}^{\mathrm{HT}}_{(1,2)}\)-morphism on \(\partial X\).
Contractibility of its cone \(\mathrm{O6}\)
(Definition~\ref{defn:oca-comparison-cone};
Theorem~\ref{thm:filtered-oca-cone-criterion}) gives the induced
quasi-isomorphism.  The optional row \(\mathrm{O7}\) is required only
for the \(E_3^{\mathrm{top}}\) closed-colour target.
Before a common deformation complex is constructed, the five components
of \(\operatorname{Ob}^{\mathrm{global},5}\) are native cohomology
classes in the following obstruction complexes.
Here \(d_{\mathrm{QME}}^{L}=Q+\{I[L],-\}_{L}+\hbar_{\mathrm{QME}}\Delta_L\) is the
Costello scale-\(L\) quantum master differential on local functionals,
\begin{equation}\label{eq:five-component-complexes}
\begin{array}{l|l}
  \mathsf{Class} & \mathsf{Native\ source}\\\hline
	  \operatorname{Ob}^{\mathrm{ch}\,E_2}_{\mathrm{Deligne}}
	    & H^2\!\left(
	        \mathfrak{Def}^{\mathrm{HT}}_{\mathrm{Mor}\mbox{-}E_2}
	        (\mathcal C^{\mathrm{op}}_\partial,b_p)\right)\\
	  \operatorname{Ob}^{\mathrm{desc}}_{\partial}
	    & H^2\!\left(
	      \operatorname{Tot}\check C^\bullet_{\mathrm{Weiss}}
	      (\mathfrak U;\mathcal E^\bullet_{\mathrm{desc}})\right)\\
  \operatorname{Ob}^{\mathrm{QME}}_{\mathrm{all\mbox{-}order}}
    & H^1(\mathcal Q^\bullet_{\mathrm{QME}}(I))\\
  \operatorname{Ob}^{\mathrm{anom}}
    & H^2(\mathrm{Conv}_{\mathrm{SC}}\!\bigl(\mathcal A^{\mathrm{cl}}_{\mathrm{bulk}},
       Z^{\mathrm{der},\mathrm{Mor}}_{E_2^{\mathrm{HT}}}\bigr))\\
	  \operatorname{Ob}^{\mathrm{loc}}
	    & H^1\!\left(
	      \operatorname{Tot}\check C^\bullet_{\mathrm{Weiss,Ran}}
	      (\mathfrak U;\mathcal E^\bullet_{\mathrm{loc}})\right)
\end{array}
\end{equation}
Put
\[
  \mathrm{CE}_{\bar A}^{\mathbf 1}
  =
  C^\bullet_{\mathrm{Lie}}(\bar A;\C\gamma_{\mathbf 1}).
\]
The QME row in \eqref{eq:five-component-complexes} is the following
filtered tower, not a single analytic theorem supplied by notation:
\begin{center}
\scriptsize
\renewcommand{\arraystretch}{1.2}
\begin{tabular}{@{}p{0.16\textwidth}|p{0.23\textwidth}|p{0.09\textwidth}|p{0.23\textwidth}|p{0.09\textwidth}@{}}
\(\mathsf{Entry}\) & \(\mathsf{Complex}\) & \(\mathsf{Degree}\) &
\(\mathsf{Primitive}\) & \(\mathsf{Status}\)\\\hline
scalar projection &
\(\operatorname{Hom}(\operatorname{gr}_{F}\mathfrak Q^\bullet,
\mathrm{CE}_{\bar A}^{\mathbf 1}[1])\) &
\(1\) &
filtration-lowering correction to \(\widehat\sigma_{\mathrm{sc}}\) &
criterion\\
non-scalar residual &
\(\mathcal Q^\bullet_{w,M}(I)=\ker\widehat\sigma_{\mathrm{sc},M}\) &
\(1\) &
\(C_{n,M}\in\mathcal Q^0_{w,M}(I)\), \(d_MC_{n,M}=-\mathfrak o^{\mathrm{ns}}_{n,M}\) &
criterion\\
tower compatibility &
\(\varprojlim\nolimits^1_MH^0(\mathcal Q^\bullet_{w,M}(I))\) &
secondary &
compatible inverse-limit primitives &
required
\end{tabular}
\end{center}
The scalar row is killed only after the successive classes
\(o_{\sigma,w}^{(r)}\) vanish; in the ordinary scalar-reduced
\(\lie{gl}_N\) branch its Capelli input is the raw Rees/Weyl class
\(\hbar_{\mathrm W}\chi_N(K)[\bar c]
=\hbar_{\mathrm W}N[\bar c]\).  The Costello loop parameter
\(\hbar_{\mathrm{QME}}\) governs the QME filtration and does not
normalize the constant-Hamiltonian extension.
The non-scalar row is killed precisely by the counterterm criterion of
Theorem~\ref{thm:wt-nonscalar-counterterm-criterion}; a closed-exchange
branch is governed by the complex and maps of
Theorem~\ref{thm:wt-closed-exchange-counterterm-criterion}.
The common direct-sum language is valid only after comparison maps from
these native sources to a master deformation complex have been fixed.
Each displayed source carries the native null-homotopy for its component.
The four-case obstruction table of
Theorem~\ref{thm:four-curvature-taxonomy} records the local curvature type
only for components already represented in a fixed deformation complex:
the formal all-arity Hamiltonian-current kernels on the polydisk are
constructed in
Theorem~\ref{app:thm:formal-polydisk-all-arity-current-operad}; the
Deligne and descent classes remain target-construction obligations
for the Morita \(E_2^{\mathrm{HT}}\)-extension complex
\(\mathfrak{Def}^{\mathrm{HT}}_{\mathrm{Mor}\mbox{-}E_2}
(\mathcal C^{\mathrm{op}}_\partial,b_p)\) and the
Costello--Gwilliam descent object.  The QME class lies in
case (2) (counterterm coherence); the anomaly class lies in case (3)
(projective coherence); and the locality class lies in case (2)
(counterterm coherence).
\end{thm}

\begin{proof}
\emph{Action before morphism.}  A morphism of
\(\mathrm{SC}^{\mathrm{HT}}_{(1,2)}\)-algebras is defined only after the
source and target carry compatible
\(\mathrm{SC}^{\mathrm{HT}}_{(1,2)}\)-algebra structures.  Row
\(\mathrm{O0}\) is therefore prior to the open--closed comparison map.
It is the Maurer--Cartan lifting problem of
Definition~\ref{defn:sc-ht-action-datum}; its obstruction tower is
\(\{\mathfrak o^{\mathrm{SC}}_{n+1}\}_{n\ge2}\).  The trace-sector
formal-Darboux theorem fixes the arity-\(\le2\) boundary condition of
this lifting problem and no higher arity.

\emph{Hochschild chart before centre morphism.}  The local theorem gives
the first Taylor coefficient
\(\tau_{N,1}^{\mathrm{tr}}\colon\mathfrak g\to
C^\bullet_{\mathrm{Hoch},\mathrm{adm}}(A,A)[1]\).  Continuous
coordinate duality gives \(c_f\mapsto\theta_f\) and
\(u_f\mapsto M_{J(f)}\) only afterward.  A continuous Hochschild
twisting cochain \(\Phi_N^{\mathrm{Hoch}}\) is the stronger datum
\(\tau_N=\mathcal C_{\Phi_N^{\mathrm{Hoch}}}\) of
Definition~\ref{defn:typed-hochschild-centre-lift-datum}.  Its
Maurer--Cartan curvature contains the Lie-bracket defect, the action
defect, product centrality, \(P_0\)-centrality, and the Stasheff
Jacobiator.  The primitives
\(H^{\mathrm{Lie}},H^{\mathrm{act}},H^{\mathrm{prod}},H^{P_0},H^3,\ldots\)
are therefore not consequences of the trace formula.  They are the
separate centrality/equivariance data that turn the chart of the Morita
centre into a genuine Hochschild-centre map; their obstruction is the
relative class
\(\operatorname{Ob}^{\mathrm{cent}/\mathrm{gen}}_{m+1}\) of
Definition~\ref{defn:relative-hochschild-centre-extension-obstruction}.

\emph{Descent before equivalence.}  Local formal-Darboux comparison maps
\(\mathrm{OCA}_{N,a}\) on a collared Weiss cover become a global
comparison morphism only after they form the total cocycle
\(\mathfrak F_{\mathrm{desc}}\) of
Definition~\ref{defn:typed-brane-link-descent-cone-datum}.  The first
curvature is the triple-overlap class
\(\delta_{\check C}H^{(1)}\); the higher curvatures are the successive
components of \(d_{\mathrm{desc}}\mathfrak F_{\le k}\).  Their
primitives are exactly the homotopy-coherent \v{C}ech--Weiss descent
cochains.  Once these cochains exist, the descended map is still only a
morphism.  The quasi-isomorphism assertion is the separate acyclicity of
\(\operatorname{Cone}(\mathrm{OCA}_N)\), supplied either by a direct
contracting homotopy or by the filtered associated-graded contraction
and Neumann-series lift of Theorem~\ref{thm:filtered-oca-cone-criterion}.

\emph{Cotangent current before QME.}  The full Hamiltonian BF source
contains the shifted-cotangent summand \(P[1]\).  The trace map
\(J(f)\) and the CE/PV coordinate theorem see its reduced
principal-part image only after primitive trace projection.  The
interval current algebra
\(\mathfrak g_I^{\mathrm{cur}}\) of
Definition~\ref{defn:typed-cotangent-current-lift-datum} is therefore
the typed source for the missing cotangent modes.  Theorem~\ref{thm:reduced-principal-part-boundary-current}
constructs the reduced support-local brackets.  Passing from those
brackets to an unreduced factorization-centre morphism is exactly the
homotopy-fibre problem defining
\(\mathcal K^{\mathrm{cot}}_\partial(I)\).  Its obstruction class
\(\operatorname{Ob}^{\mathrm{cot}}_\partial(I)\) measures the failure to
choose a primitive projection and centrality homotopies against
arbitrary open observables.  Hence \(\mathrm{O}_{\mathrm{cot}}\) is
independent of the scalar trace-sector theorem and precedes any
Costello graph or all-order QME counterterm enhancement.

\emph{Relation with the three source obstructions.}  The three source
obstructions of Theorem~\ref{thm:main-global-deligne-criterion} are the
ordered triple
\[
 \bigl(
 \operatorname{Ob}^{\mathrm{ch}\,E_2}_{\mathrm{Deligne}},
 \operatorname{Ob}^{\mathrm{desc}}_{\partial},
 \operatorname{Ob}^{\mathrm{QME}}_{\mathrm{all\mbox{-}order}}
 \bigr).
\]
The QME entry is the scalar-contact reduction and non-scalar
counterterm tower of the Costello local-functional obstruction problem
assembled in the Roos complex
\(\mathcal Q^\bullet_{\mathrm{QME}}(I)\): first the scalar projection
must lift through the \(o_{\sigma,w}^{(r)}\) classes, then the residual
finite-window scale-\(L\) classes and their RG transition defects must
vanish.  When a common
filtered obstruction complex has been constructed, the ghost-number
filtration separates the projective-anomaly representative
\[
  \operatorname{Ob}^{\mathrm{anom}}\in
  H^2(\mathrm{Conv}_{\mathrm{SC}}(
  \mathcal A^{\mathrm{cl}}_{\mathrm{bulk}},
  Z^{\mathrm{der},\mathrm{Mor}}_{E_2^{\mathrm{HT}}}))
\]
whose trace-sector scalar is the raw Rees/Weyl class
\(\hbar_{\mathrm{W}}\chi_N(K)[\bar c]
=\hbar_{\mathrm{W}}N[\bar c]\), and
the factorisation-degree filtration separates the locality representative
\[
  \operatorname{Ob}^{\mathrm{loc}}\in
  H^1\!\left(
  \operatorname{Tot}\check C^\bullet_{\mathrm{Weiss,Ran}}
  (\mathfrak U;\mathcal E^\bullet_{\mathrm{loc}})\right).
\]
Definition~\ref{defn:typed-anomaly-locality-datum} gives the two
possible anomaly interpretations.  In the strict interpretation the
entire operadic curvature has a primitive.  In the projective
interpretation only the reduced curvature is null-homotopic, while the
central scalar projection is the raw Rees/Weyl trace class
\(\hbar_{\mathrm{W}}\chi_N(K)[\bar c]
=\hbar_{\mathrm{W}}N[\bar c]\).  Locality is independent of that scalar
choice: it is the Ran disjoint-union cone measuring whether the value on
\(\coprod_i U_i\) agrees with the completed tensor product of the values
on the \(U_i\).
The residual QME entry is a filtered tower in the Roos complex
\(\mathcal Q^\bullet_{\mathrm{QME}}(I)\): scalar-projection lift
classes \(o_{\sigma,w}^{(r)}\), reduced finite-window scale-\(L\)
classes, Costello RG transition defects, and finite-window
primitive-compatibility classes.  Thus
\eqref{eq:five-component-complexes} records five native sources before
comparison, and records five graded summands only after the comparison
maps to \(\mathfrak K_{\mathcal T}\) have been fixed.

\emph{Independence.}  The action tower, the five source complexes of
\eqref{eq:five-component-complexes}, the comparison cone, and the
optional topological transfer complex are native complexes.  They become
bidegrees of a master deformation complex
\(\mathfrak K_{\mathcal T}\) only after comparison maps into that
complex have been fixed: the action tower is graded by operadic arity,
the Deligne and descent classes by their Čech--Weiss--Ran cover degree
on \(\partial X\), the QME class by the renormalisation-group filtration
in \(\ker\widehat\sigma_{\mathrm{sc}}\), the anomaly class by ghost
number in the operadic convolution complex, the locality class by
factorisation degree in
\(\mathrm{Fact}_{\mathrm{SC}^{\mathrm{HT}}_{(1,2)}}\), the cone by the
comparison filtration, and the optional topological transfer by the
\(\rho\)-weight and topological arity filtrations.  After the comparison
maps are fixed, no two complexes share both bidegree and graded piece.
The five-component obstruction
\(\operatorname{Ob}^{\mathrm{global},5}\) of
\eqref{eq:obs-five-component} is one visible summand of the typed ledger
\eqref{eq:typed-total-obstruction-ledger}.  A typed OCA datum is a
compatible system consisting of the supplied action Maurer--Cartan lift,
null-homotopies in the five global summands, a contracting homotopy of
\(\operatorname{Cone}(\mathrm{OCA}_N)\), and, when requested, the
topological transfer primitive.  The action lift gives the source and
target \(\mathrm{SC}^{\mathrm{HT}}_{(1,2)}\)-structures; the five global
null-homotopies construct the comparison morphism; the cone contraction
gives the quasi-isomorphism.

  \emph{Costello--Li control.}  The Costello--Li construction
  (Statement~\ref{stmt:costello-li-flat-bcov},
  Statement~\ref{stmt:costello-bv-construction}) supplies the model for the
  renormalised local-functional complex in its hypothesis class.  In the
  Hamiltonian BF defect sector, the same complex is a valid QME target only
  after the specialization, locality, and counterterm data have been fixed.
The anomaly and locality rows require the operadic convolution complex,
the locality cocycle, and the corresponding null-homotopies.  If
topologisation is imposed, the homotopy-transfer conditions are the
vanishing of \([\tau]\) and of the tower
\(\mathfrak o^{\mathrm{top}}_{n+1}\), together with
\((\mathrm{T1})\)--\((\mathrm{T6})\), as stated in
Theorem~\ref{thm:topologisation-by-stress-tensor}.  The Deligne and
descent classes remain governed by the mixed-HT
\(E_2^{\mathrm{HT}}\)-Deligne problem and the Costello--Gwilliam descent
spectral sequence, respectively.

\emph{Four obstruction mechanisms.}
\(\operatorname{Ob}^{\mathrm{ch}\,E_2}_{\mathrm{Deligne}}\) and
\(\operatorname{Ob}^{\mathrm{desc}}_{\partial}\) lie in the
target-replacement cell
(Theorem~\ref{thm:four-curvature-taxonomy} case~(4)) until the Morita
mixed-HT \(E_2^{\mathrm{HT}}\)-centre action and the Costello--Gwilliam descent
datum have been constructed.  On a formal-Darboux polydisk the all-arity
Hamiltonian-current kernel part is already the tree-kernel action of
Theorem~\ref{app:thm:formal-polydisk-all-arity-current-operad}; the
remaining representatives are the boundary values of the categorical
global null-homotopy problem;
\(\operatorname{Ob}^{\mathrm{QME}}_{\mathrm{all\mbox{-}order}}\) is a
counterterm class in \(\mathcal Q^\bullet_{\mathrm{QME}}(I)\) only after
the Costello local counterterm complexes, RG transition maps, and window
projections have been fixed; \(\operatorname{Ob}^{\mathrm{loc}}\) is a
factorization-locality class (case~(2)) only after the locality complex
and its homotopies have been fixed;
\(\operatorname{Ob}^{\mathrm{anom}}\) is the projective Lie-anomaly
class (case~(3)).  Its source is the unscaled canonical
constant-Hamiltonian extension
\(\widehat{\mathfrak g}_{\mathrm{cent}}\) of
\eqref{eq:cent-extension}; its anomaly representative is the raw
Rees/Weyl pushout under \(\chi_N(K)=N\).
\end{proof}

\begin{rmk}[Local trivialisation of the five-component obstruction]
\label{rmk:local-trivialisation-five-component-obstruction}
On the contractible formal-Darboux neighbourhood of a brane vacuum
\(b\), the trace-sector representatives entering the typed obstruction
ledger are explicit: the all-arity tree-kernel current operad of
Theorem~\ref{app:thm:formal-polydisk-all-arity-current-operad}, whose
binary generator is \(K^{(2)}_\Delta\), the first Taylor coefficient
\(\tau_{N,1}^{\mathrm{tr}}\) with coordinate-dual CE/PV expression
\(c_f\mapsto\theta_f,\ u_f\mapsto J(f)\), the unscaled class
\([\bar c]\), the rank character \(\chi_N(K)=N\), and the raw Rees/Weyl
trace class \(\hbar_{\mathrm W}N[\bar c]\).  The tree-kernel operad is the local representative of the
separate Swiss-cheese action row.  The five-component obstruction
\(\operatorname{Ob}^{\mathrm{global},5}\) records the Morita mixed-HT
\(E_2^{\mathrm{HT}}\)-centre lift, Costello--Gwilliam descent, all-order QME
counterterm tower, anomaly, and locality rows along \(\partial X\).  The
formal-Darboux stalk theorem (Theorem~\ref{thm:oca-formal-darboux-fibre},
\(=\) Theorem~\ref{thm:main-local}) is the trace-sector comparison
fibre at \(b\).
\end{rmk}

\subsection{Local realization and descent under hypotheses}\label{ssec:local-relative-realizations}

The formal-Darboux theorem fixes only the local fibre for a global
comparison datum whose formal restriction at one brane vacuum is the
trace-sector first Taylor boundary coefficient.  A global open--closed quasi-isomorphism
requires this fibre
and the typed data of \eqref{eq:typed-total-obstruction-ledger}: a
mixed-HT Swiss-cheese action Maurer--Cartan lift, a Morita mixed-HT
\(E_2^{\mathrm{HT}}\)-centre Maurer--Cartan lift extending the formal
tree-kernel current operad, a Hochschild-centre Maurer--Cartan lift, flat transport
over the formal Darboux coordinate torsor, an unreduced cotangent-current
lift, a Costello--Gwilliam descent null-homotopy on the brane-link
boundary, an all-scale QME counterterm tower with RG/window
compatibility primitives, anomaly and locality
null-homotopies, and a contracting homotopy of
\(\operatorname{Cone}(\mathrm{OCA}_N)\).  Topological closed colour
requires the transfer datum of
Definition~\ref{defn:topological-transfer-complex}.  Stein or affine
hypotheses simplify ordinary sheaf cohomology after these complexes have
been constructed.

\subsubsection{Local trace-sector theorem at the formal-Darboux stalk}
\label{sssec:local-trace-sector-stalk}

Theorem~\ref{thm:main-local} states the trace-sector first Taylor
boundary coefficient and, after continuous coordinate duality, the
scalar-stable CE/PV coordinate isomorphism
inside the chosen admissible category \(\mathcal B^{\mathrm{adm}}\).
The admissibility conditions are structural hypotheses on the
completion, not consequences of the finite-\(N\) brane algebra.

\begin{lem}[Admissible formal-Darboux data]
\label{lem:admissible-formal-darboux-data}
In the formal-Darboux setting
\(X=\R^2_{\mathrm{top}}\times\widehat{\C^2}_{\mathrm{hol}}\), the
data used in Theorem~\ref{thm:main-local} consist of the following
admissible structures:
\begin{enumerate}[label=\textup{(\roman*)}]
  \item the residue dual
        \(P=\mathfrak h^\vee_{\mathrm{cont}}\) and the pairing
        \(\mathfrak h\otimes P\to\C\),
        \(\langle f,\omega\rangle=\Res_0 f\omega\), identified by
        Grothendieck--Matlis local duality on
        \(H^2_{\mathfrak m}(R)\)
        \cite{hartshorne-local-cohomology}*{Ch.~III, \S2}
        \cite{kunz-residues-duality}*{Ch.~5};
  \item a weighted or pro-Matlis topology in which the diagonal Casimir
        \(C=\sum_{a+b>0}z_1^az_2^b\otimes\rho_{a,b}\) converges and the
        Hamiltonian coadjoint action satisfies the seminorm estimate
        \(p^*_n(f\cdot\rho)\leq C_n\,p_{n+r}(f)\,p^*_{n+r}(\rho)\)
        of Proposition~\ref{prop:completion-promatlis-continuity};
  \item the diagonal local-cohomology kernel
        \(K^{(2)}_\Delta=d\zeta_1d\zeta_2/(\zeta_1\zeta_2)\) is
        bracket- and kernel-admissible by
        Proposition~\ref{prop:diagonal-cohomology-operad}: the
        binary singular product of \eqref{eq:singular-product-linear} factors through
        \(K^{(2)}_\Delta\) with principal-part residue in the
        Grothendieck--Matlis dual;
  \item the completed Hochschild-centrality convolution algebra
        \(\mathcal K^{\mathrm{cent}}_{\partial,N}\) of
        Definition~\ref{defn:typed-hochschild-centre-lift-datum}, with
        operation list
        \(\mathcal O^{\mathrm{cent}}_{\partial,N}\), kernel list
        \(\mathcal K^{\mathrm{ker,cent}}_{\partial,N}\) for the
        restricted \(P_0\) row, and primitives
        \(\{H^{\mathrm{Lie}},H^{\mathrm{act}},H^{\mathrm{prod}},
        H^{P_0},H^3,\ldots\}\) lying in the displayed continuous Hom
        complex;
  \item the brane-defect QME tower datum of
        Definition~\ref{defn:typed-brane-defect-qme-tower}, with
        operation list \(\mathcal O^{\mathrm{QME}}_{w,r,M,L}\), kernel
        list \(\mathcal K^{\mathrm{ker,QME}}_{w,r,M,L}\), scalar-contact
        projection, finite-scale curvature representatives, and
        Milnor/Roos compatibility primitives for any counterterm tower
        used in a quantum or unreduced \(P_0\)-centrality extension.
\end{enumerate}
\end{lem}

\begin{proof}
The first item is Grothendieck--Matlis local duality.  The remaining
items are the explicit hypotheses under which the coordinate CE/PV
identity extends to a bracketed and completed trace-sector map.
The strict finite-\(N\) polynomial model is the algebraic quotient
before these completions, kernels, and Milnor/Roos compatibility
classes are chosen.
\end{proof}

\begin{cor}[Local trace-sector formal-Darboux fibre]
\label{cor:local-trace-sector-deligne}
With the admissible data of
Lemma~\ref{lem:admissible-formal-darboux-data}, the formal-Darboux stalk
has the trace-sector first Taylor boundary coefficient and its
coordinate-dual expression from
Theorem~\ref{thm:main-local}.  After stable scalar reduction and
admissible completion it is the continuous CE/PV coordinate isomorphism; a
Hochschild cochain map requires the centrality Maurer--Cartan lift
of Definition~\ref{defn:admissible-hochschild-centre-lift}.
\end{cor}

\begin{proof}
This is Theorem~\ref{thm:main-local} with the admissible data
spelled out in Lemma~\ref{lem:admissible-formal-darboux-data}.
\end{proof}

\subsubsection{Global criterion under hypotheses on contractible
holomorphic-symplectic targets}
\label{sssec:global-contractible}

Contractibility and Stein vanishing act on the sheaf-cohomological
rows of the obstruction complex.  The mixed-HT Swiss-cheese action, the
chiral \(E_2\) all-arity operation, the QME counterterm tower, and the
comparison-cone contraction of Definition~\ref{defn:oca-comparison-cone}
are additional chain-level data.

\begin{thm}[Contractible Stein supplied-data criterion for the global
Mixed HT Deligne map]
\label{thm:contractible-stein-descent-criterion-HS}
On a contractible Stein holomorphic-symplectic surface
\(Y_{\mathrm{hol}}\), with brane-link boundary in
\(X=\R^2_{\mathrm{top}}\times Y_{\mathrm{hol}}\), the open--closed
comparison
\[
  \mathrm{OCA}_N\colon
    \mathcal A^{\mathrm{cl}}_{\mathrm{bulk}}
    \longrightarrow
    Z^{\mathrm{der},\mathrm{Mor}}_{E_2^{\mathrm{HT}}}
      (\mathcal C^{\mathrm{op}}_\partial)
\]
is a supplied morphism in a strict-filtered model of
\(\mathcal B^{\mathrm{adm}}\).  Its quasi-isomorphism datum consists of
the following typed chain-level data:
\begin{enumerate}[label=\textup{(\roman*)}]
  \item a mixed-HT Swiss-cheese action Maurer--Cartan lift extending the
        trace-sector arity-\(\le2\) boundary value;
  \item an all-arity mixed-HT \(E_2^{\mathrm{HT}}\) Morita-centre
        diagonal-kernel action extending
        \(K^{(2)}_\Delta\);
  \item an admissible Hochschild-centre chart lift with centrality
        primitives for
        \(\operatorname{Ob}^{\mathrm{cent}}_{\mathrm{Hoch}}\);
  \item flat transport over the formal Darboux coordinate torsor, with
        primitives for
        \(\operatorname{Ob}^{\mathrm{Darb}}
        =([\Omega_\nabla],[\mathcal R_{\nabla,\Theta}])\), whenever the
        comparison is required to be functorial in Darboux coordinates;
  \item an unreduced cotangent-current lift with interval-compatible
        primitives for \(\operatorname{Ob}^{\mathrm{cot}}_\partial\)
        whenever the closed source is the full shifted-cotangent BF
        algebra;
  \item a Costello--Gwilliam descent datum on the brane-link boundary
        with representative
        \(\mathcal R^{\mathrm{desc}}_{\partial}\) and primitive
        \(h_{\mathrm{desc}}\) satisfying
        \(d_{\mathrm{desc}}h_{\mathrm{desc}}
        =-\mathcal R^{\mathrm{desc}}_{\partial}\);
  \item a brane-defect QME tower datum: finite-scale Costello
        interactions, BV Laplacians, RG transition maps, finite-window
        projections, scalar-contact projection, operation list
        \(\mathcal O^{\mathrm{QME}}_{w,r,M,L}\), kernel list
        \(\mathcal K^{\mathrm{ker,QME}}_{w,r,M,L}\), counterterm
        primitives, and Roos compatibility primitives as in
        Definition~\ref{defn:typed-brane-defect-qme-tower};
  \item a reduced anomaly primitive
        \(d_{\mathrm{Conv}}h_{\mathrm{anom}}
        =-(\mathcal R^{\mathrm{anom}})_{\mathrm{red}}\) and a locality
        primitive \(\partial_{\mathrm{loc}}h_{\mathrm{loc}}
        =-\mathcal R^{\mathrm{loc}}\);
  \item a contracting homotopy for
        \(\operatorname{Cone}(\mathrm{OCA}_N)\) in
        \(\mathcal B^{\mathrm{adm}}\) after the preceding data have
        been installed.
\end{enumerate}
After these data are installed, all components of
\(\operatorname{Ob}^{\mathrm{typed}}\) in
\eqref{eq:typed-total-obstruction-ledger}, except the optional
topologisation row, are supplied by representatives and primitives in
their governing complexes, and the comparison cone is acyclic.  Replacing
the anomaly null-homotopy in the strict anomaly/locality row by the raw
Rees/Weyl trace pushout of the canonical extension gives the projective comparison of
Theorem~\ref{thm:projective-stein-descent-criterion-HS}.
\end{thm}

\begin{proof}
The first datum supplies the source and target
\(\mathrm{SC}^{\mathrm{HT}}_{(1,2)}\)-algebra structures.  The
mixed-HT \(E_2^{\mathrm{HT}}\)-centre, Hochschild-centre,
Darboux-torsor, cotangent-current, descent, QME-counterterm,
anomaly/locality, and cone entries are exactly the non-optional rows
of \eqref{eq:typed-total-obstruction-ledger}.  Contractibility removes
monodromy for locally constant descent data, and Stein
vanishing removes coherent sheaf cohomology obstructions once the
relevant sheaves and counterterm complexes have been constructed.
Contractibility changes neither the unscaled Lie extension class
\([\bar c]\) nor its raw Rees/Weyl trace pushout
\(\hbar_{\mathrm W}\chi_N(K)[\bar c]\).  The strict anomaly/locality datum is therefore a
genuine primitive in the operadic convolution and locality complexes.
The final datum is the acyclicity statement for the comparison cone.
Steinness supplies no Swiss-cheese action, no all-arity operation, no
Hochschild-centre lift, no cotangent-current lift, no counterterm tower,
and no cone contraction.
\end{proof}

\subsubsection{Projective supplied-data criterion on Stein
holomorphic-symplectic targets}
\label{sssec:projective-stein-descent-criterion}

On a general Stein holomorphic-symplectic surface \(Y\), coherent
cohomology vanishing controls the sheaf-theoretic rows after the
Swiss-cheese action, mixed-HT \(E_2^{\mathrm{HT}}\)-centre,
Darboux-torsor, cotangent-current, descent, QME counterterm tower,
locality, and cone
complexes have been constructed.  The unscaled canonical
constant-Hamiltonian extension supplies \([\bar c]\); the rank character
supplies \(N\), and the raw Rees/Weyl commutator supplies
\(\hbar_{\mathrm W}\).  Their target pushout records the projective part
of the trace anomaly and replaces only the strict anomaly primitive.

\begin{defn}[Projective Stein residual class]
\label{defn:projective-stein-residual-class}
For a Stein holomorphic-symplectic target \(Y_{\mathrm{hol}}\) with
brane-link boundary, the projective residual coefficient group is
\[
  \mathcal H^{\mathrm{proj}}_Y
  =
  H^1(X,\C_X)
  \oplus
  H^2_{\mathrm{Lie}}(\bar A,\C)[[\hbar_{\mathrm{W}}]],
  \qquad
  X=\R^2_{\mathrm{top}}\times Y_{\mathrm{hol}} .
\]
The projective residual class is
\[
  \operatorname{Ob}^{\mathrm{proj}}_Y
  =
  \bigl(\operatorname{per}_X,
  \hbar_{\mathrm{W}}\chi_N(K)[\bar c]\bigr)
  =\bigl(\operatorname{per}_X,\hbar_{\mathrm{W}}N[\bar c]\bigr)
  \in \mathcal H^{\mathrm{proj}}_Y .
\]
Here \(\operatorname{per}_X\) is the period class of the chosen global
Hamiltonian descent datum, and \([\bar c]\) is the Lie two-cocycle class
of the unscaled constant-Hamiltonian extension.  The second component is
its raw Rees/Weyl trace pushout.  A determinant-line or modular-line
curvature interpretation requires a supplied line, connection, and
Atiyah-class pullback.  A comparison with the common
mixed-HT obstruction complex is additional structure: it is a
degree-normalized linear map
\[
  \iota^{\mathrm{proj}}_Y\colon
  \mathcal H^{\mathrm{proj}}_Y
  \longrightarrow H^1(\mathfrak K_{\mathrm{HT}})
\]
compatible with the scalar-contact projection and the Hamiltonian
period row of Definition~\ref{defn:common-ht-deformation-complex}.  The
Lie \(H^2\)-summand is shifted to obstruction degree \(1\) by the
scalar-Lie row convention in \(\mathfrak K_{\mathrm{HT}}\).  The residual
lives in the common obstruction complex exactly after
\(\iota^{\mathrm{proj}}_Y\) is part of the supplied datum; before that
comparison it is the class in \(\mathcal H^{\mathrm{proj}}_Y\).
\end{defn}

\begin{thm}[Projective Stein supplied-data criterion]
\label{thm:projective-stein-descent-criterion-HS}
On a Stein holomorphic-symplectic surface $Y_{\mathrm{hol}}$
(not necessarily contractible), let
\(X=\R^2_{\mathrm{top}}\times Y_{\mathrm{hol}}\), and let \(N\) Dirac
branes sit along \(L=\R_t\times\{0_s\}\times\{b\}\).  Suppose the mixed-HT
Swiss-cheese action Maurer--Cartan lift, the all-arity mixed-HT
\(E_2^{\mathrm{HT}}\) Morita-centre operation, the Hochschild-centre
chart lift, flat Darboux-torsor
transport when required, the unreduced cotangent-current lift when the
full shifted-cotangent BF source is used, the brane-link descent datum,
the all-order quantum master equation (QME) counterterm tower datum,
the locality homotopy for \(\operatorname{Ob}^{\mathrm{loc}}\), and the
comparison-cone contraction of Definition~\ref{defn:oca-comparison-cone}
are fixed.
Suppose also that the strict anomaly row
\(\operatorname{Ob}^{\mathrm{anom}}\) retains the raw Rees/Weyl trace
pushout of the canonical unscaled extension
\(\widehat{\mathfrak g}_{\mathrm{cent}}\) under \(\chi_N(K)=N\), rather
than killing it.
Then the supplied OCA map carries its quasi-isomorphism datum in the
projective category obtained from \(\mathcal B^{\mathrm{adm}}\) by
extension along this target pushout of
  \(\widehat{\mathfrak g}_{\mathrm{cent}}\).  The remaining projective
  residual is the class
  \(\operatorname{Ob}^{\mathrm{proj}}_Y\) of
  Definition~\ref{defn:projective-stein-residual-class}.  It becomes a
  class of the common obstruction complex only after the comparison
  \(\iota^{\mathrm{proj}}_Y\) has been constructed, in which case the
  common-complex residual is
  \[
    \iota^{\mathrm{proj}}_Y
    \bigl(\operatorname{Ob}^{\mathrm{proj}}_Y\bigr)
    \in H^1(\mathfrak K_{\mathrm{HT}}).
  \]
\end{thm}

\begin{proof}
With these data fixed, the non-anomaly rows of the typed ledger
\eqref{eq:typed-total-obstruction-ledger} have the required
null-homotopies or contractions.  The anomaly row is not null-homotopic
in the strict category; it is replaced by the target central extension
obtained by the raw Rees/Weyl trace pushout.  Stein vanishing applies
inside the relevant sheaves and renormalised complexes after those
complexes have been constructed.  The
residual scalar anomaly is
\(\hbar_{\mathrm W}\chi_N(K)[\bar c]\): the unscaled Lie central class
is evaluated by the trace character and then placed in the raw
Rees/Weyl commutator.  Together with any period class of the chosen
global Hamiltonian datum, this scalar class forms
\(\operatorname{Ob}^{\mathrm{proj}}_Y\).  Passing to the projective
category records the second summand as a projective central extension of
the trace-sector action; the first summand is the remaining Hamiltonian
period obstruction.  Mapping this pair into the common deformation
complex is exactly the extra comparison
\(\iota^{\mathrm{proj}}_Y\); without it there is no single common-complex
class.
The Stein hypothesis contributes only sheaf and descent acyclicity; it
supplies no Swiss-cheese action, no all-arity operation, no
Hochschild-centre lift, no cotangent-current lift, no counterterm tower,
no locality homotopy, and no cone contraction.
\end{proof}

\subsubsection{Concrete local realizations and global hypotheses}
\label{sssec:concrete-realizations}

\begin{cor}[Concrete local realizations and global hypotheses]
\label{cor:concrete-global-realizations}
The formal-Darboux theorem applies at every smooth point of each
target listed below.  A global open--closed quasi-isomorphism is the conclusion of
Theorem~\ref{thm:contractible-stein-descent-criterion-HS} or
Theorem~\ref{thm:projective-stein-descent-criterion-HS} only after the
data stated in those criteria have been constructed.

The following contractible Stein holomorphic-symplectic surfaces have
no topological monodromy for the local trace-sector system:
\begin{enumerate}[label=\textup{(\arabic*)}]
  \item \(Y=\C^2\) with \(\omega_Y=dz_1\wedge dz_2\) --- the full
        local model of \S\ref{ssec:local-setup};
  \item \(Y=D^2\subset\C^2\) the open polydisk
        \(\{|z_1|<1,|z_2|<1\}\) with the induced symplectic form;
  \item \(Y=T^*\C=\C_p\times\C_q\) with \(\omega_Y=dp\wedge dq\),
        the cotangent bundle of the affine line;
  \item \(Y=T^*\mathbb H\) with \(\mathbb H\) the upper half-plane,
        \(\omega_Y=dp\wedge dq\) on the cotangent;
  \item \(Y=B_r^4\) the open Stein ball of radius \(r\) in \(\C^2\)
        with \(\omega_Y=dz_1\wedge dz_2\);
  \item \(Y=\widehat{\C^2}_{b}\) the formal-Darboux completion at any
        smooth point \(b\) of any holomorphic-symplectic surface
        (the stalk case of Corollary~\ref{cor:local-trace-sector-deligne}).
\end{enumerate}

The following Stein holomorphic-symplectic surfaces carry possible
period or equivariant descent classes and therefore belong to the
projective criterion:
\begin{enumerate}[label=\textup{(\arabic*)},resume]
  \item \(Y=\C\times\C^*\) with \(\omega_Y=dz\wedge d\log w\) --- the
        \(\C^*\)-cylinder, with a period class along the circle;
  \item an affine smooth toric holomorphic-symplectic surface, with
        toric chart gluing recorded by the descent class.
\end{enumerate}

The punctured surface \(\C^2\setminus\{0\}\) is excluded from the
Stein list: by Hartogs extension its holomorphic functions extend to
\(\C^2\), so it is not Stein.
The smooth locus of an affine ADE surface
\(\mathrm{Spec}\,\C[z_1,z_2]^\Gamma\) is likewise a punctured quotient
surface, not a Stein target for the global criterion.  The singular
affine surface requires a derived or stacky target, and its minimal
resolution is listed below as a chartwise local target with compact
descent data.

The following holomorphic-symplectic surfaces have compact analytic
subvarieties or compact descent data.  The local theorem applies on
Stein charts; the global statement requires compact descent:
\begin{enumerate}[label=\textup{(\arabic*$'$)},start=7]
  \item \(Y=T^*E\) for \(E\) an elliptic curve: holomorphic-symplectic
        but globally non-Stein (compact zero section
        \(E\hookrightarrow T^*E\));
  \item \(Y=T^*\Sigma_{g\geq 1}\) for \(\Sigma_g\) a higher-genus
        Riemann surface: holomorphic-symplectic but globally non-Stein
        (compact zero section \(\Sigma_g\hookrightarrow T^*\Sigma_g\));
  \item \(Y=\widetilde{\C^2/\Gamma}\) the minimal resolution of the
        Kleinian quotient: smooth holomorphic-symplectic but globally
        non-Stein (compact ADE exceptional divisor of \(\mathbb P^1\)
        components);
  \item \(Y=\mathrm{Hilb}^N(\C^2)\) the Hilbert scheme of \(N\) points
        on \(\C^2\): smooth quasi-projective but globally non-Stein
        for \(N\geq 2\) (compact punctual locus
        \(\mathrm{Hilb}^N(\C^2)_0\hookrightarrow\mathrm{Hilb}^N(\C^2)\));
        the cyclic-framed stable branch of
        Lemma~\ref{lem:cyclic-framing-hilbert-branch} carries the
        Nakajima ADHM presentation
        \cite{nakajima-hilbert-lectures} as a GIT quotient on the
        stable framed locus where \(\C[\phi_1,\phi_2]\cdot v=\C^N\);
        on each Stein chart of the Hilbert--Chow morphism
        \(\mathrm{Hilb}^N(\C^2)\to S^N(\C^2)\) the local theorem
	        gives the trace-sector chart; descent through the
	        exceptional stratum has a scalar projective obstruction class
	        with coefficients in the target line
	        \(\C[[\hbar_{\mathrm W}]]\gamma_{\mathbf1}\), obtained from the
	        unscaled extension by the raw Rees/Weyl trace pushout.  The trace map
        \(J(f)\) descends to the
        \(\mathrm{Hilb}^N(\C^2)\)-coordinate ring through the cyclic
        framing vector \cite{nakajima-hilbert-lectures}.
\end{enumerate}
Thus the catalogue gives local trace-sector realizations and the
strict, projective, and compact-descent hypotheses for global
comparisons.
\end{cor}

\begin{proof}
Items (1)--(6) are formal or Stein contractible models.  Items
(7)--(8) are Stein or affine only after the smooth-target and descent
data are specified.  Hartogs extension excludes
\(\C^2\setminus\{0\}\) and the smooth locus of the affine ADE
surface from the Stein list.  Items (7')--(10') contain
compact analytic subvarieties and therefore cannot be Stein; the
formal theorem applies chartwise, and the descent classes carry the
global obstruction.
\end{proof}

\begin{cor}[Stein-cover descent: explicit obstruction classes for the
compact-containing realisations]\label{cor:compact-containing-descent}
For each compact-containing realisation \((7')\)--\((10')\) of
Corollary~\ref{cor:concrete-global-realizations}, choose the following
Stein atlas or universal Stein cover and let \(\Gamma_Y\) denote its
descent groupoid:
\begin{center}
\small
\begin{tabular}{@{}p{0.05\textwidth}p{0.18\textwidth}p{0.36\textwidth}p{0.30\textwidth}@{}}
\hline
\textbf{Item} & \textbf{Target \(Y\)} & \textbf{Stein atlas or cover} & \textbf{Descent groupoid \(\Gamma_Y\)}\\
\hline
\((7')\) & \(T^*E\) & \(\C^2=T^*\C\) (universal cover of \(T^*E\)) & \(\Lambda\cong\Z^2\), translation by the lattice \(E=\C/\Lambda\)\\
\((8')\) & \(T^*\Sigma_{g\geq 1}\) & \(T^*\widetilde\Sigma_g\) (\(\C^2\) for \(g=1\); \(\C\times\mathbb H\) for \(g\geq 2\)) & \(\pi_1(\Sigma_g)\)\\
\((9')\) & \(\widetilde{\C^2/\Gamma}\) & affine Stein chart cover of the resolution & gluing groupoid on the ADE exceptional divisor\\
\((10')\) & \(\mathrm{Hilb}^N(\C^2)\) & Bia\l ynicki-Birula chart cover indexed by partitions \(\lambda\vdash N\) & combinatorial gluing on the punctual stratum\\
\hline
\end{tabular}
\end{center}
Assume first that the chart-level supplied data of
Theorem~\ref{thm:projective-stein-descent-criterion-HS} have been fixed
on the Stein atlas: the mixed-HT Swiss-cheese action Maurer--Cartan
lift, the all-arity mixed-HT \(E_2^{\mathrm{HT}}\) Morita-centre
operation, the QME counterterm tower, the canonical unscaled
constant-Hamiltonian extension together with its raw Rees/Weyl pushout
under \(\chi_N(K)=N\) for the projective anomaly row, locality homotopy,
and cone contraction.  The remaining obstruction to the OCA
quasi-isomorphism on \(Y\) is the descent class
\[
  \operatorname{Ob}^{\mathrm{desc}}_{\Gamma_Y}
   \;\in\;
   \check H^2\bigl(
   \Gamma_Y;\C[[\hbar_{\mathrm W}]]\gamma_{\mathbf1}\bigr)
\]
and the OCA descends after a representative
\(\mathcal R^{\mathrm{desc}}_{\Gamma_Y}\) has a compatible descent
primitive \(h_{\Gamma_Y}\) satisfying
\(d_{\mathrm{desc}}h_{\Gamma_Y}
=-\mathcal R^{\mathrm{desc}}_{\Gamma_Y}\); a nonzero class obstructs
descent.
The projective supplied-data criterion
(Theorem~\ref{thm:projective-stein-descent-criterion-HS})
applies on each Stein chart of the atlas after the stated data are
        fixed; descent is the problem of gluing the chartwise OCA data along the
        descent groupoid.  After the coefficient groupoid and comparison
  maps are fixed, the representatives are the period class on \(H^1(E)\) for
  \((7')\), the period class on \(H^1(\Sigma_g)\) for \((8')\), the ADE
  root-lattice gluing class of the exceptional divisor for \((9')\),
  and the tautological gluing class along the Hilbert--Chow exceptional
  stratum for \((10')\).  Each is recorded as a component of the compact comparison class
\(\operatorname{Ob}_{\mathrm{comparison}}\) of
Theorem~\ref{thm:cross-target-restriction}.
\end{cor}

\begin{proof}
On each Stein chart \(U\subset\widetilde Y\), the projective descent
criterion
(Theorem~\ref{thm:projective-stein-descent-criterion-HS}) gives
chartwise projective OCA data only after the mixed-HT
\(E_2^{\mathrm{HT}}\)-centre, descent, QME counterterm tower, locality,
and cone-contraction data are fixed.  The descent
groupoid \(\Gamma_Y\) preserves the holomorphic-symplectic form on
overlaps by construction; the class
\(\operatorname{Ob}^{\mathrm{desc}}_{\Gamma_Y}\) is the obstruction
to lifting the chartwise projective OCA data to a
\(\Gamma_Y\)-descent datum, equivalently the class in
\(\check H^2(\Gamma_Y;
\C[[\hbar_{\mathrm W}]]\gamma_{\mathbf1})\)
classifying central extensions of the descent groupoid by the target
projective scalar line.  Its coefficient is
\(\hbar_{\mathrm W}\chi_N(K)[\bar c]\); the source extension remains
\(\bar cK\).  The chosen Stein atlas supplies the
\v{C}ech--Dolbeault acyclicity row; the OCA quasi-isomorphism on \(Y\)
is the chartwise projective datum together with a null-homotopy of
\(\operatorname{Ob}^{\mathrm{desc}}_{\Gamma_Y}\); a nonzero class
obstructs equivariant gluing.  The named period classes are proposed
cocycle representatives until the coefficient groupoid is fixed.
\end{proof}

\begin{rmk}[Nakajima quiver-variety reading on
\(\mathrm{Hilb}^N(\C^2)\)]\label{rmk:nakajima-quiver-OCA}
On the cyclic-framed stable branch of Corollary item~(10$'$), the
boundary algebra \(A^{\mathrm{cl}}_{\partial,N}\) restricts to the
coordinate ring of \(\mathrm{Hilb}^N(\C^2)\) through the Nakajima
ADHM presentation: the framed ADHM quiver \(\Gamma_N\) (one node
carrying the gauge \(\mathfrak{gl}_N\), one framed boundary node, two
arrows for \((\phi_1,\phi_2)\), one moment-map arrow for \(\psi\)) is
the representation-variety presentation of the underlying derived
commuting stack, and the trace map \(J(f)\) is the trace of the
arrow-monomial on the framed representation.  The OCA trace-sector
map on this branch is the restriction of
Theorem~\ref{thm:main-local} to the cyclic-framed stable locus,
where \(N=1\) gives
\(\widehat{\mathcal O}_{\mathrm{Hilb}^1(\C^2),[b]}\cong\C[[z_1,z_2]]\).
For \(N>1\), the Hilbert scheme has dimension \(2N\), and its completed
local ring is a \(2N\)-dimensional Darboux chart with tautological
rank-\(N\) algebra; \(A_b^{\mathrm{coord}}=\C[[z_1,z_2]]\) remains the
target formal disk acting through that tautological algebra.  The
completed local ring of \(\mathrm{Hilb}^N(\C^2)\) is the corresponding
tautological \(2N\)-dimensional deformation algebra.  This is a
concrete chart-level realization of the trace map; global descent
across the Hilbert--Chow exceptional locus is a compact comparison
problem.
\end{rmk}

\begin{cor}[Concrete local comparison loci]
\label{cor:additional-concrete-realizations}
At every smooth holomorphic-symplectic point of the following spaces,
Theorem~\ref{thm:main-local} gives the formal-Darboux trace-sector
map:
\begin{enumerate}[label=\textup{(\arabic*)},start=14]
  \item \(T^*S\), for a contractible Stein open Riemann surface \(S\);
  \item a contractible pseudoconvex Reinhardt domain in \(\C^2\);
  \item \(K3\setminus C\), for an ample curve \(C\subset K3\);
  \item a smooth two-dimensional character variety;
  \item a smooth two-dimensional Coulomb or Higgs branch of a
        \(3d\ \mathcal N=4\) gauge theory;
  \item a smooth two-dimensional multiplicative quiver variety;
  \item a smooth two-dimensional hypertoric variety;
  \item a smooth affine wild character variety of complex dimension
        \(2\).
\end{enumerate}
Items (14)--(15) fall under the contractible Stein criterion when the
mixed-HT \(E_2^{\mathrm{HT}}\)-centre, descent, QME counterterm tower,
and comparison-cone
data of Definition~\ref{defn:oca-comparison-cone} are fixed.
Items (16)--(21) fall under the projective Stein criterion when the
same data are fixed, the rank character \(\chi_N(K)=N\) is retained,
and the projective scalar is taken in the raw Rees/Weyl commutator.
Singular quotients require a derived or stacky replacement of the
smooth-target hypothesis.
\end{cor}

\begin{proof}
The input from the examples is the following.  In items \textup{(14)}
and \textup{(15)} the displayed contractible Stein models give the
Stein acyclicity required by
Theorem~\ref{thm:contractible-stein-descent-criterion-HS}.  In items
\textup{(16)--(21)} one works on the specified smooth
holomorphic-symplectic affine locus, or on a chosen Stein atlas of that
locus, before applying
Theorem~\ref{thm:projective-stein-descent-criterion-HS}.  At the chosen
point, the only local geometric input is the completed
holomorphic-symplectic formal Darboux neighbourhood.  A global comparison
requires the supplied data of
Theorem~\ref{thm:contractible-stein-descent-criterion-HS};
Steinness alone gives no mixed-HT \(E_2^{\mathrm{HT}}\)-centre action, no all-order QME
counterterm tower, and no comparison-cone contraction
(Definition~\ref{defn:oca-comparison-cone}).
\end{proof}

\begin{rmk}[Compact comparison loci]\label{rmk:physical-interpretations}
The comparison loci of
Corollary~\ref{cor:additional-concrete-realizations} carry only the
formal local conclusion unless the global comparison data exist:
\begin{itemize}
  \item (16) \(K3\setminus C\): the holomorphic-symplectic chart gives
        only the local formal trace-sector first Taylor boundary
        coefficient and its coordinate-dual expression;
  \item (17)--(21): the character, Coulomb, Higgs, multiplicative
        quiver, hypertoric, and wild-character targets require their
        own Stein atlas, OCA comparison, mixed-HT Swiss-cheese action
        datum, all-arity mixed-HT \(E_2^{\mathrm{HT}}\) Morita-centre
        datum, QME counterterm tower,
        locality homotopy, descent null-homotopy, and cone contraction.
\end{itemize}
In each case the formal-Darboux theorem is the local trace-sector
closed-to-open map at a smooth brane point.  The full
derived \(E_2^{\mathrm{HT}}\)-centre comparison requires the global
supplied-data criterion above.
\end{rmk}

\begin{rmk}[Slodowy slices and deformations of Kleinian
singularities]\label{rmk:slodowy-slices}
Slodowy slices \(\mathcal S_e=e+\ker\mathrm{ad}(f)\subset\mathfrak g\)
at subregular nilpotent elements $e$ in semisimple Lie algebras
$\mathfrak g$ intersected with the nilpotent cone produce the
Kleinian surface singularities:
\[
  \mathcal S_e\cap\mathcal N \simeq \C^2/\Gamma_e ,
\]
where \(\Gamma_e\subset SL_2(\C)\) is the finite ADE subgroup
McKay-corresponding to the Dynkin type of \(\mathfrak g\).  The
singular surface is affine and Poisson; its smooth locus is
holomorphic-symplectic, and the minimal resolution is the smooth
holomorphic-symplectic target of item~(9$'$) of
Corollary~\ref{cor:concrete-global-realizations}.  The formal theorem
applies on smooth formal Darboux charts; the singular point requires a
derived or stacky replacement.
\end{rmk}

\begin{rmk}[Compact and non-Stein targets]
\label{rmk:compact-nonstein-targets}
For compact holomorphic-symplectic surfaces or compact Calabi--Yau
\(2\)-folds (\(K3\) being the canonical example), Serre duality
replaces Stein vanishing: by simple-connectedness \(H^1(K3,\mathcal O)=0\), but Serre
duality and the holomorphic symplectic class
\(\sigma\in H^0(K3,\Omega^2)\) give
\[
  H^2(K3,\mathcal O)\;\cong\;H^0(K3,\Omega^2)^{\!\vee}\;\cong\;\C,
\]
the nonvanishing coherent class
(\(h^{2,0}(K3)=h^{0,2}(K3)=1\); separately \(h^{1,1}(K3)=20\) sits in
\(H^1(K3,\Omega^1)\), not \(H^1(\mathcal O)\)).  The descent
obstruction complex on \(\partial X\) has a possible Čech class
sourced by this \(H^2(\mathcal O)\)-class.
The descent class
\(\operatorname{Ob}^{\mathrm{desc}}_\partial\) and the all-order QME
counterterm obstruction class
\(\operatorname{Ob}^{\mathrm{QME}}_{\mathrm{all\mbox{-}order}}\)
therefore live in global cohomology groups not killed by Stein
vanishing; representing them and choosing null-homotopies is the
content of the Hall--Drinfeld--Borcherds filtered comparison complex
(Conjecture~\ref{conj:hall-drinfeld-borcherds-ss}).  A spectral
sequence appears only after its filtration, bidegrees, differentials,
and abutment have been fixed.
Conjecture~\ref{conj:operator-modular-lift} asks for the chain-level
filtered boundary operator algebra of
Definition~\ref{defn:six-boundary-open-closed-objects}~(C6) whose
protected trace has scalar invariant the Igusa cusp form \(\Phi_{10}\) and
its Borcherds square root \(\Delta_5=\Phi_{10}^{1/2}\).  By
Lemma~\ref{lem:scalar-non-faithfulness}, a compact scalar partition
normalisation target such as
\(Z^{K3\times E}_{\mathrm{BPS,un}}=\Phi_{10}^{-1}=\Delta_5^{-2}\)
or \(Z^{K3\times E}_{\mathrm{OP}}=-4096\,\Phi_{10}^{-1}\) acquires
operator meaning only after the scalar normalisation and such a
chain-level \(\mathrm{SC}^{\mathrm{HT}}_{(1,2)}\)-object have been
constructed; the required datum is the operator-valued lift of the modular line
(\S\ref{ssec:open-modular-lift}).  An identification of the local
projective Lie class with curvature of that line additionally requires
a connection and an Atiyah-class pullback.
\end{rmk}



The six typed objects impose the following chain-level constraints.  Each
constraint is either proved at the formal-Darboux stalk or represented
by an obstruction class.

\begin{enumerate}[label=\textup{(V\arabic*)}]
  \item \emph{Object typing.}  Every equality names its
        source object, its target object, its ambient category, and
        its comparison map.
  \item \emph{Scalar non-faithfulness.}  For any proposed
        chain-level
        \(\mathrm{SC}^{\mathrm{HT}}_{(1,2)}\)-object, adjoin the
        acyclic filtered summand \(K\) of the proof of
        Lemma~\ref{lem:scalar-non-faithfulness}; the resulting object
        has the same \(\chi\) but inequivalent chain-level structure.
        This is the obstruction to any scalar-determines-chain
        assertion.
	  \item \emph{Maurer--Cartan identity.}  The Maurer--Cartan identity
	        \(d_{\operatorname{Coder}}\Theta+\tfrac12[\Theta,\Theta]=0\) of
	        \eqref{eq:mc-bar-equation} is the bar \(A_\infty\) identity on
	        \(\bar B^{E_1}(A^{\mathrm{cl}}_{\partial,N})\).  The
	        diagonal-kernel decoration supplies a separate scalar Arnold
	        coefficient identity
	        \(\eta_{12}\eta_{23}+\eta_{23}\eta_{31}+\eta_{31}\eta_{12}=0\)
	        after the chiral coefficient action map has been fixed.
  \item \emph{Bar--cobar equivalence.}  The bar--cobar Quillen equivalence
        \(\Omega^{\mathrm{c}}\bar B^{E_1}(A^{\mathrm{cl}}_{\partial,N})
        \xrightarrow{\sim}A^{\mathrm{cl}}_{\partial,N}\) is verified
        on the Koszul locus by
        Theorem~\ref{thm:tate-bar-cobar-quillen} and
        Theorem~\ref{thm:universal-filtered-cobar-recognition}.  The
        Milnor \(\varprojlim^1\) defect of
        Theorem~\ref{thm:hadamard-mittag-leffler} is the obstruction
        off the criterion's domain.
  \item \emph{Derived centre.}  The bar-resolution computation
        \(C^\bullet_{\mathrm{ch}}(A,A)=\RHom_{A\otimes A^{\mathrm{op}}}(A,A)\)
        on the chiral side is realised on the formal-Darboux stalk by
        \eqref{eq:bar-resolution-derived-centre}; comparison with any
        proposed physical bulk \(\mathcal O^{\mathrm{phys}}_{\mathrm{bulk}}\)
        is mediated by the open--closed comparison
        \(\mathrm{OCA}\) of Definition~\ref{defn:oca-comparison}.  The
        obstruction to constructing this morphism is the non-vanishing of
        the Swiss-cheese action tower or of any component of
        \(\operatorname{Ob}^{\mathrm{global},5}\); the obstruction to a
        quasi-isomorphism is the cohomology of
        \(\operatorname{Cone}(\mathrm{OCA}_N)\).
  \item \emph{Topologisation.}  Exhibit \(T_i=[Q,G_i]\) of
        \eqref{eq:stress-tensor-null-homotopy} chain-level in
        \(\mathcal C_\rho\), then verify the six conditions
        \((\mathrm{T1})\)--\((\mathrm{T6})\) of
        Theorem~\ref{thm:topologisation-by-stress-tensor}.  The
        obstruction is a non-zero anomaly class
        \(\operatorname{Ob}^{\mathrm{anom}}\neq 0\) in the
        operadic convolution complex.
  \item \emph{Object disjointness.}  For each forbidden pair of
        objects \((C_i,C_j)\) in the type table of
        Proposition~\ref{prop:boundary-open-closed-type-table}, exhibit the
        identification that would collapse them and the corresponding
        replacement (\S\ref{ssec:six-typed-objects}).
\end{enumerate}

\subsection{Summary diagram}
\label{ssec:typed-summary-diagram}

The six boundary open--closed objects fit into the following diagram;
the upper square is the formal-Darboux operadic square, and the lower
square is defined after the operator lift and topologisation data are
fixed:
\[
  \begin{array}{ccc}
    \bar B^{E_1}(A^{\mathrm{cl}}_{\partial,N})
       & \xrightarrow{\;\;\mathrm{Verdier/Koszul}\;\;}
       & (A^{\mathrm{cl}}_{\partial,N})^!\\
    \downarrow\!{\scriptstyle \mathrm{Hochschild}}
       & & \downarrow\!{\scriptstyle \mathrm{modules/lines}}\\
    Z^{\mathrm{der},\mathrm{Mor}}_{E_2^{\mathrm{HT}}}(\mathcal C^{\mathrm{op}}_\partial)
       & \xleftarrow{\;\;\mathrm{OCA}\;\;}
       & \mathcal A^{\mathrm{cl}}_{\mathrm{bulk}}\\
    \downarrow\!{\scriptstyle \text{protected trace after operator lift}}
       & & \downarrow\!{\scriptstyle \text{topologisation }T_i=[Q,G_i]}\\
    \chi\bigl(\mathrm{Op}^{\mathrm{HT}}_{\partial X}\bigr)
       & \xrightarrow[\text{under (T1--T6) of}]
                     {\text{Theorem~\ref{thm:topologisation-by-stress-tensor}}}
       & (\mathcal A^{\mathrm{cl}}_{\mathrm{bulk}},
          E_3^{\mathrm{top}})
  \end{array}
\]
The upper square is operadic, proved at the formal-Darboux stalk
(Theorem~\ref{thm:oca-formal-darboux-fibre}) and global only after the
Swiss-cheese action tower, the five rows of
\(\operatorname{Ob}^{\mathrm{global},5}\), and the comparison-cone
contraction of Theorem~\ref{thm:typed-global-mixed-ht-deligne} have been
supplied.  The lower square is the topologisation
passage, under the operator lift and the topologisation datum
of Definition~\ref{defn:topologisation-datum}.  The scalar character
\(\chi\) is non-faithful by
Lemma~\ref{lem:scalar-non-faithfulness}; the filtered boundary
operator datum
\(\mathrm{Op}^{\mathrm{HT}}_{\partial X}\) is a separate filtered
\(\mathrm{SC}^{\mathrm{HT}}\)-object
(\S\ref{ssec:open-modular-lift}) whose chain-level
construction is controlled by the Hall--Borcherds residual class.

The diagram is the typed expression of the relative global mixed-HT
Deligne problem: every assertion factors through the named morphisms
\(\Phi^{\mathrm{HT}}\), \(\mathrm{OCA}\), and the topologisation
arrow; every scalar-to-operator lift factors through
Lemma~\ref{lem:scalar-non-faithfulness}; every operator-level lift to
the filtered boundary operator datum requires its chain-level
construction.

\begin{rmk}[Formal-Darboux fibre and global K3$\times E$ Hall--Borcherds comparison]
\label{rmk:formal-darboux-stalk-scalar-nonfaithfulness}
The formal-Darboux stalk theorem
(Theorem~\ref{thm:main-local}, equivalently
Theorem~\ref{thm:oca-formal-darboux-fibre} at the formal-Darboux
fibre) is a \emph{local} trace-sector \(E_2\)-map in the
admissible filtered nuclear Tate / pro-Matlis category
\(\mathcal B^{\mathrm{adm}}\), equipped with the unscaled
constant-Hamiltonian extension and, on the raw Rees/Weyl trace target,
its pushout by \(\chi_N(K)=N\).  Its correct global role is
twofold:
\begin{enumerate}[label=\textup{(\roman*)}]
  \item to give the \emph{local model} for the open--closed
        comparison \(\mathrm{OCA}\) at every brane vacuum
        \(b\in\widehat{\C^2}\), entering the global Costello--Gwilliam
        bulk--boundary correspondence as the formal-Darboux fibre;
  \item to give the \emph{trace map}
        \(J(f)=\varprojlim_WJ_{\leq W}(f\bmod\mathfrak m^{W+1})\) of
        Theorem~\ref{thm:capelli-projective-anomaly} as the trace
        coordinate of the local stalk; its raw projective commutator
        class is \(\hbar_{\mathrm W}\chi_N(K)[\bar c]\).
\end{enumerate}
The theorem contributes only the formal-Darboux fibre to the following
global \(K3\times E\) constructions, each governed by its own
chain-level theorem in the \(d=1\) algebraic-HT sector or the
\(d\geq 2\) Calabi--Yau-to-chiral comparison:
\begin{itemize}
  \item the compact Hall descent
        \(\widehat{\mathrm{CoHA}}^{\mathrm{red}}_{\Lambda^{2,1}_{\mathrm{II}}}(K3\times E)\)
        on the polarized Borcherds source datum;
  \item the completed Drinfeld-double extension
        \(\mathbf{D}\bigl(\mathcal Y^{+}_{\mathrm{Hall}}(K3\times E)\bigr)\)
        with Hall pairing non-degeneracy on every recognised window;
  \item the current-envelope morphism
        \(\mathrm{Cur}_E(\mathbf D(\mathcal Y^{+}_{\mathrm{Hall}}))
        \to \mathrm{Op}^{\mathrm{HT}}_{K3\times E}\) along the
        reference elliptic curve \(E\);
  \item the \(\mathrm{SC}^{\mathrm{HT}}_{(1,2)}\)-trace compatibility
        whose curvature-section trace is
        \(\Phi_{10}\in H^0(\overline{\mathcal A}_2,\lambda^{10})\);
        after the compact scalar trace datum is constructed, the
        unnormalised \(K3\times E\) BPS partition normalisation target is
        \(Z^{K3\times E}_{\mathrm{BPS,un}}=\Phi_{10}^{-1}
        =\Delta_5^{-2}\), and the OP convention inserts its separate
        normalising factor.
\end{itemize}
The four items \((\beta,\delta,\varepsilon,\tau)\) above are the
four-part chain comparison datum of the boundary-operator
Hall--Borcherds conjecture in the \(d=1\) algebraic
holomorphic-topological comparison;
each carries its own finite-window hypotheses \((W1)\)--\((W8)\),
distinct from the
seven admissibility conditions \((i)\)--\((vii)\) of
Theorem~\ref{thm:main-local}.  Confusing the formal-Darboux local
\(\mathrm{OCA}\) at a single brane vacuum with the global Hall--Borcherds
chain comparison at K3\(\times E\) is the precise collapse forbidden by
the typed separation of source complexes: \(\mathrm{(C4)}\) at one stalk does
not imply \(\mathrm{(C6)}\) on the full target geometry.
\end{rmk}
